\documentclass[UTF8,fontset=none,11pt,a4paper]{ctexart}
\errorcontextlines=40
\setCJKmainfont{Noto Serif CJK SC}[BoldFont={Noto Serif CJK SC Bold},ItalicFont={Noto Serif CJK SC},SlantedFont={Noto Serif CJK SC}]
\setCJKsansfont{Noto Sans CJK SC}
\usepackage[margin=23mm]{geometry}
\usepackage{amsmath,amssymb,mathtools,bm,graphicx,adjustbox,booktabs,longtable,pdflscape,xurl,xcolor}
\usepackage[colorlinks=true,linkcolor=blue,citecolor=teal,urlcolor=blue]{hyperref}
\setlength{\LTcapwidth}{\textwidth}
\newcommand{\fitmath}[1]{\mbox{\begin{adjustbox}{max width=.94\linewidth}$\displaystyle #1$\end{adjustbox}}}
\newcommand{\cellmath}[1]{\mbox{\begin{adjustbox}{max width=\linewidth}$\displaystyle #1$\end{adjustbox}}}
\setlength{\emergencystretch}{3em}
\setcounter{secnumdepth}{2}
\pagestyle{plain}

\title{$\alpha$、$\beta$ 和 $\gamma$ 散度族：\\灵活而稳健的相似性度量}
\author{Andrzej Cichocki \and Shun-ichi Amari}
\date{Entropy 2010, 12, 1532--1568\\中文编译版：2026 年 10 月 15 日}
\begin{document}
\maketitle
{\small\noindent 1. 日本理化学研究所脑科学研究所，高级脑信号处理实验室，日本和光市。\\2. 波兰科学院系统研究所，波兰。\\3. 日本理化学研究所脑科学研究所，数学神经科学实验室，日本和光市。\\Andrzej Cichocki 属于机构 1、2，为通讯作者；Shun-ichi Amari 属于机构 3。\\通讯邮箱：\href{mailto:a.cichocki@riken.jp}{a.cichocki@riken.jp}；电话：+81-48467-9668；传真：+81-48467-9686。\\收稿：2010 年 4 月 26 日；录用：2010 年 6 月 1 日；发表：2010 年 6 月 14 日。\par}
\begin{quote}\small 本译文依据出版社 XML 全文恢复 LaTeX 后翻译；未取得作者原始 .tex 文件。保留原文公式、表格、章节和参考文献编号。原文 DOI：\href{https://doi.org/10.3390/e12061532}{10.3390/e12061532}。原作 © 2010 作者，采用 \href{https://creativecommons.org/licenses/by/3.0/}{CC BY 3.0} 许可；本稿的改动为中文翻译及排版调整。\end{quote}
\begin{abstract}
本文扩展并综述 \(\alpha\)、\(\beta\) 和 \(\gamma\) 散度的若干广泛族，讨论它们的基本性质。现有文献通常只考虑某一种非对称散度，即 \(\alpha\)、\(\beta\) 或 \(\gamma\) 散度。本文说明，这些散度各自构成具有一致性质的族。通过适当的非线性变换，我们建立不同散度之间的联系与对应关系。例如，可以从 \(\alpha\) 散度直接生成 \(\beta\) 散度，反之亦然。我们还说明，一类新的 \(\gamma\) 散度既可以从 \(\beta\) 散度族生成，也可以从 \(\alpha\) 散度族生成。本文连接这些散度，并说明它们与 Tsallis 熵及 Rényi 熵的关系。其中大多数散度都有自然的信息论解释。
\end{abstract}
\noindent 关键词：相似性度量；广义散度；扩展的 Itakura–Saito 型距离；Csiszár–Morimoto 散度与 Bregman 散度；Tsallis 熵与 Rényi 熵
\tableofcontents
\clearpage
\setcounter{section}{0}
\section{\texorpdfstring{引言}{引言}}
\label{sec1-entropy-12-01532}
许多用于分类和聚类的机器学习算法采用不同的相异性度量。信息论、凸分析和信息几何在这些散度的构造中起着关键作用[\hyperlink{B1-entropy-12-01532}{1}--\hyperlink{B25-entropy-12-01532}{25}]。


最常用的度量包括平方欧氏距离和 Kullback–Leibler 散度。近年来，Csiszár–Morimoto 的 f 散度与 Bregman 散度等广义散度，为机器学习算法提供了其他选择[\hyperlink{B26-entropy-12-01532}{26}--\hyperlink{B34-entropy-12-01532}{34}]。本文讨论 Csiszár–Morimoto 散度与 Bregman 散度中一个稳健的参数化子类：\(\alpha\) 散度和 \(\beta\) 散度。面对离群值和加性噪声，它们可以给出更稳健、精度更高的解。我们还建立它们与一类新的稳健 \(\gamma\) 散度[\hyperlink{B35-entropy-12-01532}{35}]之间的联系，并将后者扩展为所谓的 \(\alpha\)–\(\gamma\) 散度。


本文将散度视为相似性或相异性度量。通常，它衡量两个概率分布 \(\mathbf{P}\) 与 \(\mathbf{Q}\) 之间的准距离或有向差异，也可以定义在无约束的非负多维数组和模式上。


本文假设 \(\mathbf{P}\) 和 \(\mathbf{Q}\) 是正的有限测度或密度，未必归一化。对于归一化密度，我们明确称其为概率密度。除非另有说明，均假设这些测度是连续的。信息散度度量两条概率曲线之间的距离。本文只讨论由非负信号或时间序列表示的一维概率曲线；推广到二维或更高维变量是直接的。通常，一个密度 \(\mathbf{Q}(x)\) 已知且固定，另一个密度 \(\mathbf{P}(\mathbf{x})\) 通过学习或调整，在某种意义下尽可能接近 \(\mathbf{Q}(\mathbf{x})\)。例如，离散密度 \(\mathbf{Q}\) 对应观测数据，而向量 \(\mathbf{P}\) 对应待估计的数据或期望数据，并满足所假设模型施加的约束。对于非负矩阵分解（Non-negative Matrix Factorization，NMF），\(\mathbf{Q}\) 对应数据矩阵 \(\mathbf{Y}\)，\(\mathbf{P}\) 对应估计矩阵 \(\hat{\mathbf{Y}} = \mathbf{A}\mathbf{X}\)，二者也可以交换[\hyperlink{B30-entropy-12-01532}{30}]。


两个密度之间的距离若满足以下条件，就称为度量：

\begin{enumerate}
\item \(D(\mathbf{P}\,||\,\mathbf{Q}) \geq 0\)，等号成立当且仅当 \(\mathbf{P} = \mathbf{Q}\)，即非负性与正定性；


\item \(D(\mathbf{P}\,||\,\mathbf{Q}) = D(\mathbf{Q}\,||\,\mathbf{P})\)，即对称性；


\item \(D(\mathbf{P}\,||\,\mathbf{Z}) \leq D(\mathbf{P}\,||\,\mathbf{Q}) + D(\mathbf{Q}\,||\,\mathbf{Z})\)，即次可加性或三角不等式。

\end{enumerate}





只满足条件 1 的距离不构成度量，称为非对称散度。


在图像分析、模式识别和统计机器学习等应用中，常采用信息论散度来替代平方欧氏距离或 \(l_{p}\) 范数距离[\hyperlink{B28-entropy-12-01532}{28}]。Kullback–Leibler、Hellinger 和 Jensen–Shannon 等信息散度，是估计分布相似性的重要工具，在信息几何中已有长期研究。


散度的概念不限于欧氏空间，还可以借助 Radon–Nikodym 导数扩展到抽象空间，例如文献[\hyperlink{B36-entropy-12-01532}{36}]。设 \((\mathbf{X},\mathcal{A},\mu)\) 是一个测度空间，其中 \(\mu\) 是 \((\mathbf{X},\mathcal{A})\) 上的有限测度或 \(\sigma\) 有限测度。设 \(\mathbf{P}\) 和 \(\mathbf{Q}\) 是 \((\mathbf{X},\mathcal{A})\) 上的两个概率测度，且 \(\mathbf{P} < < \mu\)、\(\mathbf{Q} < < \mu\) 关于测度 \(\mu\) 绝对连续，例如取 \(\mu = \mathbf{P} + \mathbf{Q}\)；设 \(p = \frac{d\mathbf{P}}{d\mu}\) 和 \(q = \frac{d\mathbf{Q}}{d\mu}\) 分别为 \(\mathbf{P}\) 和 \(\mathbf{Q}\) 关于 \(\mu\) 的 Radon–Nikodym 导数，即密度。使用这些记号，两个概率分布之间的 Kullback–Leibler（KL）散度可写为

\setcounter{equation}{0}
\begin{equation}
\label{FD1-entropy-12-01532}
\fitmath{D_{KL}{(\mathbf{P}\,||\,\mathbf{Q})} = \int_{\mathbf{X}}p{(\mathbf{x})}\log\left( \frac{p(\mathbf{x})}{q(\mathbf{x})} \right)d\mu{(\mathbf{x})}}
\end{equation}

它与 Shannon 熵

\setcounter{equation}{1}
\begin{equation}
\label{FD2-entropy-12-01532}
\fitmath{H_{S}{(\mathbf{P})} = - \int_{\mathbf{X}}p{(\mathbf{x})}\,\log{(p{(\mathbf{x})})}d\mu{(\mathbf{x})}}
\end{equation}

之间满足

\[
\fitmath{D_{KL}{(\mathbf{P}\,||\,\mathbf{Q})} = V_{S}{(\mathbf{P},\mathbf{Q})} - H_{S}{(\mathbf{P})}}
\]

其中

\[
\fitmath{V_{S}{(\mathbf{P},\mathbf{Q})} = - \int_{\mathbf{X}}p{(\mathbf{x})}\log{(q{(\mathbf{x})})}d\mu{(\mathbf{x})}}
\]

为 Shannon 交叉熵，前提是相应积分存在。从测度论的角度，若由 \(\mathbf{P}\) 诱导的测度关于由 \(\mathbf{Q}\) 诱导的测度绝对连续，则该积分存在。本文始终假设所涉及的积分存在。


利用一族广义对数函数，也称 \(\alpha\) 对数，可以推广 Kullback–Leibler 散度：

\setcounter{equation}{2}
\begin{equation}
\label{FD3-entropy-12-01532}
\fitmath{\log_{\;\alpha}{(x)} = \frac{1}{1 - \alpha}{(x^{1 - \alpha} - 1)}{(\text{当}\;\; x > 0)}}
\end{equation}

这是关于 x 的幂函数，幂次为 \(1 - \alpha\)；在极限 \(\left. \alpha\rightarrow 1 \right.\) 下，它退化为自然对数。利用幂函数（\hyperref[FD3-entropy-12-01532]{3}），常可构造对离群值更稳健的散度，获得更好或更灵活的性能，例如文献[\hyperlink{B37-entropy-12-01532}{37}]。


利用这种扩展，我们推导并回顾 \(\alpha\)、\(\beta\) 和 \(\gamma\) 三个散度族，它们都是 KL 散度的推广。我们还说明它们与 Tsallis 熵和 Rényi 熵的关系，见\hyperref[app1-entropy-12-01532]{附录 A}和\hyperref[app2-entropy-12-01532]{附录 B}。此外，还将说明如何利用幂函数从 Csiszár–Morimoto 的 f 散度导出 \(\alpha\) 散度，以及从 Bregman 散度导出 \(\beta\) 散度。


与 Zhang[\hyperlink{B22-entropy-12-01532}{22}--\hyperlink{B24-entropy-12-01532}{24}]以及 Hein 和 Bousquet[\hyperlink{B21-entropy-12-01532}{21}]的工作类似，我们的一个动机是说明广泛的散度类之间的紧密联系，并在同一框架内处理已知散度、中间散度和新散度。不过，本文的相似性度量与文献[\hyperlink{B21-entropy-12-01532}{21}]提出的度量不同，所用方法和结果也与文献[\hyperlink{B22-entropy-12-01532}{22}]有明显区别。


已有一些工作尝试统一散度函数，尤其是与 \(\alpha\) 散度有关的函数，包括 Zhu 和 Rohwer[\hyperlink{B10-entropy-12-01532}{10},\allowbreak \hyperlink{B11-entropy-12-01532}{11}]、Amari 和 Nagaoka[\hyperlink{B3-entropy-12-01532}{3}]、Taneya[\hyperlink{B38-entropy-12-01532}{38},\allowbreak \hyperlink{B39-entropy-12-01532}{39}]、Zhang[\hyperlink{B22-entropy-12-01532}{22}]以及 Gorban 和 Judge[\hyperlink{B40-entropy-12-01532}{40}]的研究。Zhang 在文献[\hyperlink{B22-entropy-12-01532}{22}]及其后续工作[\hyperlink{B23-entropy-12-01532}{23},\allowbreak \hyperlink{B24-entropy-12-01532}{24}]中，通过基于凸函数的统一方法，研究了信息几何与各种散度函数之间的深层联系，其中主要涉及 \(\alpha\) 散度。然而，已有工作没有明确讨论 \(\alpha\)、\(\beta\) 和 \(\gamma\) 这三个基本散度类之间的全部联系。本文还回顾并扩展它们的若干基本性质。


本文结果覆盖的广义散度函数类包含许多有用的损失函数，例如基于相对熵、广义 Kullback–Leibler 散度或 I 散度、Hellinger 距离、Jensen–Shannon 散度、J 散度、Pearson 与 Neyman 卡方散度、三角判别散度及算术–几何散度的损失函数。我们还说明如何生成新的散度。特别地，我们构造一族新的 \(\alpha\)–\(\gamma\) 散度，以及具有尺度不变性的 Itakura–Saito 型距离；后者属于更广泛的 \(\beta\) 散度类。这些新的尺度不变散度扩展了 \(\beta\) 散度族和 \(\gamma\) 散度族。本文讨论的散度函数具有灵活性：通过设定调节参数，可以生成许多已知且常用的特例。此外，通过自适应调节参数，可以优化学习算法的代价函数，并在噪声与离群值存在时估计模型参数。本文讨论的散度，尤其是 \(\beta\) 散度和 \(\gamma\) 散度，在某些参数取值下对离群值具有稳健性。


这些散度族可以为散度的选择乃至新度量的构造提供指导，并利用 Csiszár–Morimoto 散度、Bregman 散度及其基本性质，将不同散度置于同一框架。它们通常定义在未归一化的有限测度上，而不局限于归一化概率。这使我们可以分析不同大小、具有不同权重的模式，例如尺寸不同的图像或长度不同的文档。这类测度在神经计算、模式识别、学习、估计、推断和优化中也很重要。我们已将其中一部分散度作为代价函数，并在需要时加入约束或正则项，导出了用于非负矩阵分解和非负张量分解的乘法更新算法及加性投影梯度算法[\hyperlink{B30-entropy-12-01532}{30},\allowbreak \hyperlink{B31-entropy-12-01532}{31}]。


散度与概率分布流形的不变几何性质密切相关[\hyperlink{B5-entropy-12-01532}{5}--\hyperlink{B7-entropy-12-01532}{7}]。

\setcounter{section}{1}
\section{\texorpdfstring{\(\alpha\) 散度族}{α 散度族}}
\label{sec2-entropy-12-01532}
\(\alpha\) 散度可以从 Csiszár–Morimoto 的 f 散度导出；Amari 最近还说明，通过一些处理，也可以从 Bregman 散度导出它[\hyperlink{B3-entropy-12-01532}{3},\allowbreak \hyperlink{B6-entropy-12-01532}{6}]，另见\hyperref[app1-entropy-12-01532]{附录 A}。\(\alpha\) 散度由 Chernoff 提出[\hyperlink{B41-entropy-12-01532}{41}]，Amari[\hyperlink{B1-entropy-12-01532}{1},\allowbreak \hyperlink{B3-entropy-12-01532}{3},\allowbreak \hyperlink{B5-entropy-12-01532}{5},\allowbreak \hyperlink{B6-entropy-12-01532}{6}]和其他研究者对它进行了广泛研究与扩展。有关修改或推广，可参阅 Liese 和 Vajda[\hyperlink{B36-entropy-12-01532}{36}]、Minka[\hyperlink{B42-entropy-12-01532}{42}]、Taneja[\hyperlink{B38-entropy-12-01532}{38},\allowbreak \hyperlink{B39-entropy-12-01532}{39},\allowbreak \hyperlink{B43-entropy-12-01532}{43}]、Cressie–Read[\hyperlink{B44-entropy-12-01532}{44}]、Zhu–Rohwer[\hyperlink{B10-entropy-12-01532}{10},\allowbreak \hyperlink{B11-entropy-12-01532}{11}]以及 Zhang[\hyperlink{B22-entropy-12-01532}{22}--\hyperlink{B25-entropy-12-01532}{25}]的工作。我们研究 \(\alpha\) 散度族的一个动机，是为非负矩阵分解和非负张量分解开发灵活、高效的学习算法，以统一并扩展已有算法，包括期望最大化最大似然算法（Expectation Maximization Maximum Likelihood，EMML）和图像空间重建算法（Image Space Reconstruction Algorithm，ISRA）；见文献[\hyperlink{B30-entropy-12-01532}{30},\allowbreak \hyperlink{B45-entropy-12-01532}{45}]及其中的参考文献。


\setcounter{section}{2}
\setcounter{subsection}{0}
\subsection{\texorpdfstring{非对称 \(\alpha\) 散度}{非对称 α 散度}}
\label{sec2dot1-entropy-12-01532}
基本的非对称 \(\alpha\) 散度可定义为[\hyperlink{B3-entropy-12-01532}{3}]：

\setcounter{equation}{3}
\begin{equation}
\label{FD4-entropy-12-01532}
\fitmath{D_{A}^{(\alpha)}{(\mathbf{P}\,||\,\mathbf{Q})} = \frac{1}{\alpha(\alpha - 1)}\int\left( {p^{\,\alpha}{(\mathbf{x})}\, q^{\, 1 - \alpha}{(\mathbf{x})} - \alpha\, p{(\mathbf{x})} + {(\alpha - 1)}\, q{(\mathbf{x})}} \right)d\mu{(\mathbf{x})},\alpha \in \mathbb{R} \smallsetminus {\{ 0,1\}}}
\end{equation}

其中 \(p(\mathbf{x})\) 和 \(q(\mathbf{x})\) 不必归一化。


\(\alpha\) 散度可通过广义 KL 散度表示为

\setcounter{equation}{4}
\begin{equation}
\label{FD5-entropy-12-01532}
\fitmath{D_{A}^{(\alpha)}{(\mathbf{P}\,||\;\mathbf{Q})} = D_{GKL}^{(\alpha)}{(\mathbf{P}\,||\;\mathbf{Q})} = - \frac{1}{\alpha}\int\left( {p\log_{\,\alpha}\left( \frac{q}{p} \right) + p - q} \right)d\mu{(\mathbf{x})},\;\;\alpha \in \mathbb{R} \smallsetminus {\{ 0,1\}}}
\end{equation}




对于概率质量函数为 \(\mathbf{P} = \lbrack p_{1},p_{2},\ldots,p_{n}\rbrack\) 和 \(\mathbf{Q} = \lbrack q_{1},q_{2},\ldots,q_{n}\rbrack\) 的离散概率测度，离散 \(\alpha\) 散度可以写成可分离形式：

\setcounter{equation}{5}
\begin{equation}
\label{FD6-entropy-12-01532}
\fitmath{D_{A}^{(\alpha)}{(\mathbf{P}\,||\,\mathbf{Q})} = \sum\limits_{i = 1}^{n}d_{A}^{(\alpha)}(p_{i}{\,||\,}q_{i}{) =}\frac{1}{\alpha(\alpha - 1)}\sum\limits_{i = 1}^{n}\left( {p_{i}^{\alpha}\, q_{i}^{1 - \alpha} - \alpha\, p_{i} + {(\alpha - 1)}\, q_{i}} \right),\alpha \in \mathbb{R} \smallsetminus {\{ 0,1\}}}
\end{equation}

其中 \(d_{A}^{(\alpha)}(p_{i}{\,||\,}q_{i}{) =}\left( {p_{i}^{\alpha}q_{i}^{1 - \alpha} - \alpha p_{i} + {(\alpha - 1)}q_{i}} \right)/\left( {\alpha(\alpha - 1)} \right).\)。


这一 \(\alpha\) 散度形式与文献[\hyperlink{B1-entropy-12-01532}{1}]中的损失函数略有不同，后者只定义在概率分布上，即假设 \(\int p(\mathbf{x})d\mu(\mathbf{x}) = 1\) 且 \(\int q(\mathbf{x})d\mu(\mathbf{x}) = 1\)。Zhu 和 Rohwer 在文献[\hyperlink{B10-entropy-12-01532}{10},\allowbreak \hyperlink{B11-entropy-12-01532}{11}]中通过加入额外项，将其扩展到正测度；另见 Amari 和 Nagaoka[\hyperlink{B3-entropy-12-01532}{3}]。这些额外项使该定义能够处理未归一化密度，即正测度，其作用与广义 KL 散度或 I 散度中的附加项相同。


将 Kullback–Leibler 散度扩展为 \(\alpha\) 散度族，是统一理解广泛散度函数的一步，因为 \(\alpha\) 散度的非负性既可以由 Jensen 不等式解释，也可以由算术–几何平均不等式解释。正如 Jun Zhang 指出的，这一点可以充分利用凸函数不等式，或者更一般地利用他提出的基于凸性的研究方法[\hyperlink{B22-entropy-12-01532}{22}]。


对于归一化密度 \(\overline{p}{(\mathbf{x})} = p{(\mathbf{x})}/\int p{(\mathbf{x})}d\mu{(\mathbf{x})}\) 和 \(\overline{q}{(\mathbf{x})} = q{(\mathbf{x})}/\int q{(\mathbf{x})}d\mu{(\mathbf{x})}\)，\(\alpha\) 散度简化为[\hyperlink{B1-entropy-12-01532}{1},\allowbreak \hyperlink{B44-entropy-12-01532}{44}]：

\setcounter{equation}{6}
\begin{equation}
\label{FD7-entropy-12-01532}
\fitmath{D_{A}^{(\alpha)}(\overline{\mathbf{P}}{\,||\,}\overline{\mathbf{Q}}{) =}\frac{1}{\alpha(\alpha - 1)}\left( {\int{\overline{p}}^{\,\alpha}{(\mathbf{x})}\,{\overline{q}}^{\, 1 - \alpha}{(\mathbf{x})}d\mu{(\mathbf{x})} - 1} \right),\alpha \in \mathbb{R} \smallsetminus {\{ 0,1\}}}
\end{equation}

它还与 Tsallis 散度及 Tsallis 熵[\hyperlink{B46-entropy-12-01532}{46}]有关，见\hyperref[app1-entropy-12-01532]{附录 A}：

\setcounter{equation}{7}
\begin{equation}
\label{FD8-entropy-12-01532}
\fitmath{H_{T}^{(\alpha)}{(\overline{\mathbf{P}})} = \frac{1}{1 - \alpha}\left( {\int{\overline{p}}^{\,\alpha}{(\mathbf{x})}\, d\mu{(\mathbf{x})} - 1} \right) = - \int{\overline{p}}^{\,\alpha}{(\mathbf{x})}\log_{\;\alpha}\overline{p}{(\mathbf{x})}d\mu{(\mathbf{x})}}
\end{equation}

其中假设右侧积分存在。


Tsallis 熵实际上最早由 Havrda 和 Charvat 于 1967 年定义[\hyperlink{B47-entropy-12-01532}{47}]，后来几乎被遗忘；Tsallis 于 1988 年在另一背景下重新提出了它[\hyperlink{B46-entropy-12-01532}{46}]，见\hyperref[app1-entropy-12-01532]{附录 A}。


不同作者对参数 \(\alpha\) 的约定并不相同。例如，采用 Amari 的记号 \(\alpha_{A}\)，并令 \(\alpha = (1 - \alpha_{A})/2\)，则 \(\alpha\) 散度具有如下形式[\hyperlink{B1-entropy-12-01532}{1}--\hyperlink{B3-entropy-12-01532}{3},\allowbreak \hyperlink{B10-entropy-12-01532}{10},\allowbreak \hyperlink{B11-entropy-12-01532}{11},\allowbreak \hyperlink{B22-entropy-12-01532}{22}--\hyperlink{B25-entropy-12-01532}{25},\allowbreak \hyperlink{B44-entropy-12-01532}{44},\allowbreak \hyperlink{B48-entropy-12-01532}{48},\allowbreak \hyperlink{B49-entropy-12-01532}{49}]：

\setcounter{equation}{8}
\begin{equation}
\label{FD9-entropy-12-01532}
\fitmath{{\widetilde{D}}_{A}^{(\alpha_{A})}{(\underline{P}\,||\, \underline{Q})} = \frac{4}{1 - \alpha_{A}^{2}}\int\left( {\frac{1 - \alpha_{A}}{2}\, p + \frac{1 + \alpha_{A}}{2}\, q - p^{\frac{1 - \alpha_{A}}{2}}q^{\frac{1 + \alpha_{A}}{2}}} \right)d\mu{(\mathbf{x})},\;\;\alpha_{A} \in \mathbb{R} \smallsetminus {\{ \pm 1\}}}
\end{equation}




当 \(\alpha\) 从 0 变化到 1 时，\(\alpha_{A}\) 从 \(- 1\) 变化到 1。\(\alpha_{A}\) 与 \(- \alpha_{A}\) 之间存在对偶关系，即 \({\widetilde{D}}_{A}^{(\alpha_{A})}{(\mathbf{P}\,||\,\mathbf{Q})} = {\widetilde{D}}_{A}^{( - \alpha_{A})}{(\mathbf{Q}\,||\,\mathbf{P})}\)。


当 \(\alpha = 2,0.5, - 1\) 时，由式（\hyperref[FD4-entropy-12-01532]{4}）分别得到熟知的 Pearson 卡方散度、Hellinger 距离和逆 Pearson 卡方散度；最后一种也称 Neyman 卡方距离：

\setcounter{equation}{9}
\begin{equation}
\label{FD10-entropy-12-01532}
\fitmath{\begin{array}{rlr}
{D_{A}^{(2)}{(\mathbf{P}\,||\,\mathbf{Q})}} & = & {D_{P}{(\mathbf{P}\,||\,\mathbf{Q})} = \frac{1}{2}\int\frac{{(p{(\mathbf{x})} - q{(\mathbf{x})})}^{2}}{q(\mathbf{x})}d\mu{(\mathbf{x})}} \\
{D_{A}^{(1/2)}{(\mathbf{P}\,||\,\mathbf{Q})}} & = & {2D_{H}{(\mathbf{P}\,||\,\mathbf{Q})} = 2\int{(\sqrt{p(\mathbf{x})} - \sqrt{q(\mathbf{x})})}^{2}d\mu{(\mathbf{x})}} \\
{D_{A}^{( - 1)}{(\mathbf{P}\,||\,\mathbf{Q})}} & = & {D_{N}{(\mathbf{P}\,||\,\mathbf{Q})} = \frac{1}{2}\int\frac{{(p{(\mathbf{x})} - q{(\mathbf{x})})}^{2}}{p(\mathbf{x})}d\mu{(\mathbf{x})}}
\end{array}}
\end{equation}




对于奇异参数值 \(\alpha = 1\) 和 \(\alpha = 0\)，\(\alpha\) 散度（\hyperref[FD4-entropy-12-01532]{4}）必须分别通过 \(\left. \alpha\rightarrow 1 \right.\) 和 \(\left. \alpha\rightarrow 0 \right.\) 的极限来定义。当 \(\left. \alpha\rightarrow 1 \right.\) 时，利用洛必达法则计算极限，得到 Kullback–Leibler 散度：

\setcounter{equation}{10}
\begin{equation}
\label{FD11-entropy-12-01532}
\fitmath{\begin{array}{cl}
{D_{KL}{(\mathbf{P}\,||\,\mathbf{Q})}} & {= \lim\limits_{\alpha\rightarrow 1}D_{A}^{(\alpha)}{(\mathbf{P}\,||\,\mathbf{Q})} = \lim\limits_{\alpha\rightarrow 1}\int\frac{p^{\,\alpha}q^{\, 1 - \alpha}\log p - p^{\,\alpha}q^{1 - \alpha}\log q - p + q}{2\alpha - 1}d\mu{(\mathbf{x})}} \\
 & {= \int\left( {p\;\log\left( \frac{p}{q} \right) - p + q} \right)d\mu{(\mathbf{x})}}
\end{array}}
\end{equation}

其中约定 \(0/0 = 0,\; 0\log(0) = 0\)，且对于 \(p > 0\)，有 \(p/0 = \infty\)。


由不等式 \(p\log p \geq p - 1\) 可知，KL 的 I 散度非负，且仅在 \(\mathbf{P} = \mathbf{Q}\) 时为零。


类似地，当 \(\left. \alpha\rightarrow 0 \right.\) 时，得到反向 Kullback–Leibler 散度：

\setcounter{equation}{11}
\begin{equation}
\label{FD12-entropy-12-01532}
\fitmath{D_{KL}{(\mathbf{Q}||\mathbf{P})} = \lim\limits_{\alpha\rightarrow 0}D_{A}^{(\alpha)}{(\mathbf{P}\,||\,\mathbf{Q})} = \int\left( {q\log\left( \frac{q}{p} \right) - q + p} \right)d\mu{(\mathbf{x})}}
\end{equation}




因此，\(\alpha\) 散度可更明确地写为

\setcounter{equation}{12}
\begin{equation}
\label{FD13-entropy-12-01532}
\fitmath{D_{A}^{(\alpha)}{(\mathbf{P}\,||\,\mathbf{Q})} = \left\{ \begin{matrix}
{\frac{1}{\alpha(\alpha - 1)}\int\left( {q{(\mathbf{x})}\left\lbrack {\left( \frac{p(\mathbf{x})}{q(\mathbf{x})} \right)^{\alpha} - 1} \right\rbrack - \alpha\,{\lbrack p{(\mathbf{x})} - q{(\mathbf{x})}\rbrack}} \right)d\mu{(\mathbf{x})},\hspace{3.33333pt}\hspace{3.33333pt}\alpha \neq 0,1} \\
 \\
{\int\left( {q{(\mathbf{x})}\log\frac{q(\mathbf{x})}{p(\mathbf{x})} + p{(\mathbf{x})} - q{(\mathbf{x})}} \right)d\mu{(\mathbf{x})},\alpha = 0} \\
 \\
{\int\left( {p{(\mathbf{x})}\log\frac{p(\mathbf{x})}{q(\mathbf{x})} - p{(\mathbf{x})} + q{(\mathbf{x})}} \right)d\mu{(\mathbf{x})},\alpha = 1}
\end{matrix} \right.}
\end{equation}




\(\alpha\) 散度平滑地连接 I 散度 \(D_{KL}{(\mathbf{P}\,||\,\mathbf{Q})}\) 与反向 I 散度 \(D_{KL}{(\mathbf{Q}\,||\,\mathbf{P})}\)，并经过 Hellinger 距离[\hyperlink{B50-entropy-12-01532}{50}]。它还平滑地连接 Pearson 卡方散度与 Neyman 卡方散度，并经过这两个 I 散度[\hyperlink{B10-entropy-12-01532}{10},\allowbreak \hyperlink{B11-entropy-12-01532}{11}]。


\(\alpha\) 散度是 Csiszár–Morimoto 的 f 散度[\hyperlink{B17-entropy-12-01532}{17},\allowbreak \hyperlink{B19-entropy-12-01532}{19}]的特例；Ali 和 Silvey 后来也提出了这类散度[\hyperlink{B20-entropy-12-01532}{20}]，Morimoto 则独立地给出了定义[\hyperlink{B51-entropy-12-01532}{51}]。Csiszár–Morimoto 的 f 散度与函数 \(f(u)\) 对应，其中该函数在 \((0,\infty)\) 上为凸函数，并满足 \(f(1) = 0\)：

\setcounter{equation}{13}
\begin{equation}
\label{FD14-entropy-12-01532}
\fitmath{D_{f}{(\mathbf{P}\,||\,\mathbf{Q})} = \int q{(\mathbf{x})}f\left( \frac{p(\mathbf{x})}{q(\mathbf{x})} \right)d\mu{(\mathbf{x})}}
\end{equation}




定义 \(0f(0/0) = 0\) 和
\[
\fitmath{0f{(a/0)} = \lim_{t\rightarrow 0}tf{(a/t)} = \lim_{u\rightarrow\infty}f{(u)}/u}
\]
。取 \(f{(u)} = {(u^{\alpha} - \alpha u + \alpha - 1)}{)/}{(\alpha^{2} - \alpha)}\)，即可得到式（\hyperref[FD4-entropy-12-01532]{4}）。


Csiszár–Morimoto 的 f 散度具有以下性质[\hyperlink{B6-entropy-12-01532}{6},\allowbreak \hyperlink{B17-entropy-12-01532}{17},\allowbreak \hyperlink{B18-entropy-12-01532}{18},\allowbreak \hyperlink{B20-entropy-12-01532}{20},\allowbreak \hyperlink{B22-entropy-12-01532}{22},\allowbreak \hyperlink{B36-entropy-12-01532}{36},\allowbreak \hyperlink{B52-entropy-12-01532}{52},\allowbreak \hyperlink{B53-entropy-12-01532}{53}]：

\begin{itemize}
\item 非负性：Csiszár–Morimoto 的 f 散度始终非负，且仅在概率密度 \(p(\mathbf{x})\) 与 \(q(\mathbf{x})\) 相同时为零。对于归一化密度，这可直接由 Jensen 不等式得到：

\setcounter{equation}{14}
\begin{equation}
\label{FD15-entropy-12-01532}
\fitmath{D_{f}{(\mathbf{P}||\mathbf{Q})} = \int qf\left( \frac{p}{q} \right)d\mu{(\mathbf{x})} \geq f\left( {\int q{(\frac{p}{q})}d\mu{(\mathbf{x})}} \right) = f{(1)} = 0.}
\end{equation}




\item 广义熵：该散度对应如下形式的广义 f 熵：

\setcounter{equation}{15}
\begin{equation}
\label{FD16-entropy-12-01532}
\fitmath{H_{f}{(\mathbf{P})} = - \int f{(p{(\mathbf{x})})}d\mu{(\mathbf{x})}}
\end{equation}

当 \(f(p) = p\log(p)\) 时，得到 Shannon 熵这一特例。\(H_{f}\) 为凹函数，而 f 为凸函数。


\item 凸性：对于任意 \(0 \leq \lambda \leq 1\)，有

\setcounter{equation}{16}
\begin{equation}
\label{FD17-entropy-12-01532}
\fitmath{D_{f}{(\lambda}\mathbf{P}_{1} + {(1 - \lambda)}\mathbf{P}_{2}{\,||\,\lambda}\mathbf{Q}_{1} + {(1 - \lambda)}\mathbf{Q}_{2}{) \leq \lambda}D_{f}(\mathbf{P}_{1}{||}\mathbf{Q}_{1}{) + {(1 - \lambda)}}D_{f}(\mathbf{P}_{2}{||}\mathbf{Q}_{2}{).}}
\end{equation}




\item 缩放性质：对于任意正常数 \(c > 0\)，有

\setcounter{equation}{17}
\begin{equation}
\label{FD18-entropy-12-01532}
\fitmath{cD_{f}{(\mathbf{P}||\mathbf{Q})} = D_{cf}{(\mathbf{P}||\mathbf{Q})}}
\end{equation}




\item 不变性：f 散度在双射变换下保持不变[\hyperlink{B6-entropy-12-01532}{6},\allowbreak \hyperlink{B20-entropy-12-01532}{20}]。也就是说，当 \(\mathbf{x}\) 通过

\setcounter{equation}{18}
\begin{equation}
\label{FD19-entropy-12-01532}
\fitmath{\mathbf{y} = \mathbf{k}(\mathbf{x})}
\end{equation}

双射变换为 \(\mathbf{y}\) 时，概率分布 \(\mathbf{P}(\mathbf{x})\) 变为 \(\widetilde{\mathbf{P}}{(\mathbf{y})}\)，但仍有

\setcounter{equation}{19}
\begin{equation}
\label{FD20-entropy-12-01532}
\fitmath{D_{f}{(\mathbf{P}||\mathbf{Q})} = D_{f}(\widetilde{\mathbf{P}}{||}\widetilde{\mathbf{Q}})}
\end{equation}

此外，对于任意常数 c 以及 \(\widetilde{f}{(u)} = f{(u)} - c{(u - 1)}\)，有

\setcounter{equation}{20}
\begin{equation}
\label{FD21-entropy-12-01532}
\fitmath{D_{f}{(\mathbf{P}||\mathbf{Q})} = D_{\widetilde{f}}{(\mathbf{P}||\mathbf{Q})}}
\end{equation}

并且

\setcounter{equation}{21}
\begin{equation}
\label{FD22-entropy-12-01532}
\fitmath{D_{f}{(\mathbf{P}||\mathbf{Q})} = D_{f^{*}}{(\mathbf{Q}||\mathbf{P})}}
\end{equation}

其中 \(f^{*}{(u)} = uf{(1/u)}\) 称为共轭函数。


\item 对称性：对于任意 Csiszár–Morimoto 的 f 散度，当 \(f_{sym}{(u)} = f{(u)} + f^{*}{(u)}\) 时，可以构造对称散度。


\item 有界性：正测度或正密度上的 Csiszár–Morimoto 的 f 散度具有下列界，前提是相关极限存在且有限[\hyperlink{B28-entropy-12-01532}{28},\allowbreak \hyperlink{B36-entropy-12-01532}{36}]：

\setcounter{equation}{22}
\begin{equation}
\label{FD23-entropy-12-01532}
\fitmath{0 \leq D_{f}{(\mathbf{P}||\mathbf{Q})} \leq \lim\limits_{u\rightarrow 0^{+}}\left\{ {f{(u)} + uf{(\frac{1}{u})}} \right\}}
\end{equation}

此外，还有[\hyperlink{B54-entropy-12-01532}{54}]：

\setcounter{equation}{23}
\begin{equation}
\label{FD24-entropy-12-01532}
\fitmath{0 \leq D_{f}{(\mathbf{P}||\mathbf{Q})} \leq \int{(p - q)}f'{(\frac{p}{q})}d\mu{(\mathbf{x})}}
\end{equation}



\end{itemize}





利用 Csiszár–Morimoto 的 f 散度的基本性质，可以建立 \(\alpha\) 散度的基本性质[\hyperlink{B3-entropy-12-01532}{3},\allowbreak \hyperlink{B6-entropy-12-01532}{6},\allowbreak \hyperlink{B22-entropy-12-01532}{22},\allowbreak \hyperlink{B40-entropy-12-01532}{40},\allowbreak \hyperlink{B42-entropy-12-01532}{42}]。


\(\alpha\) 散度（\hyperref[FD4-entropy-12-01532]{4}）具有以下基本性质：

\begin{enumerate}
\item 凸性：\(D_{A}^{(\alpha)}{(\mathbf{P}||\mathbf{Q})}\) 关于 \(\mathbf{P}\) 和 \(\mathbf{Q}\) 均为凸函数。


\item 严格正性：\(D_{A}^{(\alpha)}{(\mathbf{P}||\mathbf{Q})} \geq 0\)，且 \(D_{A}^{(\alpha)}{(\mathbf{P}||\mathbf{Q})} = 0\) 当且仅当 \(\mathbf{P} = \mathbf{Q}\)。


\item 连续性：\(\alpha\) 散度作为实变量 \(\alpha\) 的函数，在包括奇异参数值在内的整个取值范围上连续。


\item 对偶性：\(D_{A}^{(\alpha)}{(\mathbf{P}||\mathbf{Q})} = D_{A}^{(1 - \alpha)}{(\mathbf{Q}||\mathbf{P})}\)。


\item 排他性与包容性[\hyperlink{B42-entropy-12-01532}{42}]：

\begin{itemize}
\item 当 \(\left. \alpha\rightarrow - \infty \right.\) 时，用于逼近 \(p(\mathbf{x})\) 的 \(q(\mathbf{x})\) 的估计具有排他性，即对所有 \(\mathbf{x}\) 都有 \(q(\mathbf{x}) \leq p(\mathbf{x})\)。对 \(q(\mathbf{x})\) 最小化 \(D_{A}^{(\alpha)}{(\mathbf{P}||\mathbf{Q})}\)，将迫使 \(q(\mathbf{x})\) 成为排他性近似，即 \(q(\mathbf{x}\) 的质量位于 \(p(\mathbf{x})\) 之内；详细说明和图示见文献[\hyperlink{B42-entropy-12-01532}{42}]。


\item 当 \(\left. \alpha\rightarrow\infty \right.\) 时，用于逼近 \(p(\mathbf{x})\) 的 \(q(\mathbf{x})\) 的估计具有包容性，即对所有 \(\mathbf{x}\) 都有 \(q(\mathbf{x}) \geq p(\mathbf{x})\)。换言之，\(q(\mathbf{x})\) 的质量覆盖 \(p(\mathbf{x})\) 的全部质量。

\end{itemize}





\item 强制为零与避免为零的性质[\hyperlink{B42-entropy-12-01532}{42}]。

这里考虑 \(p(\mathbf{x})\) 与 \(q(\mathbf{x})\) 未必相互绝对连续的情形。此时，散度可能发散到无穷大，但仍具有以下两个性质：

\begin{itemize}
\item 当 \(\alpha \leq 0\) 时，用于逼近 \(p(\mathbf{x})\) 的 \(q(\mathbf{x})\) 的估计具有强制为零的性质，即 \(p(\mathbf{x}) = 0\) 将迫使 \(q(\mathbf{x}) = 0\)。


\item 当 \(\alpha \geq 1\) 时，用于逼近 \(p(\mathbf{x})\) 的 \(q(\mathbf{x})\) 的估计具有避免为零的性质，即 \(p(\mathbf{x}) > 0\) 蕴含 \(q(\mathbf{x}) > 0\)。

\end{itemize}




\end{enumerate}





\(\alpha\) 散度的一个重要性质是：它关于 \(\mathbf{P}\) 和 \(\mathbf{Q}\) 为凸函数，并在 \(\mathbf{P} = \mathbf{Q}\) 时具有唯一的极小值，例如文献[\hyperlink{B22-entropy-12-01532}{22}]。


\setcounter{section}{2}
\setcounter{subsection}{1}
\subsection{\texorpdfstring{\(\alpha\)–Rényi 散度}{α–Rényi 散度}}
\label{sec2dot2-entropy-12-01532}
\(\alpha\) 散度与 Rényi 散度密切相关。定义 \(\alpha\)–Rényi 散度如下：

\setcounter{equation}{24}
\begin{equation}
\label{FD25-entropy-12-01532}
\fitmath{\begin{array}{ccrl}
{D_{AR}^{(\alpha)}{(\mathbf{P}\,||\,\mathbf{Q})}} & & = & {\frac{1}{\alpha(\alpha - 1)}\log\left( {1 + \alpha{(\alpha - 1)}D_{A}^{(\alpha)}{(\mathbf{P}||\mathbf{Q}}} \right)} \\
 & & = & {\frac{1}{\alpha(\alpha - 1)}\log\left( {\int\left( {p^{\,\alpha}q^{1 - \alpha} - \alpha p + {(\alpha - 1)}\, q} \right)d\mu{(\mathbf{x})} + 1} \right),\;\;\alpha \in \mathbb{R} \smallsetminus {\{ 0,1\}}}
\end{array}}
\end{equation}




对于 \(\alpha = 0\) 和 \(\alpha = 1\)，\(\alpha\)–Rényi 散度分别简化为 Kullback–Leibler 散度：

\setcounter{equation}{25}
\begin{equation}
\label{FD26-entropy-12-01532}
\fitmath{\begin{array}{ccrl}
{D_{KL}{(\mathbf{P}\,||\,\mathbf{Q})} = \lim\limits_{\alpha\rightarrow 1}D_{AR}^{(\alpha)}{(\mathbf{P}\,||\,\mathbf{Q})}} & & = & {\lim\limits_{\alpha\rightarrow 1}\frac{1}{2\alpha - 1}\frac{\int(p^{\,\alpha}q^{\, 1 - \alpha}{(\log p - \log q)} - p + q)d\mu{(\mathbf{x})}}{\int(p^{\,\alpha}q^{\, 1 - \alpha} - \alpha p - {(\alpha - 1)}q)d\mu{(\mathbf{x})} + 1}} \\
 & & = & {\int\left( {p\;\log\left( \frac{p}{q} \right) - p + q} \right)d\mu{(\mathbf{x})}}
\end{array}}
\end{equation}

以及

\setcounter{equation}{26}
\begin{equation}
\label{FD27-entropy-12-01532}
\fitmath{\begin{array}{r}
{D_{KL}{(\mathbf{Q}||\mathbf{P})} = \lim\limits_{\alpha\rightarrow 0}D_{AR}^{(\alpha)}{(\mathbf{P}\,||\,\mathbf{Q})} = \int\left( {q\log\left( \frac{q}{p} \right) - q + p} \right)d\mu{(\mathbf{x})}}
\end{array}}
\end{equation}




对于归一化概率密度，定义 \(\alpha\)–Rényi 散度为

\setcounter{equation}{27}
\begin{equation}
\label{FD28-entropy-12-01532}
\fitmath{D_{AR}^{(\alpha)}(\overline{\mathbf{P}}{\,||\,}\overline{\mathbf{Q}}{) =}\frac{1}{\alpha(\alpha - 1)}\log\left( {\int{\overline{p}}^{\alpha}{\overline{q}}^{1 - \alpha}d\mu{(\mathbf{x})}} \right) = \frac{1}{\alpha - 1}\log\left( {\int\overline{q}\left( \frac{\overline{p}}{\overline{q}} \right)^{\alpha}d\mu{(\mathbf{x})}} \right)^{\frac{1}{\alpha}}}
\end{equation}

它对应于 Rényi 熵[\hyperlink{B55-entropy-12-01532}{55}--\hyperlink{B57-entropy-12-01532}{57}]，见\hyperref[app2-entropy-12-01532]{附录 B}：

\setcounter{equation}{28}
\begin{equation}
\label{FD29-entropy-12-01532}
\fitmath{H_{R}^{(\alpha)}{(\overline{\mathbf{P}})} = - \frac{1}{\alpha - 1}\log\left( {\int{\overline{p}}^{\alpha}{(\mathbf{x})}d\mu{(\mathbf{x})}} \right)}
\end{equation}




这里采用的缩放因子与原始 Rényi 散度中的因子不同[\hyperlink{B56-entropy-12-01532}{56},\allowbreak \hyperlink{B57-entropy-12-01532}{57}]。一般而言，\(\alpha\)–Rényi 散度只对归一化概率密度有意义；否则，对于某些未归一化的正测度，式（\hyperref[FD25-entropy-12-01532]{25}）给出的函数可能取复数值。


此外，对任意参数 \(\alpha \in \mathbb{R}\)，\(\alpha\) 散度关于正密度 p 和 q 为凸函数；\(\alpha\)–Rényi 散度（\hyperref[FD28-entropy-12-01532]{28}）则在 \(\alpha \in \lbrack 0,1\rbrack\) 时关于 p、q 联合为凸函数，在 \(\alpha > 1\) 时不具有凸性。凸性意味着，当 p 和 q 固定时，\(D_{AR}^{(\alpha)}\) 随 \(\alpha\) 增大。实际上，Rényi 散度在整个集合 \((0,\infty)\) 上都随 \(\alpha\) 增大[\hyperlink{B58-entropy-12-01532}{58}--\hyperlink{B60-entropy-12-01532}{60}]。


\setcounter{section}{2}
\setcounter{subsection}{2}
\subsection{\texorpdfstring{扩展的 \(\alpha\) 散度族}{扩展的 α 散度族}}
\label{sec2dot3-entropy-12-01532}
非对称 \(\alpha\) 散度有多种扩展方式。例如，可以用 q 与 p 的平均值替代 q：如果 p 和 q 相似，则二者应当都接近其平均值。据此定义修正的 \(\alpha\) 散度为

\setcounter{equation}{29}
\begin{equation}
\label{FD30-entropy-12-01532}
\fitmath{D_{Am1}^{(\alpha)}{(\mathbf{P}\,||\,}\widetilde{\mathbf{Q}}{) =}\frac{1}{\alpha(\alpha - 1)}\int\left( {\widetilde{q}\left\lbrack {\left( \frac{p}{\widetilde{q}} \right)^{\,\alpha} - 1} \right\rbrack - \alpha\;{(p - \widetilde{q})}} \right)d\mu{(\mathbf{x})}}
\end{equation}

相应的伴随 \(\alpha\) 散度为

\setcounter{equation}{30}
\begin{equation}
\label{FD31-entropy-12-01532}
\fitmath{D_{Am2}^{(\alpha)}(\widetilde{\mathbf{Q}}{\,||\,\mathbf{P})} = \frac{1}{\alpha(\alpha - 1)}\int\left( {p\left\lbrack {\left( \frac{\widetilde{q}}{p} \right)^{\alpha} - 1} \right\rbrack + \alpha\;{(p - \widetilde{q})}} \right)d\mu{(\mathbf{x})}}
\end{equation}

其中 \(\widetilde{\mathbf{Q}} = {(\mathbf{P} + \mathbf{Q})}/2\) 且 \(\widetilde{q} = {(p + q)}/2\)。


对于奇异参数值 \(\alpha = 1\) 和 \(\alpha = 0\)，\(\alpha\) 散度（\hyperref[FD30-entropy-12-01532]{30}）分别为

\[
\fitmath{\lim\limits_{\alpha\rightarrow 0}D_{Am1}^{(\alpha)}{(\mathbf{P}\,||}\widetilde{\mathbf{Q}}{) = \int}\left( {\widetilde{q}\;\log\left( \frac{\widetilde{q}}{p} \right) + p - \widetilde{q}} \right)d\mu{(\mathbf{x})} = \int\left( {\frac{p + q}{2}\;\log\left( \frac{p + q}{2p} \right) + \frac{p - q}{2}} \right)d\mu{(\mathbf{x})}}
\]

以及

\[
\fitmath{\lim\limits_{\alpha\rightarrow 1}D_{Am1}^{(\alpha)}{(\mathbf{P}\,||\,}\widetilde{Q}{) = \int}\left( {p\log\left( \frac{p}{\widetilde{q}} \right) - p + \widetilde{q}} \right)d\mu{(\mathbf{x})} = \int\left( {p\log\left( \frac{2p}{p + q} \right) + \frac{q - p}{2}} \right)d\mu{(\mathbf{x})}}
\]

下面列出式（\hyperref[FD31-entropy-12-01532]{31}）的几个重要特例。


1. 三角判别散度（Triangular Discrimination，TD）[\hyperlink{B61-entropy-12-01532}{61}]：

\setcounter{equation}{31}
\begin{equation}
\label{FD32-entropy-12-01532}
\fitmath{D_{Am2}^{( - 1)}(\widetilde{\mathbf{Q}}{\,||\,\mathbf{P})} = \frac{1}{4}D_{TD}{(\mathbf{P}\,||\,\mathbf{Q})} = \frac{1}{4}\int\frac{{(p - q)}^{2}}{p + q}d\mu{(\mathbf{x})}}
\end{equation}




2. 相对 Jensen–Shannon 散度[\hyperlink{B62-entropy-12-01532}{62}--\hyperlink{B64-entropy-12-01532}{64}]：

\setcounter{equation}{32}
\begin{equation}
\label{FD33-entropy-12-01532}
\fitmath{\lim\limits_{\alpha\rightarrow 0}D_{Am2}^{(\alpha)}(\widetilde{\mathbf{Q}}{\,||\,\mathbf{P})} = D_{RJS}{(\mathbf{P}\,||\,\mathbf{Q})} = \int\left( {p\log\left( \frac{2p}{p + q} \right) - p + q} \right)d\mu{(\mathbf{x})}}
\end{equation}




3. Taneya 提出的相对算术–几何散度[\hyperlink{B38-entropy-12-01532}{38},\allowbreak \hyperlink{B39-entropy-12-01532}{39},\allowbreak \hyperlink{B43-entropy-12-01532}{43}]：

\setcounter{equation}{33}
\begin{equation}
\label{FD34-entropy-12-01532}
\fitmath{\lim\limits_{\alpha\rightarrow 1}D_{Am2}^{(\alpha)}(\widetilde{\mathbf{Q}}{\,||\,\mathbf{P})} = \frac{1}{2}D_{RAG}{(\mathbf{P}\,||\,\mathbf{Q})} = \int\left( {{(p + q)}\log\left( \frac{p + q}{2p} \right) + p - q} \right)d\mu{(\mathbf{x})}}
\end{equation}




3. Neyman 卡方散度：

\setcounter{equation}{34}
\begin{equation}
\label{FD35-entropy-12-01532}
\fitmath{D_{Am2}^{(2)}(\widetilde{\mathbf{Q}}{\,||\,\mathbf{P})} = \frac{1}{8}D_{\chi^{2}}{(\mathbf{P}\,||\,\mathbf{Q})} = \frac{1}{8}\int\frac{{(p - q)}^{2}}{p}d\mu{(\mathbf{x})}}
\end{equation}




非对称 \(\alpha\) 散度可以形式化地表示为 Csiszár–Morimoto 的 f 散度，见\hyperref[entropy-12-01532-t001]{表 1}[\hyperlink{B30-entropy-12-01532}{30}]。


\clearpage\begin{landscape}\small\setlength{\LTcapwidth}{\linewidth}
\setcounter{table}{0}
\begin{longtable}{@{}p{0.48\linewidth}p{0.48\linewidth}@{}}
\caption[]{非对称 \(\alpha\) 散度及相应的凸 Csiszár–Morimoto 函数[\protect\hyperlink{B30-entropy-12-01532}{30}]。}\label{entropy-12-01532-t001}\\
\toprule
散度 \cellmath{D_{A}^{(\alpha)}{({\mathbf{P}||\mathbf{Q}})} = \int qf^{(\alpha)}\left( \frac{p}{q} \right)d\mu{(\mathbf{x})}} & Csiszár 函数 \cellmath{f^{(\alpha)}{(u)}, u = p/q} \\ \midrule
\cellmath{\left\{ \begin{array}{l}
{\frac{1}{\alpha(\alpha - 1)}\int\left( {q\left\lbrack {\left( \frac{p}{q} \right)^{\alpha} - 1} \right\rbrack - \alpha{(p - q)}} \right)d\mu{(\mathbf{x})},} \\
{\int\left( {q\log\frac{q}{p} + p - q} \right)d\mu{(\mathbf{x})},} \\
{\int\left( {p\log\frac{p}{q} - p + q} \right)d\mu{(\mathbf{x})}.}
\end{array} \right.} & \cellmath{\left\{ \begin{matrix}
{\frac{u^{\alpha} - 1 - \alpha{(u - 1)}}{\alpha(\alpha - 1)},\alpha \neq 0,1} \\
 \\
{u - 1 - \log u,\alpha = 0,} \\
 \\
{1 - u + u\log u,\;\alpha = 1.}
\end{matrix} \right.} \\ \midrule
\cellmath{\left\{ \begin{array}{l}
{\frac{1}{\alpha(\alpha - 1)}\int\left( {p\left\lbrack {\left( \frac{q}{p} \right)^{\alpha} - 1} \right\rbrack + \alpha{(p - q)}} \right)d\mu{(\mathbf{x})},} \\
{\int\left( {p\log\frac{p}{q} - p + q} \right)d\mu{(\mathbf{x})},} \\
{\int\left( {q\log\frac{q}{p} + p - q} \right)d\mu{(\mathbf{x})}.}
\end{array} \right.} & \cellmath{\left\{ \begin{matrix}
{\frac{u^{1 - \alpha} + {(\alpha - 1)}u - \alpha}{\alpha(\alpha - 1)},\alpha \neq 0,1,} \\
 \\
{1 - u + u\log u,\alpha = 0} \\
 \\
{u - 1 - \log u,\alpha = 1.}
\end{matrix} \right.} \\ \midrule
\cellmath{\left\{ \begin{array}{l}
{\frac{1}{\alpha(\alpha - 1)}\int\left( {q\left( \frac{p + q}{2q} \right)^{\alpha} - q - \alpha\frac{p - q}{2}} \right)d\mu{(\mathbf{x})},} \\
{\int\left( {q\log\left( \frac{2q}{p + q} \right) + \frac{p - q}{2}} \right)d\mu{(\mathbf{x})},} \\
{\int\left( {\frac{p + q}{2}\log\left( \frac{p + q}{2q} \right) - \frac{p - q}{2}} \right)d\mu{(\mathbf{x})}.}
\end{array} \right.} & \cellmath{\left\{ \begin{array}{l}
{\frac{\left( \frac{u + 1}{2} \right)^{\alpha} - 1 - \alpha\left( \frac{u - 1}{2} \right)}{\alpha(\alpha - 1)},\;\;\alpha \neq 0,1,} \\
{\frac{u - 1}{2} + \log\left( \frac{2}{u + 1} \right),\;\;\alpha = 0,} \\
{\frac{1 - u}{2} + \frac{u + 1}{2}\log\left( \frac{u + 1}{2} \right),\alpha = 1.}
\end{array} \right.} \\ \midrule
\cellmath{\left\{ \begin{array}{l}
{\frac{1}{\alpha(\alpha - 1)}\int\left( {p\left( \frac{p + q}{2p} \right)^{\alpha} - p + \alpha\frac{p - q}{2}} \right)d\mu{(\mathbf{x})},} \\
{\int\left( {p\log\left( \frac{2p}{p + q} \right) - \frac{p - q}{2}} \right)d\mu{(\mathbf{x})},} \\
{\int\left( {\frac{p + q}{2}\log\left( \frac{p + q}{2p} \right) + \frac{p - q}{2}} \right)d\mu{(\mathbf{x})}.}
\end{array} \right.} & \cellmath{\left\{ \begin{array}{l}
{\frac{u\left( \frac{u + 1}{2u} \right)^{\alpha} - u - \alpha\left( \frac{1 - u}{2} \right)}{\alpha(\alpha - 1)},\;\alpha \neq 0,1,} \\
{\frac{1 - u}{2} - u\log\left( \frac{u + 1}{2u} \right),\;\alpha = 0,} \\
{\frac{u - 1}{2} + \left( \frac{u + 1}{2} \right)\log\left( \frac{u + 1}{2u} \right),\alpha = 1.}
\end{array} \right.} \\ \midrule
\cellmath{\left\{ \begin{array}{l}
{\frac{1}{\alpha - 1}\int{(p - q)}\left\lbrack {\left( \frac{p + q}{2q} \right)^{\alpha - 1} - 1} \right\rbrack d\mu{(\mathbf{x})},} \\
{\int{(p - q)}\log\left( \frac{p + q}{2q} \right)d\mu{(\mathbf{x})}.}
\end{array} \right.} & \cellmath{\left\{ \begin{array}{l}
{\frac{{(u - 1)}\left\lbrack {\left( \frac{u + 1}{2} \right)^{\alpha - 1} - 1} \right\rbrack}{\alpha - 1},\alpha \neq 1,} \\
{{(u - 1)}\log\left( \frac{u + 1}{2} \right),\alpha = 1.}
\end{array} \right.} \\ \midrule
\cellmath{\left\{ \begin{array}{l}
{\frac{1}{\alpha - 1}\int{(q - p)}\left\lbrack {\left( \frac{p + q}{2p} \right)^{\alpha - 1} - 1} \right\rbrack d\mu{(\mathbf{x})},} \\
{\int{(q - p)}\log\left( \frac{p + q}{2p} \right)d\mu{(\mathbf{x})}.}
\end{array} \right.} & \cellmath{\left\{ \begin{array}{l}
{\frac{{(1 - u)}\left\lbrack {\left( \frac{u + 1}{2u} \right)^{\alpha - 1} - 1} \right\rbrack}{\alpha - 1},\alpha \neq 1,} \\
{{(1 - u)}\log\left( \frac{u + 1}{2u} \right),\alpha = 1.}
\end{array} \right.} \\ \midrule
\bottomrule\end{longtable}\end{landscape}\clearpage
\setcounter{section}{2}
\setcounter{subsection}{3}
\subsection{\texorpdfstring{对称化的 \(\alpha\) 散度}{对称化的 α 散度}}
\label{sec2dot4-entropy-12-01532}
基本的 \(\alpha\) 散度是非对称的，即 \(D_{A}^{(\alpha)}{(\mathbf{P}\,||\,\mathbf{Q})} \neq D_{A}^{(\alpha)}{(\mathbf{Q}\,||\,\mathbf{P})}\)。


通常有两种散度对称化方法。第一类为

\setcounter{equation}{35}
\begin{equation}
\label{FD36-entropy-12-01532}
\fitmath{D_{S1}{(\mathbf{P}\,||\,\mathbf{Q})} = \frac{1}{2}\left\lbrack {D_{A}{(\mathbf{P}\,||\,\mathbf{Q})} + D_{A}{(\mathbf{Q}\,||\,\mathbf{P})}} \right\rbrack}
\end{equation}

第二类为

\setcounter{equation}{36}
\begin{equation}
\label{FD37-entropy-12-01532}
\fitmath{D_{S2}{(\mathbf{P}\,||\,\mathbf{Q})} = \frac{1}{2}\left\lbrack {D_{A}\left( {{\mathbf{P}\,||\,}\frac{\mathbf{P} + \mathbf{Q}}{2}} \right) + D_{A}\left( {{\mathbf{Q}\,||\,}\frac{\mathbf{P} + \mathbf{Q}}{2}} \right)} \right\rbrack}
\end{equation}




第一类对称 \(\alpha\) 散度可定义为下式。为简化表达，省略缩放因子 \(1/2\)：

\setcounter{equation}{37}
\begin{equation}
\label{FD38-entropy-12-01532}
\fitmath{D_{AS1}^{(\alpha)}{(\mathbf{P}\,||\,\mathbf{Q})} = D_{A}^{(\alpha)}{(\mathbf{P}\,||\,\mathbf{Q})} + D_{A}^{(\alpha)}{(\mathbf{Q}\,||\,\mathbf{P})} = \int\frac{\left( {p^{\alpha} - q^{\alpha}} \right)\left( {p^{1 - \alpha} - q^{1 - \alpha}} \right)}{\alpha(1 - \alpha)}\, d\mu{(\mathbf{x})}}
\end{equation}

作为特例，可以得到几个已知的对称散度。


1. 对称卡方散度[\hyperlink{B54-entropy-12-01532}{54}]：

\setcounter{equation}{38}
\begin{equation}
\label{FD39-entropy-12-01532}
\fitmath{D_{AS1}^{( - 1)}{(\mathbf{P}\,||\,\mathbf{Q})} = D_{AS1}^{(2)}{(\mathbf{P}\,||\,\mathbf{Q})} = \frac{1}{2}D_{\chi^{2}}{(\mathbf{P}\,||\,\mathbf{Q})} = \frac{1}{2}\int\frac{{(p - q)}^{2}{(p + q)}}{pq}d\mu{(\mathbf{x})}}
\end{equation}




2. 对称化的 KL 散度，也称 J 散度，对应于 Jeffreys 熵最大化[\hyperlink{B65-entropy-12-01532}{65},\allowbreak \hyperlink{B66-entropy-12-01532}{66}]：

\setcounter{equation}{39}
\begin{equation}
\label{FD40-entropy-12-01532}
\fitmath{\lim\limits_{\alpha\rightarrow 0}D_{AS1}^{(\alpha)}{(\mathbf{P}\,||\,\mathbf{Q})} = \lim\limits_{\alpha\rightarrow 1}D_{AS1}^{(\alpha)}{(\mathbf{P}\,||\,\mathbf{Q})} = D_{J}{(\mathbf{P}\,||\,\mathbf{Q})} = \int{(p - q)}\log\left( \frac{p}{q} \right)d\mu{(\mathbf{x})}}
\end{equation}




另一类广泛的对称散度，可以由第二类对称 \(\alpha\) 散度表示：

\[
\fitmath{\begin{array}{ccrl}
{D_{AS2}^{(\alpha)}{(\mathbf{P}\,||\,\mathbf{Q})}} & & = & {D_{A}^{(\alpha)}\left( {{\mathbf{P}\,||\,}\frac{\mathbf{P} + \mathbf{Q}}{2}} \right) + D_{A}^{(\alpha)}\left( {{\mathbf{Q}\,||\,}\frac{\mathbf{P} + \mathbf{Q}}{2}} \right)} \\
 & & = & {\frac{1}{\alpha(\alpha - 1)}\int\left( {{(p^{1 - \alpha} + q^{1 - \alpha})}\left( \frac{p + q}{2} \right)^{\alpha} - {(p + q)}} \right)d\mu{(\mathbf{x})}}
\end{array}}
\]




上述度量包括以下重要特例：

\begin{enumerate}
\item 三角判别散度[\hyperlink{B30-entropy-12-01532}{30},\allowbreak \hyperlink{B38-entropy-12-01532}{38}]：

\setcounter{equation}{40}
\begin{equation}
\label{FD42-entropy-12-01532}
\fitmath{D_{AS2}^{( - 1)}{(\mathbf{P}\,||\,\mathbf{Q})} = \frac{1}{2}D_{TD}{(\mathbf{P}\,||\,\mathbf{Q})} = \frac{1}{2}\int\frac{{(p - q)}^{2}}{p + q}d\mu{(\mathbf{x})}}
\end{equation}




\item 对称 Jensen–Shannon 散度[\hyperlink{B62-entropy-12-01532}{62},\allowbreak \hyperlink{B64-entropy-12-01532}{64}]：

\setcounter{equation}{41}
\begin{equation}
\label{FD43-entropy-12-01532}
\fitmath{\lim\limits_{\alpha\rightarrow 0}D_{AS2}^{(\alpha)}{(\mathbf{P}\,||\,\mathbf{Q})} = D_{JS}{(\mathbf{P}\,||\,\mathbf{Q})} = \int\left( {p\log\left( \frac{2p}{p + q} \right) + q\log\left( \frac{2q}{p + q} \right)} \right)d\mu{(\mathbf{x})}}
\end{equation}

Jensen–Shannon 散度是 Kullback–Leibler 散度经过对称化和平滑化得到的变体，可以解释为两个分布分别到其平均分布的 KL 散度的平均。对于归一化概率密度，Jensen–Shannon 散度与 Shannon 熵具有以下关系：

\setcounter{equation}{42}
\begin{equation}
\label{FD44-entropy-12-01532}
\fitmath{D_{JS} = H_{S}{({(\mathbf{P} + \mathbf{Q})}/2)} - {(H_{S}{(\mathbf{P})} + H_{S}{(\mathbf{Q})})}/2}
\end{equation}

其中 \(H_{S}{(\mathbf{P})} = - \int p{(\mathbf{x})}\log p{(\mathbf{x})}d\mu{(\mathbf{x})}\)。


\item 算术–几何散度[\hyperlink{B39-entropy-12-01532}{39}]：

\setcounter{equation}{43}
\begin{equation}
\label{FD45-entropy-12-01532}
\fitmath{\lim\limits_{\alpha\rightarrow 1}D_{AS2}^{(\alpha)}{(\mathbf{P}\,||\,\mathbf{Q})} = \int{(p + q)}\log\left( \frac{p + q}{2\sqrt{pq}} \right)d\mu{(\mathbf{x})}}
\end{equation}




\item 对称卡方散度[\hyperlink{B54-entropy-12-01532}{54}]：

\setcounter{equation}{44}
\begin{equation}
\label{FD46-entropy-12-01532}
\fitmath{D_{AS2}^{(2)}{(\mathbf{P}\,||\,\mathbf{Q})} = \frac{1}{8}D_{\chi^{2}}{(\mathbf{P}\,||\,\mathbf{Q})} = \frac{1}{8}\int\frac{{(p - q)}^{2}{(p + q)}}{pq}d\mu{(\mathbf{x})}}
\end{equation}



\end{enumerate}





上述 \(\alpha\) 散度关于参数 \(\mathbf{P}\) 和 \(\mathbf{Q}\) 对称；即使 \(\mathbf{P}\) 和 \(\mathbf{Q}\) 不相互绝对连续，它也具有良好定义。例如，在离散情形中，即使某些指标 \(p_{i}\) 上，一个分量为零而 \(q_{i}\) 不为零，或者 \(q_{i}\) 为零而 \(p_{i}\) 不为零，\(D_{AS2}^{(\alpha)}{(\mathbf{P}||\mathbf{Q})}\) 仍有良好定义[\hyperlink{B54-entropy-12-01532}{54}]。它还有下界和上界；例如，Jensen–Shannon 散度的取值位于 0 与 2 之间[\hyperlink{B36-entropy-12-01532}{36}]。

\setcounter{section}{2}
\section{\texorpdfstring{\(\beta\) 散度族}{β 散度族}}
\label{sec3-entropy-12-01532}
基本的 \(\beta\) 散度由 Basu 等[\hyperlink{B67-entropy-12-01532}{67}]以及 Minami 和 Eguchi[\hyperlink{B15-entropy-12-01532}{15}]提出，许多研究者研究了它的应用，包括文献[\hyperlink{B8-entropy-12-01532}{8},\allowbreak \hyperlink{B13-entropy-12-01532}{13},\allowbreak \hyperlink{B30-entropy-12-01532}{30}--\hyperlink{B34-entropy-12-01532}{34},\allowbreak \hyperlink{B37-entropy-12-01532}{37},\allowbreak \hyperlink{B37-entropy-12-01532}{37},\allowbreak \hyperlink{B68-entropy-12-01532}{68}--\hyperlink{B72-entropy-12-01532}{72}]及其中的参考文献。从应用角度，研究 \(\beta\) 散度的主要动机，是为聚类、特征提取、分类和盲源分离开发对离群值高度稳健的学习算法。\(\beta\) 散度已成功用于稳健主成分分析（Principal Component Analysis，PCA）与聚类[\hyperlink{B71-entropy-12-01532}{71}]、稳健独立成分分析（Independent Component Analysis，ICA）[\hyperlink{B15-entropy-12-01532}{15},\allowbreak \hyperlink{B68-entropy-12-01532}{68},\allowbreak \hyperlink{B69-entropy-12-01532}{69}]，以及稳健 NMF 和非负张量分解（Non-negative Tensor Factorization，NTF）[\hyperlink{B30-entropy-12-01532}{30},\allowbreak \hyperlink{B70-entropy-12-01532}{70},\allowbreak \hyperlink{B73-entropy-12-01532}{73}--\hyperlink{B76-entropy-12-01532}{76}]。


先定义两个未归一化密度函数之间的基本非对称 \(\beta\) 散度：

\setcounter{equation}{45}
\begin{equation}
\label{FD47-entropy-12-01532}
\fitmath{D_{B}^{(\beta)}{(\mathbf{P}\,||\;\mathbf{Q})} = \int\left( {p{(\mathbf{x})}\;\frac{p^{\beta - 1}{(\mathbf{x})} - q^{\beta - 1}{(\mathbf{x})}}{\beta - 1} - \frac{p^{\beta}{(\mathbf{x})} - q^{\beta}{(\mathbf{x})}}{\beta}} \right)d\mu{(\mathbf{x})},\;\;\beta \in \mathbb{R} \smallsetminus {\{ 0,1\}}}
\end{equation}

其中 \(\beta\) 为实数；在 \(\beta = 0,1\) 时，通过连续性定义，详见下文。


对于概率质量函数为 \(\mathbf{P} = \lbrack p_{1},p_{2},\ldots,p_{n}\rbrack\) 和 \(\mathbf{Q} = \lbrack q_{1},q_{2},\ldots,q_{n}\rbrack\) 的离散概率测度，离散 \(\beta\) 散度定义为

\setcounter{equation}{46}
\begin{equation}
\label{FD48-entropy-12-01532}
\fitmath{D_{B}^{(\beta)}{(\mathbf{P}\,||\;\mathbf{Q})} = \sum\limits_{i = 1}^{n}d_{B}^{\;(\beta)}(p_{i}{\,||\,}q_{i}{) =}\sum\limits_{i = 1}^{n}\left( {p_{i}\;\frac{p_{i}^{\beta - 1} - q_{i}^{\beta - 1}}{\beta - 1} - \frac{p_{i}^{\beta} - q_{i}^{\beta}}{\beta}} \right)\;\;\beta \in \mathbb{R} \smallsetminus {\{ 0,1\}}}
\end{equation}

\(\beta\) 散度也可通过广义 KL 散度表示：

\setcounter{equation}{47}
\begin{equation}
\label{FD49-entropy-12-01532}
\fitmath{D_{B}^{(\beta)}{(\mathbf{P}\,||\;\mathbf{Q})} = D_{GKL}^{(\beta^{- 1})}(\mathbf{P}^{\beta}{\,||\;}\mathbf{Q}^{\beta}{) = -}\frac{1}{\beta}\int\left( {p^{\beta}\log_{(\frac{1}{\beta})}\left( \frac{q^{\beta}}{p^{\beta}} \right) + p^{\beta} - q^{\beta}} \right)d\mu{(\mathbf{x})}\;\;\beta \in \mathbb{R} \smallsetminus {\{ 0,1\}}}
\end{equation}

这一表示说明，在某些调节参数 \(\beta\) 的取值下，\(\beta\) 散度对离群值具有稳健性，因而更适合某些应用。例如，在声音处理中，语音功率谱可以用指数族密度建模；当 \(\beta = 0\) 时，\(\beta\) 散度正是 Itakura–Saito 距离，也称 Itakura–Saito 散度、Itakura–Saito 失真度量或 Burg 交叉熵[\hyperlink{B12-entropy-12-01532}{12},\allowbreak \hyperlink{B13-entropy-12-01532}{13},\allowbreak \hyperlink{B30-entropy-12-01532}{30},\allowbreak \hyperlink{B76-entropy-12-01532}{76}--\hyperlink{B79-entropy-12-01532}{79}]。具体地，\(\beta\) 散度在 \(\left. \beta\rightarrow 0 \right.\) 的极限情形下，必须定义为 Itakura–Saito 距离：

\setcounter{equation}{48}
\begin{equation}
\label{FD50-entropy-12-01532}
\fitmath{D_{IS}{(\mathbf{P}\,||\,\mathbf{Q})} = \lim\limits_{\beta\rightarrow 0}D_{B}^{(\beta)}{(\mathbf{P}\,||\,\mathbf{Q})} = \int\left( {\log\frac{q}{p} + \frac{p}{q} - 1} \right)d\mu{(\mathbf{x})}}
\end{equation}

Itakura 和 Saito 的距离源于语音谱的最大似然（Maximum Likelihood，ML）估计[\hyperlink{B77-entropy-12-01532}{77}]，用于衡量两个频谱之间的失真或拟合优度。它具有尺度不变性，重建信号的感知效果较好，因此常被语音处理领域用作标准度量。尺度不变性使 p 的低能量分量与高能量分量具有相同的相对重要性。当 p 的系数具有很大动态范围时，例如短时音频谱，这一点尤其重要[\hyperlink{B30-entropy-12-01532}{30},\allowbreak \hyperlink{B76-entropy-12-01532}{76},\allowbreak \hyperlink{B79-entropy-12-01532}{79}]。


当 \(\beta = 2\) 时，得到标准的平方欧氏距离，即 \(L_{2}\) 范数距离；对于奇异情形 \(\beta = 1\)，则得到 KL 的 I 散度：

\setcounter{equation}{49}
\begin{equation}
\label{FD51-entropy-12-01532}
\fitmath{D_{KL}{(\mathbf{P}\,||\,\mathbf{Q})} = \lim\limits_{\beta\rightarrow 1}D_{B}^{(\beta)}{(\mathbf{P}\,||\,\mathbf{Q})} = \int\left( {p\;\log\frac{p}{q} - p + q} \right)d\mu{(\mathbf{x})}}
\end{equation}

这里使用了恒等式 \(\lim_{\beta\rightarrow 0}\frac{p^{\beta} - q^{\beta}}{\beta} = \log{(p/q)}\) 和 \(\lim_{\beta\rightarrow 0}\frac{p^{\beta} - 1}{\beta} = \log p\)。


因此，\(\beta\) 散度可以更明确地表示为

\setcounter{equation}{50}
\begin{equation}
\label{FD52-entropy-12-01532}
\fitmath{D_{B}^{(\beta)}{(\mathbf{P}\,||\,\mathbf{Q})} = \left\{ \begin{array}{l}
{\frac{1}{\beta(\beta - 1)}\int\left( {p^{\beta}{(\mathbf{x})} + {(\beta - 1)}q^{\beta}{(\mathbf{x})} - \beta p{(\mathbf{x})}q^{\beta - 1}{(\mathbf{x})}} \right)d\mu{(\mathbf{x})},\beta \neq 0,1} \\
{\int\left( {p{(\mathbf{x})}\;\log{(\frac{p(\mathbf{x})}{q(\mathbf{x})})} - p{(\mathbf{x})} + q{(\mathbf{x})}} \right)d\mu{(\mathbf{x})},\;\;\beta = 1} \\
{\int\left( {\log{(\frac{q(\mathbf{x})}{p(\mathbf{x})})} + \frac{p(\mathbf{x})}{q(\mathbf{x})} - 1} \right)d\mu{(\mathbf{x})},\;\;\;\beta = 0}
\end{array} \right.}
\end{equation}

基本的 \(\beta\) 散度平滑地连接 Itakura–Saito 距离和平方欧氏 \(L_{2}\) 范数距离，并经过 KL 的 I 散度 \(D_{KL}{(\mathbf{P}\,||\,\mathbf{Q})}\)。\(\alpha\) 散度族无法给出这样的参数化连接。


调节参数 \(\beta\) 的选择取决于数据集的统计分布。例如，对正态分布，最优参数为 \(\beta = 2\)；对伽马分布，为 \(\beta = 0\)；对泊松分布，为 \(\beta = 1\)；对复合泊松分布，为 \(\beta \in (1,2){}\)[\hyperlink{B15-entropy-12-01532}{15},\allowbreak \hyperlink{B31-entropy-12-01532}{31}--\hyperlink{B34-entropy-12-01532}{34},\allowbreak \hyperlink{B68-entropy-12-01532}{68},\allowbreak \hyperlink{B69-entropy-12-01532}{69}]。


\(\beta\) 散度可以从 Bregman 散度导出。Bregman 散度是衡量两个密度函数值 \(p(\mathbf{x})\) 和 \(q(\mathbf{x})\) 之间差异的伪距离[\hyperlink{B9-entropy-12-01532}{9},\allowbreak \hyperlink{B16-entropy-12-01532}{16},\allowbreak \hyperlink{B80-entropy-12-01532}{80}]：

\setcounter{equation}{51}
\begin{equation}
\label{FD53-entropy-12-01532}
\fitmath{d_{\Phi}{(p||q)} = \Phi{(p)} - \Phi{(q)} - {(p - q)}\Phi'{(q)}}
\end{equation}

其中 \(\Phi(t)\) 为严格凸的实值函数，\(\Phi'{(q)}\) 是关于 q 的导数。两个函数 \(p(\mathbf{x})\) 和 \(q(\mathbf{x})\) 之间的总差异为

\setcounter{equation}{52}
\begin{equation}
\label{FD54-entropy-12-01532}
\fitmath{D_{\Phi}{(\mathbf{P}||\mathbf{Q})} = \int\left\lbrack {\Phi{(p{(\mathbf{x})})} - \Phi{(q{(\mathbf{x})})} - {(p{(\mathbf{x})} - q{(\mathbf{x})})}\Phi'{(q{(\mathbf{x})})}} \right\rbrack d\mu{(\mathbf{x})}}
\end{equation}

它对应于连续概率测度 \(p(\mathbf{x}) \geq 0\) 的 \(\Phi\) 熵，定义为

\setcounter{equation}{53}
\begin{equation}
\label{FD55-entropy-12-01532}
\fitmath{H_{\Phi}{(\mathbf{P})} = - \int\Phi{(p{(x)})}d\mu{(\mathbf{x})}}
\end{equation}




说明：散度与熵的概念密切相关。设 \(\mathbf{Q}_{0}\) 为均匀分布，满足

\setcounter{equation}{54}
\begin{equation}
\label{FD56-entropy-12-01532}
\fitmath{\mathbf{Q}_{0}{(\mathbf{x})} = const}
\end{equation}

当 \(\mathbf{x}\) 是无限空间 \(\mathbb{R}^{n}\) 时，它可能不是概率分布，但仍是一个测度。于是

\setcounter{equation}{55}
\begin{equation}
\label{FD57-entropy-12-01532}
\fitmath{H{(\mathbf{P})} = - D(\mathbf{P}||\mathbf{Q}_{0}) + const}
\end{equation}

可视为相应的熵，即从 \(\mathbf{P}\) 到均匀分布的散度的负值。反过来，给定凹的熵函数 \(H(\mathbf{P})\)，可由凸函数 \(\Phi(\mathbf{P}) = - H(\mathbf{P})\) 导出 Bregman 散度，作为相应的散度。


若 \(\mathbf{x}\) 在某个空间上取离散值，可分离 Bregman 散度定义为
\[
\fitmath{D_{\Phi}{(\mathbf{P}||\mathbf{Q})} = \sum_{i = 1}d_{\Phi}(p_{i}{||}q_{i}{) =}\sum_{i = 1}^{n}\left\lbrack {\Phi{(p_{i})} - \Phi{(q_{i})} - {(p_{i} - q_{i})}\Phi'{(q_{i})})} \right\rbrack}
\]
，其中 \(\Phi'{(q)}\) 表示关于 q 的导数。对于一般的不可分离情形，两个向量 \(\mathbf{P}\) 和 \(\mathbf{Q}\) 之间的 Bregman 散度定义为
\[
\fitmath{D_{\Phi}{(\mathbf{P}||\mathbf{Q})} = \Phi{(\mathbf{P})} - \Phi{(\mathbf{Q})} - {(\mathbf{P} - \mathbf{Q})}^{T}\nabla\Phi{(\mathbf{Q})}}
\]
，其中 \(\nabla\Phi(\mathbf{Q})\) 为 \(\Phi\) 在 \(\mathbf{Q}\) 处的梯度。


\(D_{\Phi}{(\mathbf{P}||\mathbf{Q})}\) 等于 \(\Phi(\mathbf{P})\) 在 \(\mathbf{Q}\) 处一阶 Taylor 展开的余项。Bregman 散度包含许多重要的相异性度量，例如平方欧氏距离、Mahalanobis 距离、广义 Kullback–Leibler 散度和 Itakura–Saito 距离。


取如下严格凸的连续函数，即可从 Bregman 散度生成 \(\beta\) 散度[\hyperlink{B36-entropy-12-01532}{36},\allowbreak \hyperlink{B81-entropy-12-01532}{81}]：

\setcounter{equation}{56}
\begin{equation}
\label{FD58-entropy-12-01532}
\fitmath{\Phi{(t)} = \left\{ \begin{array}{l}
{\frac{1}{\beta(\beta - 1)}\left( {t^{\beta} - \beta t + \beta - 1} \right),\beta \neq 0,1} \\
{t\log(t) - t + 1,\beta = 1} \\
{t - \log(t) - 1,\;\;\;\beta = 0}
\end{array} \right.}
\end{equation}

同一个生成函数 \(f{(u)} = \Phi{(t)}_{|t = u}\)，在 \(\alpha = \beta\) 的条件下，也可以通过 Csiszár–Morimoto 的 f 散度 \(D_{f}{(\mathbf{P}||\mathbf{Q})} = \int q\, f{(p/q)}\, d\mu{(\mathbf{x})}\) 生成 \(\alpha\) 散度。


\(\beta\) 散度还可以由一种广义 f 散度生成：

\setcounter{equation}{57}
\begin{equation}
\label{FD59-entropy-12-01532}
\fitmath{{\widetilde{D}}_{f}{(\mathbf{P}||\mathbf{Q})} = \int\, q^{\beta}f{(p/q)}\, d\mu{(\mathbf{x})} = \int\, q^{\beta}\widetilde{f}\left( \frac{p^{\beta}}{q^{\beta}} \right)\, d\mu{(\mathbf{x})}}
\end{equation}

其中 \(f{(u)} = \Phi{(t)}_{|t = u}\)，并且 \(u = p/q\)；函数

\setcounter{equation}{58}
\begin{equation}
\label{FD60-entropy-12-01532}
\fitmath{\widetilde{f}{(\widetilde{u})} = \frac{1}{1 - \beta}\left( {{\widetilde{u}}^{\frac{1}{\beta}} - \frac{1}{\beta}\widetilde{u} + \frac{1}{\beta} - 1} \right)}
\end{equation}

是凸生成函数，满足 \(\widetilde{u} = p^{\beta}/q^{\beta}\)。


Bregman 散度与 \(\beta\) 散度之间的联系很重要，因为 Bregman 散度的许多基本性质同样适用于 \(\beta\) 散度[\hyperlink{B28-entropy-12-01532}{28},\allowbreak \hyperlink{B82-entropy-12-01532}{82}]：

\begin{enumerate}
\item 凸性：Bregman 散度 \(D_{\Phi}{(\mathbf{P}||\mathbf{Q})}\) 关于第一个变量 \(\mathbf{P}\) 始终为凸函数，但关于第二个变量 \(\mathbf{Q}\) 常常不是凸函数。


\item 非负性：Bregman 散度非负，即 \(D_{\Phi}{(\mathbf{P}||\mathbf{Q})} \geq 0\)，其零点满足 \(\mathbf{P} = \mathbf{Q}\)。


\item 线性：Bregman 散度的任意正系数线性组合仍为 Bregman 散度，即

\setcounter{equation}{59}
\begin{equation}
\label{FD61-entropy-12-01532}
\fitmath{D_{c_{1}\Phi_{1} + c_{2}\Phi_{2}}{(\mathbf{P}||\mathbf{Q})} = c_{1}D_{\Phi_{1}}{(\mathbf{P}||\mathbf{Q})} + c_{2}D_{\Phi_{2}}{(\mathbf{P}||\mathbf{Q})}}
\end{equation}

其中 \(c_{1},c_{2}\) 为正常数，\(\Phi_{1},\Phi_{2}\) 为严格凸函数。


\item 不变性：对于正测度 \(\mathbf{P}\) 和 \(\mathbf{Q}\)，泛函形式的 Bregman 散度在给生成函数加入线性项和任意常数项的仿射变换 \(\Gamma{(\mathbf{Q})} = \Phi{(\mathbf{Q})} + \int a{(x - x')}q{(x')}dx' + c\) 下保持不变[\hyperlink{B28-entropy-12-01532}{28},\allowbreak \hyperlink{B82-entropy-12-01532}{82}]，即

\setcounter{equation}{60}
\begin{equation}
\label{FD62-entropy-12-01532}
\fitmath{D_{\Gamma}{(\mathbf{P}||\mathbf{Q})} = D_{\Phi}{(\mathbf{P}||\mathbf{Q})}}
\end{equation}




\item 三点性质推广了余弦定理：

\setcounter{equation}{61}
\begin{equation}
\label{FD63-entropy-12-01532}
\fitmath{D_{\Phi}{(\mathbf{P}\,||\,\mathbf{Q})} = D_{\Phi}{(\mathbf{P}\,||\,\mathbf{Z})} + D_{\Phi}{(\mathbf{Z}\,||\,\mathbf{Q})} - {(\mathbf{P} - \mathbf{Z})}^{T}{(\frac{\delta}{\delta\mathbf{Q}}\Phi{(\mathbf{Q})} - \frac{\delta}{\delta\mathbf{Z}}\Phi{(\mathbf{Z})})}}
\end{equation}




\item 广义勾股定理：

\setcounter{equation}{62}
\begin{equation}
\label{FD64-entropy-12-01532}
\fitmath{D_{\Phi}{(\mathbf{P}\,||\,\mathbf{Q})} \geq D_{\Phi}{(\mathbf{P}\,||\,}P_{\Omega}{{(\mathbf{Q})}) +}D_{\Phi}(P_{\Omega}{{(\mathbf{Q})}\,||\,\mathbf{Q})}}
\end{equation}

其中 \(P_{\Omega}{(\mathbf{Q})} = \underset{\mathbf{\omega} \in \Omega}{\arg\min}D_{\Phi}{(\mathbf{\omega}||\mathbf{Q})}\) 是到凸集 \(\Omega\) 的 Bregman 投影，并且 \(\mathbf{P} \in \Omega\)。若 \(\Omega\) 是仿射集，则等号成立。从信息几何的角度，这就是广义勾股关系。

\end{enumerate}





\(\beta\) 散度（\hyperref[FD47-entropy-12-01532]{46}）关于 \(\mathbf{Q}\) 的一阶和二阶 Fréchet 导数为[\hyperlink{B28-entropy-12-01532}{28},\allowbreak \hyperlink{B76-entropy-12-01532}{76}]：

\setcounter{equation}{63}
\begin{equation}
\label{FD65-entropy-12-01532}
\fitmath{\frac{\delta D_{B}^{(\beta)}}{\delta q} = q^{\beta - 2}{(q - p)},\text{\quad\quad}\frac{\delta^{2}D_{B}^{(\beta)}}{\delta q^{2}} = q^{\beta - 3}{({(\beta - 1)}q - {(\beta - 2)}p)}}
\end{equation}

因此，\(\beta\) 散度在 \(\mathbf{P} = \mathbf{Q}\) 时具有唯一的全局极小值零，并随 \(|p - q|\) 增大。对于 \(q(\mathbf{x}) > 0\)，\(\beta\) 散度仅在 \(\beta \in \lbrack 1,2\rbrack\) 时严格凸。当 \(\beta = 0\)，即 Itakura–Saito 距离的情形，它在 \(\int q^{- 3}{(2p - q)}d\mu{(\mathbf{x})} \geq 0\)，也就是 \(q/p \leq 2\) 时为凸函数[\hyperlink{B78-entropy-12-01532}{78}]。


\setcounter{section}{3}
\setcounter{subsection}{0}
\subsection{\texorpdfstring{从 \(\alpha\) 散度族直接生成 \(\beta\) 散度族}{从 α 散度族直接生成 β 散度族}}
\label{sec3dot1-entropy-12-01532}
原始工作[\hyperlink{B15-entropy-12-01532}{15},\allowbreak \hyperlink{B67-entropy-12-01532}{67}--\hyperlink{B69-entropy-12-01532}{69}]只讨论了 \(\beta \geq 1\) 时的 \(\beta\) 散度函数，没有覆盖参数 \(\beta\) 的全部非正取值，尤其没有讨论 \(\beta = 0\) 时重要的 Itakura–Saito 距离。此外，与 \(\alpha\) 散度类似，还存在相应的 \(\beta\) 散度族，其中包含一族广义 Itakura–Saito 型距离。由此需要说明如何生成整个 \(\beta\) 散度族，以及 \(\alpha\) 散度与 \(\beta\) 散度之间的联系和对应关系。根据前面的讨论，可以从 \(\alpha\) 散度得到完整的 \(\beta\) 散度族；反过来，也可以直接从 \(\beta\) 散度得到 \(\alpha\) 散度。


要从一个 \(\alpha\) 散度得到与之对应的 \(\beta\) 散度，需要采用以下非线性变换：

\setcounter{equation}{64}
\begin{equation}
\label{FD66-entropy-12-01532}
\fitmath{\left. p\rightarrow p^{\beta}, q\rightarrow q^{\beta}\text{且}\alpha = \beta^{- 1} \right.}
\end{equation}




例如，对基本的非对称 \(\alpha\) 散度（\hyperref[FD4-entropy-12-01532]{4}）作上述变换或代换，并假设 \(\alpha = {(\beta)}^{- 1}\)，得到

\setcounter{equation}{65}
\begin{equation}
\label{FD67-entropy-12-01532}
\fitmath{D_{A}^{(\beta)}{(\mathbf{P}\,||\,\mathbf{Q})} = \beta^{2}\int\left( \frac{- \beta p\, q^{\beta - 1} + p^{\beta} + {(\beta - 1)}\, q^{\beta}}{\beta(\beta - 1)} \right)d\mu{(\mathbf{x})}}
\end{equation}




忽略缩放因子 \(\beta^{2}\) 后，就得到式（\hyperref[FD47-entropy-12-01532]{46}）定义的基本非对称 \(\beta\) 散度。


整个 \(\alpha\) 散度族与 \(\beta\) 散度族之间都存在这种联系，见\hyperref[entropy-12-01532-t002]{表 2}。


例如，可以从第一类对称 \(\alpha\) 散度（\hyperref[FD38-entropy-12-01532]{38}）导出对称 \(\beta\) 散度：

\[
\fitmath{\begin{array}{clr}
{D_{BS1}^{(\beta)}{(\mathbf{P}\,||\,\mathbf{Q})}} & = & {D_{B}^{(\beta)}{(\mathbf{P}\,||\,\mathbf{Q})} + D_{B}^{(\beta)}{(\mathbf{Q}\,||\,\mathbf{P})}} \\
 & = & {\frac{1}{\beta - 1}\int\left( {{(p - q)}{(p^{\beta - 1} - q^{\beta - 1})}} \right)d\mu{(\mathbf{x})}}
\end{array}}
\]




在若干特例中，可以得到以下结果。


对称 KL 散度，即 J 散度[\hyperlink{B65-entropy-12-01532}{65}]：

\setcounter{equation}{66}
\begin{equation}
\label{FD69-entropy-12-01532}
\fitmath{D_{BS1}^{(1)} = \lim\limits_{\beta\rightarrow 1}D_{BS1}^{(\beta)} = \int{(p - q)}\log\left( \frac{p}{q} \right)d\mu{(\mathbf{x})}}
\end{equation}

以及对称卡方散度[\hyperlink{B54-entropy-12-01532}{54}]：

\setcounter{equation}{67}
\begin{equation}
\label{FD70-entropy-12-01532}
\fitmath{D_{BS1}^{(0)} = \lim\limits_{\beta\rightarrow 0}D_{BS1}^{(\beta)} = \int\frac{{(p - q)}^{2}}{pq}d\mu{(\mathbf{x})}}
\end{equation}




类似地，从第二类对称 \(\alpha\) 散度可得

\setcounter{equation}{68}
\begin{equation}
\label{FD71-entropy-12-01532}
\fitmath{D_{BS2}^{(\beta)}{(\mathbf{P}\,||\,\mathbf{Q})} = \frac{1}{\beta - 1}\int\left( {{p^{\beta} + q^{\beta}} - {(p^{\beta - 1} + q^{\beta - 1})}\left( \frac{p^{\beta} + q^{\beta}}{2} \right)^{\frac{1}{\beta}}} \right)d\mu{(\mathbf{x})}}
\end{equation}



\clearpage\begin{landscape}\small\setlength{\LTcapwidth}{\linewidth}
\setcounter{table}{1}
\begin{longtable}{@{}p{0.48\linewidth}p{0.48\linewidth}@{}}
\caption[]{\(\alpha\) 散度族及其对应的 \(\beta\) 散度族。使用的变换为 \(\left. p\rightarrow p^{\beta},\;\; q\rightarrow q^{\beta},\;\;\alpha = 1/\beta \right.\)。其中 \(D_{A}^{(1)}{({\mathbf{P}||\mathbf{Q}})} = D_{B}^{(1)}{({\mathbf{P}||\mathbf{Q}})}\)，它们在 \(\alpha = \beta = 1\) 时给出扩展的 KL 散度族；在 \(\beta = 0\) 时，\(\beta\) 散度给出广义的、扩展的 Itakura–Saito 型距离族。}\label{entropy-12-01532-t002}\\
\toprule
\cellmath{\alpha} 散度 \cellmath{D_{A}^{(\alpha)}{({\mathbf{P}||\mathbf{Q}})}} & \cellmath{\beta} 散度 \cellmath{D_{B}^{(\beta)}{({\mathbf{P}||\mathbf{Q}})}} \\ \midrule
\cellmath{\left\{ \begin{array}{l}
{\frac{1}{\alpha(\alpha - 1)}\int\left( {p^{\alpha}\, q^{1 - \alpha} - \alpha p + {(\alpha - 1)}\, q} \right)d\mu{(\mathbf{x})},} \\
{\int\left( {q\log\left( \frac{q}{p} \right) + p - q} \right)d\mu{(\mathbf{x})},\alpha = 0} \\
{\int\left( {p\log\left( \frac{p}{q} \right) - p + q} \right)d\mu{(\mathbf{x})},\alpha = 1}
\end{array} \right.} & \cellmath{\left\{ \begin{array}{l}
{\frac{1}{\beta(\beta - 1)}\int\left( {p^{\beta} + {(\beta - 1)}\, q^{\beta} - \beta\, p\, q^{\beta - 1}} \right)d\mu{(\mathbf{x})},} \\
{\int\left( {\log\left( \frac{q}{p} \right) + \frac{p}{q} - 1} \right)d\mu{(\mathbf{x})},\beta = 0} \\
{\int\left( {p\log\left( \frac{p}{q} \right) - p + q} \right)d\mu{(\mathbf{x})},\beta = 1}
\end{array} \right.} \\ \midrule
\cellmath{\left\{ \begin{array}{l}
{\frac{\int\left( {{(\frac{p + q}{2})}^{\alpha}q^{1 - \alpha} - \frac{\alpha}{2}p + {(\frac{\alpha}{2} - 1)}\, q} \right)d\mu{(\mathbf{x})}}{\alpha(\alpha - 1)},} \\
{\int\left( {q\log\left( \frac{2q}{p + q} \right) + \frac{p - q}{2}} \right)d\mu{(\mathbf{x})},\alpha = 0} \\
{\int\left( {\frac{p + q}{2}\log\left( \frac{p + q}{2q} \right) - \frac{p - q}{2}} \right)d\mu{(\mathbf{x})},\alpha = 1}
\end{array} \right.} & \cellmath{\left\{ \begin{array}{l}
{\frac{\int\left( {p^{\beta} + {(2\beta - 1)}q^{\beta} - 2\beta q^{\beta - 1}\left( \frac{p^{\beta} + q^{\beta}}{2} \right)^{\frac{1}{\beta}}} \right)d\mu{(\mathbf{x})}}{\beta(\beta - 1)},} \\
{\int\left( {\log{(\frac{q}{p})} + 2{(\sqrt{\frac{p}{q}} - 1)}} \right)\, d\mu{(\mathbf{x})},\beta = 0} \\
{\int\left( {\frac{p + q}{2}\log\left( \frac{p_{1} + q}{2q} \right) - \frac{p - q}{2}} \right)d\mu{(\mathbf{x})},\beta = 1}
\end{array} \right.} \\ \midrule
\cellmath{\left\{ \begin{array}{l}
{\frac{1}{\alpha - 1}\int{(p - q)}\left\lbrack {\left( \frac{p + q}{2q} \right)^{\alpha - 1} - 1} \right\rbrack d\mu{(\mathbf{x})},} \\
{\int\frac{{(p - q)}^{2}}{p + q}d\mu{(\mathbf{x})},\alpha = 0} \\
{\int{(p - q)}\log\left( \frac{p + q}{2q} \right)d\mu{(\mathbf{x})},\alpha = 1}
\end{array} \right.} & \cellmath{\left\{ \begin{array}{l}
{\frac{1}{\beta(\beta - 1)}\int{(p^{\beta} - q^{\beta})}\left( {1 - \frac{2^{\frac{\beta - 1}{\beta}}q^{\beta - 1}}{\left( {p^{\beta} + q^{\beta}} \right)^{\frac{\beta - 1}{\beta}}}} \right)d\mu{(\mathbf{x})},} \\
{\int\left( {\sqrt{\frac{p}{q}} - 1} \right)\log\left( \frac{p}{q} \right)d\mu{(\mathbf{x})},\beta = 0} \\
{\int{(p - q)}\log\left( \frac{p + q}{2q} \right)d\mu{(\mathbf{x})},\beta = 1}
\end{array} \right.} \\ \midrule
\cellmath{\left\{ \begin{array}{l}
{\frac{1}{\alpha(\alpha - 1)}\int\left( {p^{\alpha}q^{1 - \alpha} + p^{1 - \alpha}q^{\alpha} - p - q} \right)d\mu{(\mathbf{x})},} \\
{\int{(p - q)}\log\left( \frac{p}{q} \right)d\mu{(\mathbf{x})},\alpha = 0,1}
\end{array} \right.} & \cellmath{\left\{ \begin{array}{l}
{\frac{1}{\beta - 1}\int\left( {p^{\beta} + q^{\beta} - pq^{\beta - 1} - p^{\beta - 1}q} \right)d\mu{(\mathbf{x})},} \\
{\int\frac{{(p - q)}^{2}}{pq}d\mu{(\mathbf{x})},\beta = 0} \\
{\int{(p - q)}\log\left( \frac{p}{q} \right)d\mu{(\mathbf{x})},\beta = 1}
\end{array} \right.} \\ \midrule
\cellmath{\left\{ \begin{array}{l}
{\frac{1}{1 - \alpha}\int\left( {\frac{p + q}{2} - \left( \frac{p^{\alpha} + q^{\alpha}}{2} \right)^{\frac{1}{\alpha}}} \right)d\mu{(\mathbf{x})},} \\
{\int\left( \frac{\sqrt{p} - \sqrt{q})^{2}}{2} \right)d\mu{(\mathbf{x})},\;\;\alpha = 0} \\
{H_{S}\left( \frac{\mathbf{P} + \mathbf{Q}}{2} \right) - \frac{H_{S}{(\mathbf{P})} + H_{S}{(\mathbf{Q})}}{2},\alpha = 1}
\end{array} \right.} & \cellmath{\left\{ \begin{array}{l}
{\frac{1}{\beta(\beta - 1)}\int\left( {\frac{p^{\beta} + q^{\beta}}{2} - \left( \frac{p + q}{2} \right)^{\beta}} \right)d\mu{(\mathbf{x})},} \\
{\int\log\left( \frac{p + q}{2\sqrt{pq}} \right)d\mu{(\mathbf{x})},\beta = 0} \\
{H_{S}\left( \frac{\mathbf{P} + \mathbf{Q}}{2} \right) - \frac{H_{S}{(\mathbf{P})} + H_{S}{(\mathbf{Q})}}{2},\beta = 1}
\end{array} \right.} \\ \midrule
\cellmath{\left\{ \begin{array}{l}
{\frac{\int\left( {{(p^{1 - \alpha} + q^{1 - \alpha})}\left( \frac{p + q}{2} \right)^{\alpha} - p - q} \right)d\mu{(\mathbf{x})}}{\alpha(\alpha - 1)},} \\
{\int\left( {p\log\left( \frac{2p}{p + q} \right) + q\log\left( \frac{2q}{p + q} \right)} \right)d\mu{(\mathbf{x})},\;\;\alpha = 0} \\
{\int{(p + q)}\log\left( \frac{p + q}{2\sqrt{pq}} \right)d\mu{(\mathbf{x})},\alpha = 1}
\end{array} \right.} & \cellmath{\left\{ \begin{array}{l}
{\frac{1}{\beta - 1}\int\left( {p^{\beta} + q^{\beta} - {(p^{\beta - 1} + q^{\beta - 1})}\left( \frac{p^{\beta} + q^{\beta}}{2} \right)^{\frac{1}{\beta}}} \right)d\mu{(\mathbf{x})},} \\
{\int\left( {\sqrt{\frac{p}{q}} + \sqrt{\frac{q}{p}} - 2} \right)d\mu{(\mathbf{x})},\beta = 0} \\
{\int{(p + q)}\log\left( \frac{p + q}{2\sqrt{pq}} \right)d\mu{(\mathbf{x})},\beta = 1}
\end{array} \right.} \\ \midrule
\bottomrule\end{longtable}\end{landscape}\clearpage

在若干特例中，得到以下结果。


算术–几何散度[\hyperlink{B38-entropy-12-01532}{38},\allowbreak \hyperlink{B39-entropy-12-01532}{39}]：

\setcounter{equation}{69}
\begin{equation}
\label{FD72-entropy-12-01532}
\fitmath{D_{BS2}^{(1)} = \lim\limits_{\beta\rightarrow 1}D_{BS2}^{(\beta)} = \int{(p + q)}\log\left( \frac{p + q}{2\sqrt{pq}} \right)d\mu{(\mathbf{x})},}
\end{equation}

以及对称化的 Itakura–Saito 距离，也称 COSH 距离[\hyperlink{B12-entropy-12-01532}{12},\allowbreak \hyperlink{B13-entropy-12-01532}{13}]：

\setcounter{equation}{70}
\begin{equation}
\label{FD73-entropy-12-01532}
\fitmath{D_{BS2}^{(0)} = \lim\limits_{\beta\rightarrow 0}D_{BS2}^{(\beta)} = \int\left( {\sqrt{\frac{p}{q}} + \sqrt{\frac{q}{p}} - 2} \right)d\mu{(\mathbf{x})} = \int\frac{{(\sqrt{p} - \sqrt{q})}^{2}}{\sqrt{pq}}d\mu{(\mathbf{x})}}
\end{equation}

\setcounter{section}{3}
\section{\texorpdfstring{从 \(\beta\) 散度和 \(\alpha\) 散度生成的 \(\gamma\) 散度族}{从 β 散度和 α 散度生成的 γ 散度族}}
\label{sec4-entropy-12-01532}
Fujisawa 和 Eguchi 最近提出了基本的非对称 \(\gamma\) 散度[\hyperlink{B35-entropy-12-01532}{35}]，作为对离群值高度稳健的相似性度量：

\setcounter{equation}{71}
\begin{equation}
\label{FD74-entropy-12-01532}
\fitmath{D_{G}^{(\gamma)}{(\mathbf{P}\,||\,\mathbf{Q})} = \frac{1}{\gamma(\gamma - 1)}\log\;\left( \frac{\left( {\int p^{\gamma}{(\mathbf{x})}d\mu{(\mathbf{x})}} \right)\left( {\int q^{\gamma}{(\mathbf{x})}d\mu{(\mathbf{x})}} \right)^{\gamma - 1}}{\left( {\int p{(\mathbf{x})}\; q^{\,\gamma - 1}{(\mathbf{x})}d\mu{(\mathbf{x})}} \right)^{\gamma}} \right)}
\end{equation}

\(\gamma\) 散度对累积模式施加非线性对数变换，其中各项 \(p,q\) 不可分离。使用 \(\gamma\) 散度的主要动机，是在离群值存在时，对某些参数进行高度稳健的估计。作者说明，即使污染十分严重，离群值引起的偏差也可以变得足够小，而且某些污染可以自然、自动地被忽略[\hyperlink{B35-entropy-12-01532}{35},\allowbreak \hyperlink{B37-entropy-12-01532}{37}]。


本文说明，既可以从 \(\alpha\) 散度，也可以从 \(\beta\) 散度，直接构造完整的 \(\gamma\) 散度族。要从 \(\alpha\) 散度或 \(\beta\) 散度得到稳健的 \(\gamma\) 散度，使用以下变换，另见\hyperref[entropy-12-01532-t003]{表 3}：

\setcounter{equation}{72}
\begin{equation}
\label{FD75-entropy-12-01532}
\fitmath{\left. c_{0}\int p^{\, c_{1}}{(\mathbf{x})}\, q^{\, c_{2}}{(\mathbf{x})}\, d\mu{(\mathbf{x})}\rightarrow\log\left( {\int p^{\, c_{1}}{(\mathbf{x})}\, q^{\, c_{2}}{(\mathbf{x})}\, d\mu{(\mathbf{x})}} \right)^{c_{0}} \right.}
\end{equation}

其中 \(c_{0},c_{1}\) 和 \(c_{2}\) 为实常数，且 \(\gamma = \alpha\)。


对基本 \(\alpha\) 散度（\hyperref[FD4-entropy-12-01532]{4}）中的所有单项式施加上述变换，得到一个新的散度，本文称为 \(\alpha\)–\(\gamma\) 散度：

\[
\fitmath{\begin{array}{ccrl}
{D_{AG}^{(\gamma)}{(\mathbf{P}||\mathbf{Q})}} & & = & {\log\left( {\int p^{\gamma}{(\mathbf{x})}q^{1 - \gamma}{(\mathbf{x})}\, d\mu{(\mathbf{x})}} \right)^{\frac{1}{\gamma(\gamma - 1)}} - \log\left( {\int p(\mathbf{x})d\mu(\mathbf{x})} \right)^{\frac{1}{\gamma - 1}} + \log\left( {\int q(\mathbf{x})d\mu(\mathbf{x})} \right)^{\frac{1}{\gamma}}} \\
 & & = & {\frac{1}{\gamma(\gamma - 1)}\log\left( \frac{\int p^{\gamma}{(\mathbf{x})}q^{1 - \gamma}{(\mathbf{x})}d\mu{(\mathbf{x})}}{\left( {\int p(\mathbf{x})d\mu(\mathbf{x})} \right)^{\gamma}\left( {\int q(\mathbf{x})d\mu(\mathbf{x})} \right)^{1 - \gamma}} \right)}
\end{array}}
\]




非对称 \(\alpha\)–\(\gamma\) 散度具有以下重要性质：

\begin{enumerate}
\item \(D_{AG}^{(\gamma)}{(\mathbf{P}\,||\,\mathbf{Q})} \geq 0\)。等号成立当且仅当存在正常数 c，使 \(\mathbf{P} = c\mathbf{Q}\)。


\item 对任意 \(\gamma\)，它都具有尺度不变性：对于任意正的缩放常数 \(c_{1},c_{2}\)，有 \(D_{AG}^{(\gamma)}{(\mathbf{P}\,||\,\mathbf{Q})} = D_{AG}^{(\gamma)}(c_{1}{\mathbf{P}\,||\,}c_{2}{\mathbf{Q})}\)。


\item \(\alpha\)–\(\gamma\) 散度等价于归一化的 \(\alpha\)–Rényi 散度（\hyperref[FD25-entropy-12-01532]{25}），即

\[
\fitmath{\begin{array}{ccrl}
{D_{AG}^{(\gamma)}{(\mathbf{P}||\mathbf{Q})}} & & = & {\frac{1}{\gamma(\gamma - 1)}\log\left( \frac{\int p^{\gamma}{(\mathbf{x})}\, q^{1 - \gamma}{(\mathbf{x})}d\mu{(\mathbf{x})}}{\left( {\int p(\mathbf{x})d\mu(\mathbf{x})} \right)^{\gamma}\left( {\int q(\mathbf{x})d\mu(\mathbf{x})} \right)^{1 - \gamma}} \right)} \\
 & & = & {\frac{1}{\gamma(\gamma - 1)}\log\left( {\int{\overline{p}}^{\gamma}{(\mathbf{x})}\,{\overline{q}}^{1 - \gamma}{(\mathbf{x})}d\mu{(\mathbf{x})}} \right)} \\
 & & = & {\frac{1}{\gamma - 1}\log\left( {\int\overline{q}{(\mathbf{x})}\,\left( \frac{\overline{p}{(\mathbf{x})}}{\overline{q}{(\mathbf{x})}} \right)^{\gamma}d\mu{(\mathbf{x})}} \right)^{\frac{1}{\gamma}} = D_{AR}^{(\gamma)}(\overline{\mathbf{P}}{||}\overline{\mathbf{Q}})}
\end{array}}
\]

其中 \(\alpha = \gamma\)，归一化密度为 \(\overline{p}{{(\mathbf{x})} = p{(\mathbf{x})}/\int(p{(\mathbf{x})}d\mu{(\mathbf{x})}}\) 和 \(\overline{q}{{(\mathbf{x})} = q{(\mathbf{x})}/\int(q{(\mathbf{x})}d\mu{(\mathbf{x})}}\)。


\item 它可以表示为广义加权平均：

\setcounter{equation}{73}
\begin{equation}
\label{FD78-entropy-12-01532}
\fitmath{D_{AG}^{(\gamma)}{(\mathbf{P}||\mathbf{Q})} = \frac{1}{\gamma - 1}\log\left( {{\overline{M}}_{\gamma}\left\{ {\overline{q};\frac{\overline{p}}{\overline{q}}} \right\}} \right)}
\end{equation}

其中加权平均定义为 \({\overline{M}}_{\gamma}\left\{ {\overline{q};\frac{\overline{p}}{\overline{q}}} \right\} = \left( {\int\overline{q}{(\mathbf{x})}\,\left( \frac{\overline{p}{(\mathbf{x})}}{\overline{q}{(\mathbf{x})}} \right)^{\gamma}d\mu{(\mathbf{x})}} \right)^{\frac{1}{\gamma}}\)。


\item 当 \(\left. \gamma\rightarrow 0 \right.\) 时，\(\alpha\)–\(\gamma\) 散度变为 Kullback–Leibler 散度：

\setcounter{equation}{74}
\begin{equation}
\label{FD79-entropy-12-01532}
\fitmath{\lim\limits_{\gamma\rightarrow 0}D_{AG}^{(\gamma)}{(\mathbf{P}\,||\,\mathbf{Q})} = D_{KL}(\overline{\mathbf{P}}{\,||\,}\overline{\mathbf{Q}}{) = \int}\overline{p}{(\mathbf{x})}\,\log\frac{\overline{p}{(\mathbf{x})}}{\overline{q}{(\mathbf{x})}}d\mu{(\mathbf{x})}}
\end{equation}




\item 对于 \(\left. \gamma\rightarrow 1 \right.\)，\(\alpha\)–\(\gamma\) 散度可以通过反向 Kullback–Leibler 散度表示：

\setcounter{equation}{75}
\begin{equation}
\label{FD80-entropy-12-01532}
\fitmath{\lim\limits_{\gamma\rightarrow 1}D_{AG}^{(\gamma)}{(\mathbf{P}\,||\,\mathbf{Q})} = D_{KL}(\overline{\mathbf{Q}}{\,||\,}\overline{\mathbf{P}}{) = \int}\overline{q}{(\mathbf{x})}\,\log\frac{\overline{q}{(\mathbf{x})}}{\overline{p}{(\mathbf{x})}}\, d\mu{(\mathbf{x})}}
\end{equation}



\end{enumerate}





类似地，可以从 \(\alpha\) 散度族生成整个 \(\alpha\)–\(\gamma\) 散度族，结果汇总于\hyperref[entropy-12-01532-t003]{表 3}。


在变换（\hyperref[FD75-entropy-12-01532]{73}）中取 \(\gamma = \beta\)，还可以生成另一族 \(\gamma\) 散度，本文称为 \(\beta\)–\(\gamma\) 散度。


具体地，对基本的非对称 \(\beta\) 散度（\hyperref[FD47-entropy-12-01532]{46}）采用非线性变换（\hyperref[FD75-entropy-12-01532]{73}），得到 \(\gamma\) 散度（\hyperref[FD74-entropy-12-01532]{72}）[\hyperlink{B35-entropy-12-01532}{35}]。本文将其记为 \(\beta\)–\(\gamma\) 散度（\(D_{G}^{(\gamma)}{(\mathbf{P}||\mathbf{Q})} = D_{BG}^{(\gamma)}{(\mathbf{P}||\mathbf{Q})}\)）：

\[
\fitmath{\begin{array}{ccrl}
{D_{BG}^{(\gamma)}{(\mathbf{P}\,||\,\mathbf{Q})}} & & = & {\frac{1}{\gamma(\gamma - 1)}\left\lbrack {\log{(\int p^{\gamma}d\mu{(\mathbf{x})})} + {(\gamma - 1)}\log{(\int q^{\gamma}d\mu{(\mathbf{x})})} - \gamma\log{(\int p\; q^{\gamma - 1}d\mu{(\mathbf{x})})}} \right\rbrack} \\
 & & = & {\log\;\left( \frac{\left( {\int p\; q^{\,\gamma - 1}d\mu{(\mathbf{x})}} \right)^{\frac{1}{1 - \gamma}}}{\left( {\int p^{\gamma}d\mu{(\mathbf{x})}} \right)^{\frac{1}{\gamma(1 - \gamma)}}\left( {\int q^{\gamma}d\mu{(\mathbf{x})}} \right)^{\frac{- 1}{\gamma}}} \right)} \\
 & & = & {\frac{1}{1 - \gamma}\log\left( {\int{\widetilde{q}}^{\;\gamma}{(\mathbf{x})}\left( {\,\frac{\widetilde{p}{(\mathbf{x})}}{\widetilde{q}{(\mathbf{x})}}} \right)\; d\mu{(\mathbf{x})}} \right)}
\end{array}}
\]

其中

\setcounter{equation}{76}
\begin{equation}
\label{FD82-entropy-12-01532}
\fitmath{\widetilde{p}{(\mathbf{x})} = \frac{p(\mathbf{x})}{\left( {\int p^{\gamma}{(\mathbf{x})}d\mu{(\mathbf{x})}} \right)^{\frac{1}{\gamma}}},\widetilde{q}{(\mathbf{x})} = \frac{q(\mathbf{x})}{\left( {\int q^{\gamma}{(\mathbf{x})}d\mu{(\mathbf{x})}} \right)^{\frac{1}{\gamma}}}}
\end{equation}




类似地，对于离散密度，\(\beta\)–\(\gamma\) 散度也可以通过广义幂平均，又称幂平均或 Hölder 平均，表示为

\setcounter{equation}{77}
\begin{equation}
\label{FD83-entropy-12-01532}
\fitmath{D_{BG}^{(\gamma)}{(\mathbf{P}\,||\,\mathbf{Q})} = - \log\;\frac{\left( {\sum\limits_{i = 1}^{n}p_{i}\; q_{i}^{\,\gamma - 1}} \right)^{\frac{1}{\gamma - 1}}}{\left( {\sum\limits_{i = 1}^{n}p_{i}^{\gamma}} \right)^{\frac{1}{\gamma(\gamma - 1)}}\left( {\sum\limits_{i = 1}^{n}q_{i}^{\gamma}} \right)^{\frac{1}{\gamma}}}}
\end{equation}




因此有

\[
\fitmath{D_{BG}^{(\gamma)}{(\mathbf{P}\,||\,\mathbf{Q})} = - \log\;\frac{\left( {\sum\limits_{i = 1}^{n}\frac{p_{i}}{q_{i}}\; q_{i}^{\,\gamma}} \right)^{\frac{1}{\gamma - 1}}}{\left\lbrack \left( {\sum\limits_{i = 1}^{n}p_{i}^{\gamma}} \right)^{\frac{1}{\gamma}} \right\rbrack^{\frac{1}{\gamma - 1}}\left( {\sum\limits_{i = 1}^{n}q_{i}^{\gamma}} \right)^{\frac{1}{\gamma}}} = - \log\left( \frac{\frac{1}{n}\sum\limits_{i = 1}^{n}\frac{p_{i}}{q_{i}}\; q_{i}^{\,\gamma}}{\left( {\frac{1}{n}\sum\limits_{i = 1}^{n}p_{i}^{\gamma}} \right)^{\frac{1}{\gamma}}\left\lbrack \left( {\frac{1}{n}\sum\limits_{i = 1}^{n}q_{i}^{\gamma}} \right)^{\frac{1}{\gamma}} \right\rbrack^{\gamma - 1}} \right)^{\frac{1}{\gamma - 1}}}
\]

最终得到

\setcounter{equation}{78}
\begin{equation}
\label{FD84-entropy-12-01532}
\fitmath{D_{BG}^{(\gamma)}{(\mathbf{P}\,||\,\mathbf{Q})} = \frac{1}{1 - \gamma}\log\,\frac{\frac{1}{n}\;\sum\limits_{i = 1}^{n}\frac{p_{i}}{q_{i}}\; q_{i}^{\,\gamma}}{M_{\gamma}{\{ p_{i}\}}\,\left\lbrack {M_{\gamma}{\{ q_{i}\}}} \right\rbrack^{\gamma - 1}}}
\end{equation}

其中 \(\gamma\) 阶广义幂平均定义为

\setcounter{equation}{79}
\begin{equation}
\label{FD85-entropy-12-01532}
\fitmath{M_{\gamma}{\{ p_{i}\}} = \left( {\frac{1}{n}\sum\limits_{i = 1}^{n}p_{i}^{\gamma}} \right)^{\frac{1}{\gamma}}}
\end{equation}

它的特例包括通常的调和平均（\(\gamma = - 1)\)）、几何平均（\(\gamma = 0\)）、算术平均（\(\gamma = 1\)）和均方根 \(\gamma = 2\)，并满足以下关系：

\setcounter{equation}{80}
\begin{equation}
\label{FD86-entropy-12-01532}
\fitmath{M_{- \infty}{\{ p_{i}\}} \leq M_{- 1}{({\{ p_{i}\}})} \leq M_{0}{\{ p_{i}\}} \leq M_{1}{\{ p_{i}\}} \leq M_{2}{\{ p_{i}\}} \leq M_{\infty}{\{ p_{i}\}},}
\end{equation}

其中 \(M_{0}{\{ p_{i}\}} = \lim_{\gamma\rightarrow 0}M_{\gamma}{\{ p_{i}\}} = \left( {\prod_{i = 1}^{n}p_{i}} \right)^{1/n}\)、\(M_{- \infty}{\{ p_{i}\}} = \min_{i}{\{ p_{i}\}}\) 且 \(M_{\infty}{\{ p_{i}\}} = \max_{i}{\{ p_{i}\}}\)。


\clearpage\begin{landscape}\small\setlength{\LTcapwidth}{\linewidth}
\setcounter{table}{2}
\begin{longtable}{@{}p{0.48\linewidth}p{0.48\linewidth}@{}}
\caption[]{\(\alpha\) 散度族及其对应的稳健 \(\alpha\)–\(\gamma\) 散度；\(\overline{p}{(\mathbf{x})} = p{(\mathbf{x})}/{(\int p{(\mathbf{x})}d\mu{(\mathbf{x})})},\;\;\overline{q}{(\mathbf{x})} = q{(\mathbf{x})}/{(\int q{(\mathbf{x})}d\mu{(\mathbf{x})})}\)。当 \(\gamma = 0,1\) 时，得到广义的稳健 KL 散度。\(\gamma\) 散度可以通过广义幂平均紧凑地表示；更直接的表示见\hyperref[entropy-12-01532-t004]{表 4}。}\label{entropy-12-01532-t003}\\
\toprule
\cellmath{\alpha} 散度 \cellmath{D_{A}^{(\alpha)}{({\mathbf{P}||\mathbf{Q}})}} & 稳健的 \cellmath{\alpha}–\cellmath{\gamma} 散度 \cellmath{D_{AG}^{(\gamma)}{({c\mathbf{P}||c\mathbf{Q}})}} \\ \midrule
\cellmath{\left\{ \begin{array}{l}
{\frac{1}{\alpha(\alpha - 1)}\int\left( {p^{\alpha}q^{1 - \alpha} - \alpha p + {(\alpha - 1)}q} \right)d\mu{(\mathbf{x})},} \\
{\int\left( {q\log\left( \frac{q}{p} \right) + p - q} \right)d\mu{(\mathbf{x})},\alpha = 0} \\
{\int\left( {p\log\left( \frac{p}{q} \right) - p + q} \right)d\mu{(\mathbf{x})},\alpha = 1}
\end{array} \right.} & \cellmath{\left\{ \begin{array}{l}
{\frac{1}{\gamma - 1}\log\left( {\int\overline{q}\left( \frac{\overline{p}}{\overline{q}} \right)^{\gamma}d\mu{(\mathbf{x})}} \right)^{\frac{1}{\gamma}},} \\
{\int\overline{q}\log\left( \frac{\overline{q}}{\overline{p}} \right)d\mu{(\mathbf{x})},\gamma = 0} \\
{\int\overline{p}\log\left( \frac{\overline{p}}{\overline{q}} \right)d\mu{(\mathbf{x})},\gamma = 1}
\end{array} \right.} \\ \midrule
\cellmath{\left\{ \begin{array}{l}
{\frac{1}{\alpha(\alpha - 1)}\int\left( {p^{1 - \alpha}q^{\alpha} + {(\alpha - 1)}p - \alpha q} \right)d\mu{(\mathbf{x})},} \\
{\int\left( {p\log\left( \frac{p}{q} \right) - p + q} \right)d\mu{(\mathbf{x})},\alpha = 0} \\
{\int\left( {q\log\left( \frac{q}{p} \right) + p - q} \right)d\mu{(\mathbf{x})},\alpha = 1}
\end{array} \right.} & \cellmath{\left\{ \begin{array}{l}
{\frac{1}{\gamma - 1}\log\left( {\int\overline{p}\left( \frac{\overline{q}}{\overline{p}} \right)^{\gamma}d\mu{(\mathbf{x})}} \right)^{\frac{1}{\gamma}},} \\
{\int\overline{p}\log\left( \frac{\overline{p}}{\overline{q}} \right)d\mu{(\mathbf{x})},\gamma = 0} \\
{\int\overline{q}\log\left( \frac{\overline{q}}{\overline{p}} \right)d\mu{(\mathbf{x})},\gamma = 1}
\end{array} \right.} \\ \midrule
\cellmath{\left\{ \begin{array}{l}
{\frac{\int\left( {\left( \frac{p + q}{2} \right)^{\alpha}q^{1 - \alpha} - \frac{\alpha}{2}p - {(1 - \frac{\alpha}{2})}\, q} \right)d\mu{(\mathbf{x})}}{\alpha(\alpha - 1)},} \\
{\int\left( {q\log\left( \frac{2q}{p + q} \right) + \frac{p - q}{2}} \right)d\mu{(\mathbf{x})},\alpha = 0} \\
{\int\left( {\frac{p + q}{2}\log\left( \frac{p + q}{2q} \right) - \frac{p - q}{2}} \right)d\mu{(\mathbf{x})},\alpha = 1}
\end{array} \right.} & \cellmath{\left\{ \begin{array}{l}
{\frac{1}{\gamma - 1}\log\left\lbrack {\left( \frac{\int qd\mu(\mathbf{x})}{\int pd\mu(\mathbf{x})} \right)^{\frac{1}{2}}\left( {\int\overline{q}\left( \frac{p + q}{2q} \right)^{\gamma}d\mu{(\mathbf{x})}} \right)^{\frac{1}{\gamma}}} \right\rbrack} \\
{\int\left( {(\overline{q}\log\left( \frac{2q}{p + q} \right)} \right)d\mu{(\mathbf{x})},\gamma = 0} \\
{\int\left( {\frac{\overline{p} + \overline{q}}{2}\log\left( \frac{p + q}{2q} \right)} \right)d\mu{(\mathbf{x})},\gamma = 1}
\end{array} \right.} \\ \midrule
\cellmath{\left\{ \begin{array}{l}
{\frac{1}{\alpha - 1}\int\left\lbrack {\left( \frac{p + q}{2q} \right)^{\alpha - 1} - 1} \right\rbrack{(p - q)}d\mu{(\mathbf{x})},} \\
{\int{(p - q)}\log\left( \frac{p + q}{2q} \right)d\mu{(\mathbf{x})},\alpha = 1}
\end{array} \right.} & \cellmath{\left\{ \begin{array}{l}
{\log\left( \frac{\int{(\overline{p}\left( \frac{p + q}{2q} \right)^{\gamma - 1})}d\mu{(\mathbf{x})}}{\int{(\overline{q}\left( \frac{p + q}{2q} \right)^{\gamma - 1})}d\mu{(\mathbf{x})}} \right)^{\frac{1}{\gamma - 1}}} \\
{\int{(\overline{p} - \overline{q})}\log\left( \frac{p + q}{2q} \right)d\mu{(\mathbf{x})},\gamma = 1}
\end{array} \right.} \\ \midrule
\cellmath{\left\{ \begin{array}{l}
{\frac{1}{\alpha(\alpha - 1)}\int\left( {p^{\alpha}q^{1 - \alpha} + p^{1 - \alpha}q^{\alpha} - p - q} \right)d\mu{(\mathbf{x})},} \\
{\int{(p - q)}\log\left( \frac{p}{q} \right)d\mu{(\mathbf{x})},\alpha = 0,1}
\end{array} \right.} & \cellmath{\left\{ \begin{array}{l}
{\frac{1}{\gamma - 1}\log\left\lbrack {\int\overline{q}\left( \frac{p}{q} \right)^{\gamma}d\mu{(\mathbf{x})}\int\overline{p}\left( \frac{q}{p} \right)^{\gamma}d\mu{(\mathbf{x})}} \right\rbrack^{\frac{1}{\gamma}}} \\
{\int{(\overline{p} - \overline{q})}\log\left( \frac{p}{q} \right)d\mu{(\mathbf{x})},\gamma = 0,1}
\end{array} \right.} \\ \midrule
\cellmath{\left\{ \begin{array}{l}
{\frac{\int\left( {{(p^{1 - \alpha} + q^{1 - \alpha})}\left( \frac{p + q}{2} \right)^{\alpha} - p - q} \right)d\mu{(\mathbf{x})}}{\alpha(\alpha - 1)},} \\
{\int\left( {p\log\left( \frac{2p}{p + q} \right) + q\log\left( \frac{2q}{p + q} \right)} \right)d\mu{(\mathbf{x})},\alpha = 0} \\
{\int{(p + q)}\log\left( \frac{p + q}{2\sqrt{pq}} \right)d\mu{(\mathbf{x})},\alpha = 1}
\end{array} \right.} & \cellmath{\left\{ \begin{array}{l}
{\frac{1}{\gamma - 1}{\log\left\lbrack {\int\overline{p}\left( \frac{p + q}{2p} \right)^{\gamma}d\mu{(\mathbf{x})}\,\int\overline{q}\left( \frac{p + q}{2q} \right)^{\gamma}d\mu{(\mathbf{x})}} \right\rbrack^{\frac{1}{\gamma}}}} \\
{\int\left( {\overline{p}\log\left( \frac{2p}{p + q} \right) + \overline{q}\log\left( \frac{2q}{p + q} \right)} \right)d\mu{(\mathbf{x})},\gamma = 0} \\
{\int{(\overline{p} + \overline{q})}\log\left( \frac{p + q}{2\sqrt{pq}} \right)d\mu{(\mathbf{x})},\gamma = 1.}
\end{array} \right.} \\ \midrule
\bottomrule\end{longtable}\end{landscape}\clearpage
\clearpage\begin{landscape}\small\setlength{\LTcapwidth}{\linewidth}
\setcounter{table}{3}
\begin{longtable}{@{}p{0.48\linewidth}p{0.48\linewidth}@{}}
\caption[]{基本的 \(\alpha\) 散度、\(\beta\) 散度，以及由它们直接生成的对应 \(\gamma\) 散度。\(\gamma\) 散度的幂平均表示另见\hyperref[entropy-12-01532-t003]{表 3}。}\label{entropy-12-01532-t004}\\
\toprule
散度 \cellmath{D_{A}^{(\alpha)}{({\mathbf{P}||\mathbf{Q}})}} 或 \cellmath{D_{B}^{(\beta)}{({\mathbf{P}||\mathbf{Q}})}} & \cellmath{\gamma} 散度 \cellmath{D_{AG}^{(\gamma)}{({c\mathbf{P}||c\mathbf{Q}})}} 和 \cellmath{D_{BG}^{(\gamma)}{({c\mathbf{P}||c\mathbf{Q}})}} \\ \midrule
\cellmath{\frac{1}{\alpha(1 - \alpha)}\int\left( {\alpha\, p + {(1 - \alpha)}\, q - p^{\alpha}\, q^{1 - \alpha}} \right)d\mu{(\mathbf{x})}} & \cellmath{ - \log\frac{\left( {\int p^{\gamma}\, q^{1 - \gamma}d\mu{(\mathbf{x})}} \right)^{\; 1/(\gamma(1 - \gamma))}}{\left( {\int pd\mu(\mathbf{x})} \right)^{1/(1 - \gamma)}\;\left( {\int qd\mu(\mathbf{x})} \right)^{1/\gamma}}} \\ \midrule
\cellmath{\frac{1}{\beta(\beta - 1)}\int\left( {p^{\beta} + {(\beta - 1)}\, q^{\beta} - \beta\, p\, q^{\beta - 1}} \right)d\mu{(\mathbf{x})}} & \cellmath{ - \log\frac{\left( {\int p\, q^{\gamma - 1}d\mu{(\mathbf{x})}} \right)^{\; 1/(\gamma - 1)}}{\left( {\int p^{\gamma}d\mu{(\mathbf{x})}} \right)^{1/(\gamma(\gamma - 1))}\;\left( {\int q^{\gamma}d\mu{(\mathbf{x})}} \right)^{1/\gamma}}} \\ \midrule
\cellmath{\frac{1}{\alpha(1 - \alpha)}\int\left( {p + q - p^{\alpha}\, q^{1 - \alpha} - p^{1 - \alpha}\, q^{\alpha}} \right)d\mu{(\mathbf{x})}} & \cellmath{\frac{- 1}{\gamma(1 - \gamma)}\log\frac{\left( {\int p^{\gamma}\, q^{1 - \gamma}d\mu{(\mathbf{x})}} \right)\,\left( {\int p^{1 - \gamma}\, q^{\gamma}d\mu{(\mathbf{x})}} \right)}{\int pd\mu(\mathbf{x})\;\int qd\mu(\mathbf{x})}} \\ \midrule
\cellmath{\frac{1}{\beta - 1}\int\left( {p^{\beta} + q^{\beta} - p\, q^{\beta - 1} - p^{\beta - 1}\, q} \right)d\mu{(\mathbf{x})}} & \cellmath{\frac{- 1}{\gamma - 1}\log\frac{\left( {\int p\, q^{\gamma - 1}d\mu{(\mathbf{x})}} \right)\left( {\,\int p^{\gamma - 1}\, qd\mu{(\mathbf{x})}} \right)}{\left( {\int p^{\gamma}d\mu{(\mathbf{x})}} \right)\left( {\int q^{\gamma}d\mu{(\mathbf{x})}} \right)}} \\ \midrule
\cellmath{\frac{\int\left( {\frac{\alpha}{2}\, p + {(1 - \frac{\alpha}{2})}\, q - \left( \frac{p + q}{2} \right)^{\alpha}q^{\alpha - 1}} \right)d\mu{(\mathbf{x})}}{\alpha(1 - \alpha)}} & \cellmath{\frac{- 1}{\gamma(1 - \gamma)}\log\frac{\int\left( \frac{p + q}{2} \right)^{\gamma - 1}q^{\gamma - 1}d\mu{(\mathbf{x})}}{\left( {\int pd\mu(\mathbf{x})} \right)^{\gamma/2}\;\left( {\int qd\mu(\mathbf{x})} \right)^{1 - \gamma/2}}} \\ \midrule
\cellmath{\frac{1}{1 - \alpha}\int{(p - q)}\left( {1 - \frac{p + q}{2q}} \right)^{\alpha - 1}d\mu{(\mathbf{x})}} & \cellmath{\frac{1}{1 - \gamma}\log\frac{\int pd\mu{(\mathbf{x})}\,\int q\left( \frac{p + q}{2q} \right)^{\gamma - 1}d\mu{(\mathbf{x})}}{\int qd\mu{(\mathbf{x})}\,\int p\left( \frac{p + q}{2q} \right)^{\gamma - 1}d\mu{(\mathbf{x})}}} \\ \midrule
\cellmath{\frac{\int\left( {p + q - {(p^{1 - \alpha} + q^{1 - \alpha})}\left( \frac{p + q}{2} \right)^{\alpha}} \right)d\mu{(\mathbf{x})}}{\alpha(1 - \alpha)}} & \cellmath{\;\log\left( \frac{\int p^{1 - \gamma}\left( \frac{p + q}{2} \right)^{\,\gamma}d\mu{(\mathbf{x})}\int q^{1 - \gamma}\left( \frac{p + q}{2} \right)^{\,\gamma}d\mu{(\mathbf{x})}}{\int pd\mu(\mathbf{x})\;\int qd\mu(\mathbf{x})} \right)^{\frac{1}{\gamma(\gamma - 1)}}} \\ \midrule
\bottomrule\end{longtable}\end{landscape}\clearpage
非对称 \(\beta\)–\(\gamma\) 散度具有以下性质[\hyperlink{B30-entropy-12-01532}{30},\allowbreak \hyperlink{B35-entropy-12-01532}{35}]：

\begin{enumerate}
\item \(D_{BG}^{(\gamma)}{(\mathbf{P}\,||\,\mathbf{Q})} \geq 0\)。等号成立当且仅当存在正常数 c，使 \(\mathbf{P} = c\mathbf{Q}\)。


\item 它具有尺度不变性，即对于任意正的缩放常数 \(c_{1},c_{2}\)，有 \(D_{BG}^{(\gamma)}{(\mathbf{P}\,||\,\mathbf{Q})} = D_{BG}^{(\gamma)}(c_{1}{\mathbf{P}\,||\,}c_{2}{\mathbf{Q})}\)。


\item 当 \(\left. \gamma\rightarrow 1 \right.\) 时，\(\gamma\) 散度变为 Kullback–Leibler 散度：

\setcounter{equation}{81}
\begin{equation}
\label{FD87-entropy-12-01532}
\fitmath{\lim\limits_{\gamma\rightarrow 1}D_{BG}^{(\gamma)}{(\mathbf{P}\,||\,\mathbf{Q})} = D_{KL}(\overline{\mathbf{P}}{\,||\,}\overline{\mathbf{Q}}{) = \int}\overline{p}\,\log\frac{\overline{p}}{\overline{q}}d\mu{(\mathbf{x})}}
\end{equation}

其中 \(\overline{p} = p/\int pd\mu{(\mathbf{x})}\) 且 \(\overline{q} = q/\int qd\mu{(\mathbf{x})}\)。


\item 当 \(\left. \gamma\rightarrow 0 \right.\) 时，\(\gamma\) 散度可以表示为

\setcounter{equation}{82}
\begin{equation}
\label{FD88-entropy-12-01532}
\fitmath{\lim\limits_{\gamma\rightarrow 0}D_{BG}^{(\gamma)}{(\mathbf{P}\,||\,\mathbf{Q})} = \int\log\frac{q(\mathbf{x})}{p(\mathbf{x})}\, d\mu{(\mathbf{x})} + \log\left( {\int\frac{p(\mathbf{x})}{q(\mathbf{x})}\, d\mu{(\mathbf{x})}} \right)}
\end{equation}

对于离散 \(\gamma\) 散度，相应公式为

\setcounter{equation}{83}
\begin{equation}
\label{FD89-entropy-12-01532}
\fitmath{D_{BG}^{(0)}{(\mathbf{P}\,||\,\mathbf{Q})} = \frac{1}{n}\sum\limits_{i = 1}^{n}\left( {\log\frac{q_{i}}{p_{i}}} \right) + \log\left( {\sum\limits_{i = 1}^{n}\frac{p_{i}}{q_{i}}} \right) - \log{(n)} = \log\frac{\frac{1}{n}\sum\limits_{i = 1}^{n}\frac{p_{i}}{q_{i}}}{\left( {\prod\limits_{i = 1}^{n}\frac{p_{i}}{q_{i}}} \right)^{1/n}}}
\end{equation}



\end{enumerate}





与 \(\alpha\) 散度和 \(\beta\) 散度类似，也可以定义对称 \(\beta\)–\(\gamma\) 散度：

\[
\fitmath{D_{BGS}^{(\gamma)}{(\mathbf{P}\,||\,\mathbf{Q})} = D_{BG}^{(\gamma)}{(\mathbf{P}\,||\,\mathbf{Q})} + D_{BG}^{(\gamma)}{(\mathbf{Q}\,||\,\mathbf{P})} = \frac{1}{1 - \gamma}\;\log\left\lbrack \frac{\left( {\int p^{\gamma - 1}\; q\, d\mu{(\mathbf{x})}} \right)\left( {\int p\; q^{\gamma - 1}\, d\mu{(\mathbf{x})}} \right)}{\left( {\int p^{\gamma}d\mu{(\mathbf{x})}} \right)\left( {\int q^{\gamma}d\mu{(\mathbf{x})}} \right)} \right\rbrack}
\]




对称 \(\gamma\) 散度具有与非对称 \(\gamma\) 散度类似的性质：

\begin{enumerate}
\item \(D_{BGS}^{(\gamma)}{(\mathbf{P}\,||\,\mathbf{Q})} \geq 0\)。等号成立当且仅当存在正常数 c，使 \(\mathbf{P} = c\mathbf{Q}\)；特别地，有 \(p = q,\;\;\forall i\)。


\item 它具有尺度不变性，即

\setcounter{equation}{84}
\begin{equation}
\label{FD90-entropy-12-01532}
\fitmath{D_{BGS}^{(\gamma)}{(\mathbf{P}\,||\,\mathbf{Q})} = D_{BGS}^{(\gamma)}(c_{1}{\mathbf{P}\,||\,}c_{2}{\mathbf{Q})}}
\end{equation}

其中 \(c_{1},c_{2}\) 为任意正的缩放常数。


\item 当 \(\left. \gamma\rightarrow 1 \right.\) 时，它退化为对称 Kullback–Leibler 散度的一种特殊形式，也称 J 散度：

\setcounter{equation}{85}
\begin{equation}
\label{FD91-entropy-12-01532}
\fitmath{\lim\limits_{\gamma\rightarrow 1}D_{BGS}^{(\gamma)}{(\mathbf{P}\,||\,\mathbf{Q})} = \int{(\overline{p} - \overline{q})}\log\frac{\overline{p}}{\overline{q}}\, d\mu{(\mathbf{x})}}
\end{equation}

其中 \(\overline{p} = p/\int pd\mu{(\mathbf{x})}\) 且 \(\overline{q} = q/\int qd\mu{(\mathbf{x})}\)。


\item 当 \(\gamma = 0\) 时，得到一个可由加权算术平均表示的简单散度：

\setcounter{equation}{86}
\begin{equation}
\label{FD92-entropy-12-01532}
\fitmath{D_{BGS}^{(0)}{(\mathbf{P}\,||\,\mathbf{Q})} = \log\left( {\int w\,\frac{p(\mathbf{x})}{q(\mathbf{x})}\, d\mu{(\mathbf{x})}\;\int w\,\frac{q(\mathbf{x})}{p(\mathbf{x})}\, d\mu{(\mathbf{x})}} \right)}
\end{equation}

其中权重函数 \(w > 0\) 满足 \(\int wd\mu(\mathbf{x}) = 1\)。

对于离散 \(\beta\)–\(\gamma\) 散度，简称 \(\gamma\) 散度，得到

\[
\fitmath{\begin{array}{clr}
{D_{BGS}^{(0)}{(\mathbf{P}\,||\,\mathbf{Q})}} & = & {\log\left( {\sum\limits_{i = 1}^{n}\frac{p_{i}}{q_{i}}\sum\limits_{i = 1}^{n}\frac{q_{i}}{p_{i}}} \right) - \log{(n)}^{2}} \\
 & = & {\log\left( {\left( {\frac{1}{n}\sum\limits_{i = 1}^{n}\frac{p_{i}}{q_{i}}} \right)\left( {\frac{1}{n}\sum\limits_{i = 1}^{n}\frac{q_{i}}{p_{i}}} \right)} \right)} \\
 & = & {\log\left( {M_{1}\left\{ \frac{p_{i}}{q_{i}} \right\}\, M_{1}\left\{ \frac{q_{i}}{p_{i}} \right\}} \right)}
\end{array}}
\]

当 \(\left. n\rightarrow\infty \right.\) 时，离散的对称 \(\gamma\) 散度可用期望表示为 \(D_{BGS}^{(1)}{(\mathbf{P}\,||\,\mathbf{Q})} = \log{(E{\{\mathbf{u}\}}\; E{\{\mathbf{u}^{- 1}\}})}\)，其中 \(u_{i} = {\{ p_{i}/q_{i}\}}\) 且 \(u_{i}^{- 1} = {\{ q_{i}/p_{i}\}}\)。


\item 当 \(\gamma = 2\) 时，非对称 \(\gamma\) 散度与对称 \(\gamma\) 散度相等，退化为 Principe 提出的 Cauchy–Schwarz 散度[\hyperlink{B83-entropy-12-01532}{83}]：

\[
\fitmath{\begin{array}{ccrl}
{D_{BG}^{(2)}{(\mathbf{P}\,||\,\mathbf{Q})}} & & = & {D_{BG}^{(2)}{(\mathbf{Q}\,||\,\mathbf{P})} = \frac{1}{2}D_{BGS}^{(2)}{(\mathbf{P}\,||\,\mathbf{Q})} = \frac{1}{2}D_{CS}{(\mathbf{P}\,||\,\mathbf{Q})}} \\
 & & = & {- \log\frac{\int p(\mathbf{x})\; q(\mathbf{x})\; d\mu(\mathbf{x})}{\left( {\int p^{2}{(\mathbf{x})}d\mu{(\mathbf{x})}} \right)^{1/2}\left( {\int q^{2}{(\mathbf{x})}d\mu{(\mathbf{x})}} \right)^{1/2}}}
\end{array}}
\]



\end{enumerate}





由 \(\beta\) 散度导出的基本非对称 \(\beta\)–\(\gamma\) 散度，与文献[\hyperlink{B35-entropy-12-01532}{35}]定义的 \(\gamma\) 散度完全等价；由 \(\alpha\) 散度族导出的 \(\alpha\)–\(\gamma\) 散度具有不同的表达式，但性质相似。

\setcounter{section}{4}
\section{\texorpdfstring{非对称散度之间的关系及其统一表示}{非对称散度之间的关系及其统一表示}}
\label{sec5-entropy-12-01532}
本文讨论的广义散度之间的基本联系与类比汇总于\hyperref[entropy-12-01532-t005]{表 5}。


\(\alpha\)、\(\beta\) 和 \(\gamma\) 散度的基本区别在于缩放性质：

\setcounter{equation}{87}
\begin{equation}
\label{FD95-entropy-12-01532}
\fitmath{D_{A}^{(\alpha)}{(c\,\mathbf{P}\,||\, c\,\mathbf{Q})} = c\, D_{A}^{(\alpha)}{(\mathbf{P}\,||\,\mathbf{Q})}}
\end{equation}



\setcounter{equation}{88}
\begin{equation}
\label{FD96-entropy-12-01532}
\fitmath{D_{B}^{(\beta)}{(c\,\mathbf{P}\,||\, c\,\mathbf{Q})} = c^{\;\beta}\, D_{B}^{(\beta)}{(\mathbf{P}\,||\,\mathbf{Q})}}
\end{equation}



\clearpage\begin{landscape}\small\setlength{\LTcapwidth}{\linewidth}
\setcounter{table}{4}
\begin{longtable}{@{}p{0.48\linewidth}p{0.48\linewidth}@{}}
\caption[]{基本非对称散度及其相互关系。通过变换 \(\left. p\rightarrow p^{\beta}, q\rightarrow q^{\beta}\text{且}\alpha = \beta^{- 1} \right.\)，可由 \(\alpha\) 散度得到 \(\beta\) 散度；通过对偶变换 \(\left. p\rightarrow p^{\alpha}, q\rightarrow q^{\alpha}\text{且}\beta = \alpha^{- 1} \right.\)，可由 \(\beta\) 散度得到 \(\alpha\) 散度。通过变换 \(\left. c_{0}\int p^{\, c_{1}}{(\mathbf{x})}\, q^{\, c_{2}}{(\mathbf{x})}\, d\mu{(\mathbf{x})}\rightarrow c_{0}\,\log\int p^{\, c_{1}}{(\mathbf{x})}\, q^{\, c_{2}}{(\mathbf{x})}\, d\mu{(\mathbf{x})} \right.\)，还可以由 \(\alpha\) 散度和 \(\beta\) 散度生成 \(\gamma\) 散度。此外，通过变换 \(\left. p\rightarrow p^{\gamma_{new}}, q\rightarrow q^{\gamma_{new}}\text{且}\gamma = \gamma_{new}^{- 1} \right.\)，可由 \(\alpha\)–\(\gamma\) 散度生成 \(\beta\)–\(\gamma\) 散度。}\label{entropy-12-01532-t005}\\
\toprule
散度名称 & 公式 \\ \midrule
\cellmath{\alpha} 散度 & \cellmath{D_{A}^{(\alpha)}{({\mathbf{P}||\mathbf{Q}})} = \frac{\int(p^{\alpha}\; q^{1 - \alpha} - \alpha\; p + {(\alpha - 1)}\; q)d\mu{(\mathbf{x})}}{\alpha(\alpha - 1)}} \\ \midrule
\cellmath{\beta} 散度 & \cellmath{D_{B}^{(\beta)}{({\mathbf{P}||\mathbf{Q}})} = \frac{\int(p^{\beta} + {(\beta - 1)}\; q^{\beta} - \beta\; p\; q^{\beta - 1})d\mu{(\mathbf{x})}}{\beta(\beta - 1)}} \\ \midrule
\cellmath{\gamma} 散度 & \cellmath{D_{AG}^{(\gamma)}{({\mathbf{P}||\mathbf{Q}})} = \frac{\log\;{(\int p^{\gamma}\; q^{1 - \gamma}d\mu{(\mathbf{x})})} - \gamma\log\;{(\int pd\mu{(\mathbf{x})})} + {(\gamma - 1)}\log\;{(\int qd\mu{(\mathbf{x})})}}{\gamma(\gamma - 1)}} \\ \midrule
 & \cellmath{ = \frac{1}{\gamma - 1}\log\left( {\int\overline{q}\,\left( \frac{\overline{p}}{\overline{q}} \right)^{\gamma}d\mu{(\mathbf{x})}} \right)^{\frac{1}{\gamma}}} \\ \midrule
 & \cellmath{D_{BG}^{(\gamma)}{({\mathbf{P}||\mathbf{Q}})} = \frac{\log\;{(\int p^{\gamma}d\mu{(\mathbf{x})})} + {(\gamma - 1)}\log\;{(\int q^{\gamma}d\mu{(\mathbf{x})})} - \gamma\log\;{(\int p\; q^{\gamma - 1}d\mu{(\mathbf{x})})}}{\gamma(\gamma - 1)}} \\ \midrule
 & \cellmath{ = \frac{1}{1 - \gamma}\log\left( {\int{\widetilde{q}}^{\;\gamma}\,\left( \frac{\widetilde{p}}{\widetilde{q}} \right)^{}d\mu{(\mathbf{x})}} \right)} )) \\ \midrule
\cellmath{\alpha}–Rényi 散度 & \cellmath{D_{AR}^{(\alpha)}{({\mathbf{P}||\mathbf{Q}})} = \frac{\log(\int{(p^{\,\alpha}q^{\, 1 - \alpha} - \alpha p + {(\alpha - 1)}\, q)}d\mu{(\mathbf{x})} + 1)}{\alpha(\alpha - 1)}} \\ \midrule
Bregman 散度 & \cellmath{D_{\Phi}{({\mathbf{P}||\mathbf{Q}})} = \int\left( {\Phi{(p)} - \Phi{(q)} - \frac{\delta\Phi}{\delta q}{(p - q)}} \right)d\mu{(\mathbf{x})}}，见式（\hyperref[FD58-entropy-12-01532]{57}）。 \\ \midrule
Csiszár–Morimoto 散度 & \cellmath{D_{f}{({\mathbf{P}||\mathbf{Q}})} = \int qf\left( \frac{p}{q} \right)d\mu{(\mathbf{x})}}，见式（\hyperref[FD58-entropy-12-01532]{57}），其中 \cellmath{f{(u)} = \Phi{(t)}_{|t = u}}。 \\ \midrule
\bottomrule\end{longtable}\end{landscape}\clearpage



\setcounter{equation}{89}
\begin{equation}
\label{FD97-entropy-12-01532}
\fitmath{D_{G}^{(\gamma)}{(c\,\mathbf{P}\,||\, c\,\mathbf{Q})} = D_{G}^{(\gamma)}{(\mathbf{P}\,||\,\mathbf{Q})}}
\end{equation}

其中 \(c > 0\) 可任意取值。


所有 \(\alpha\) 散度的值都与缩放因子成正比，而且这种缩放关系与参数 \(\alpha\) 无关；所有 \(\beta\) 散度的缩放因子，则以指数形式依赖参数 \(\beta\)。\(\beta\) 散度只有在 \(\beta = 0\)，即 Itakura–Saito 距离这一特例下，才对缩放因子保持不变；\(\gamma\) 散度则对任意参数 \(\gamma\) 都具有尺度不变性。


根据前面各节的分析与讨论，广泛的非对称散度可以用下列广义函数描述：

\setcounter{equation}{90}
\begin{equation}
\label{FD98-entropy-12-01532}
\fitmath{D_{AC}^{(\alpha,\beta,r)}{(\mathbf{P}||\mathbf{Q})} = \frac{1}{\alpha(\beta - 1)\lambda\, r}\left\lbrack {\left( {{\widetilde{D}}_{AC}^{(\alpha,\beta)} + 1} \right)^{r} - 1} \right\rbrack}
\end{equation}

其中 \(r > 0\)，并且

\setcounter{equation}{91}
\begin{equation}
\label{FD99-entropy-12-01532}
\fitmath{\begin{array}{rcl}
{{\widetilde{D}}_{AC}^{(\alpha,\beta)}{(\mathbf{P} \parallel \mathbf{Q})}} & = & {\int\left( {\alpha p^{\lambda} + {(\beta - 1)}q^{\lambda} - \lambda p^{\alpha}q^{\beta - 1}} \right)d\mu{(\mathbf{x})}}
\end{array}}
\end{equation}

这里 \(\alpha \neq 0,\;\;\beta \neq 1\) 且 \(\lambda = \alpha + \beta - 1 \neq 0\)。


这样定义的函数与广义对数（\hyperref[FD3-entropy-12-01532]{3}）具有类似结构。只有在参数满足 \({\widetilde{D}}_{AC}^{(\alpha,\beta)}{(\mathbf{P} \parallel \mathbf{Q})} \geq 0\)，且函数仅在 \(\mathbf{P} = \mathbf{Q}\) 时为零的条件下，上述函数才构成散度。


在若干特例中，得到以下散度：

\begin{itemize}
\item 当 \(r = \lambda = 1\) 且 \(\alpha + \beta = 2\) 时，得到 \(\alpha\) 散度（\hyperref[FD4-entropy-12-01532]{4}）。


\item 当 \(r = 1\) 且 \(\alpha = 1\)，即 \(\;(\beta = \lambda)\) 时，得到 \(\beta\) 散度（\hyperref[FD47-entropy-12-01532]{46}）。


\item 当 \(\left. r\rightarrow 0 \right.\) 且 \(\alpha + \beta = 2\)，即 \(\;(\lambda = 1)\) 时，得到 \(\alpha\)–Rényi 散度（\hyperref[FD25-entropy-12-01532]{25}）；对于归一化概率密度 \(\overline{p} = p/{(\int pd\mu{(\mathbf{x})})}\) 和 \(\overline{q} = q/{(\int qd\mu{(\mathbf{x})})}\)，它退化为 \(\alpha\)–\(\gamma\) 散度。


\item 当 \(\left. r\rightarrow 0 \right.\) 且 \(\alpha = 1\)，即 \(\;(\beta = \lambda = \gamma\)，并采用归一化密度 \(\widetilde{p} = p/{(\int p^{\gamma}d\mu{(\mathbf{x})})}^{1/\gamma}\) 和 \(\widetilde{q} = q/{(\int q^{\,\gamma}d\mu{(\mathbf{x})})}^{1/\gamma}\) 时，得到 Fujisawa 和 Eguchi 最近提出的 \(\gamma\) 散度（\hyperref[FD74-entropy-12-01532]{72}）[\hyperlink{B35-entropy-12-01532}{35}]，本文称其为 \(\beta\)–\(\gamma\) 散度。

\end{itemize}





函数（\hyperref[FD98-entropy-12-01532]{91}）还具有许多一般性质，这些性质超出了本文的讨论范围。

\subsection*{\texorpdfstring{对偶性}{对偶性}}
\(\alpha\) 散度和 \(\alpha\)–\(\gamma\) 散度具有以下对偶关系：

\setcounter{equation}{92}
\begin{equation}
\label{FD100-entropy-12-01532}
\fitmath{D_{A}^{(\alpha)}{(\mathbf{P}\,||\,\mathbf{Q})} = c^{- 1}D_{A}^{(1 - \alpha)}{(c\,\mathbf{Q}\,||\, c\,\mathbf{P})},\quad D_{AG}^{(\gamma)}{(\mathbf{P}\,||\,\mathbf{Q})} = D_{AG}^{(1 - \gamma)}(c_{1}{\mathbf{Q}\,||\,}c_{2}{\mathbf{P})}}
\end{equation}

要为 \(\beta\) 散度建立对偶性，需要先将其写成稍一般的双参数散度形式：

\setcounter{equation}{93}
\begin{equation}
\label{FD101-entropy-12-01532}
\fitmath{D_{AB}^{(\beta_{1},\beta_{2})}{(\mathbf{P}\,||\,\mathbf{Q})} = \frac{1}{\beta_{1}\,\beta_{2}{(\beta_{1} + \beta_{2})}}\int\left( {\beta_{1}\, p^{\beta_{1} + \beta_{2}} + \beta_{2}\, q^{\beta_{1} + \beta_{2}} - {(\beta_{1} + \beta_{2})}p^{\beta_{1}}q^{\beta_{2}}} \right)d\mu{(\mathbf{x})}}
\end{equation}

此时，\(\beta\) 散度的对偶性可表示为

\setcounter{equation}{94}
\begin{equation}
\label{FD102-entropy-12-01532}
\fitmath{D_{B}^{(\beta)}{(\mathbf{P}\,||\,\mathbf{Q})} = D_{AB}^{(1,\beta - 1)}{(\mathbf{P}\,||\,\mathbf{Q})} = c^{- \beta}D_{AB}^{(\beta - 1,1)}{(c\,\mathbf{Q}\,||\, c\,\mathbf{P})}}
\end{equation}

或更一般地表示为

\setcounter{equation}{95}
\begin{equation}
\label{FD103-entropy-12-01532}
\fitmath{D_{AB}^{(\beta_{1},\beta_{2})}{(\mathbf{P}\,||\,\mathbf{Q})} = c^{- (\beta_{1} + \beta_{2})}D_{AB}^{(\beta_{2},\beta_{1})}{(c\,\mathbf{Q}\,||\, c\,\mathbf{P})}}
\end{equation}

类似地，将非对称 \(\gamma\) 散度族定义为更一般的双参数形式：

\setcounter{equation}{96}
\begin{equation}
\label{FD104-entropy-12-01532}
\fitmath{\begin{array}{r}
{D_{ABG}^{(\gamma_{1},\gamma_{2})}{(\mathbf{P} \parallel \mathbf{Q})} = \log\frac{\left( {\int p^{\gamma_{1} + \gamma_{2}}d\mu{(\mathbf{x})}} \right)^{\frac{1}{\gamma_{2}{(\gamma_{1} + \gamma_{2})}}}\left( {\int q^{\gamma_{1} + \gamma_{2}}d\mu{(\mathbf{x})}} \right)^{\frac{1}{\gamma_{1}{(\gamma_{1} + \gamma_{2})}}}}{\left( {\int p^{\gamma_{1}}q^{\gamma_{2}}d\mu{(\mathbf{x})}} \right)^{\frac{1}{\gamma_{1}\gamma_{2}}}}}
\end{array}}
\end{equation}




于是，\(\beta\)–\(\gamma\) 散度，即基本 \(\gamma\) 散度，具有以下对偶性质：

\setcounter{equation}{97}
\begin{equation}
\label{FD105-entropy-12-01532}
\fitmath{D_{BG}^{(\gamma)}{(\mathbf{P}\,||\,\mathbf{Q})} = D_{ABG}^{(1,\gamma - 1)}{(\mathbf{P}\,||\,\mathbf{Q})} = D_{ABG}^{(\gamma - 1,1)}(c_{1}{\,\mathbf{Q}\,||\,}c_{2}{\,\mathbf{P})}}
\end{equation}

或更一般地表示为

\setcounter{equation}{98}
\begin{equation}
\label{FD106-entropy-12-01532}
\fitmath{D_{ABG}^{(\gamma_{1},\gamma_{2})}{(\mathbf{P}\,||\,\mathbf{Q})} = D_{ABG}^{(\gamma_{2},\gamma_{1})}(c_{1}{\mathbf{Q}\,||\,}c_{2}{\mathbf{P})}}
\end{equation}




\setcounter{section}{5}
\setcounter{subsection}{0}
\subsection{\texorpdfstring{结论与讨论}{结论与讨论}}
本文的主要目标是建立并扩展几类散度之间的联系。主要结果汇总于\hyperref[entropy-12-01532-t002]{表 2}、\hyperref[entropy-12-01532-t003]{表 3}、\hyperref[entropy-12-01532-t004]{表 4}和\hyperref[entropy-12-01532-t005]{表 5}。我们扩展并统一了几个重要的散度类，给出了它们之间的联系与对应关系。本文提出的参数化广义散度族，尤其是 \(\alpha\) 散度、\(\alpha\)–Rényi 散度和 \(\beta\) 散度，使我们能够在同一框架下考察各种常用基本散度之间的平滑连接。例如，\(\beta\) 散度族通过广义 KL 的 I 散度，平滑地连接平方欧氏距离，即 \(L_{2}\) 范数距离，与广义扩展的 Itakura–Saito 型距离。我们还说明，在 \(\beta = 0\) 时，可以从 \(\beta\) 散度生成整个扩展的稳健 Itakura–Saito 型距离族；它们具有尺度不变性，在语音处理等领域具有潜在应用。


我们还说明，可以分别从 \(\beta\) 散度和 \(\alpha\) 散度生成两个不同的 \(\gamma\) 散度族。\(\gamma\) 散度具有与参数 \(\gamma\) 的取值无关的全局尺度不变性，以及高度稳健性。本文回顾并扩展了 \(\gamma\) 散度的基本性质；所用方法也能够生成广泛的新散度，尤其是 \(\gamma\) 散度及扩展的 Itakura–Saito 型距离族。本文特别讨论了调节参数取奇异值 \(0,1\) 时的散度。


这些散度族具有灵活性、对离群值的稳健性等性质，而且只涉及一个调节参数，即 \(\alpha\)、\(\beta\) 或 \(\gamma\)。信息几何的分析可以进一步阐明这些散度的基本结构与几何性质[\hyperlink{B1-entropy-12-01532}{1},\allowbreak \hyperlink{B3-entropy-12-01532}{3},\allowbreak \hyperlink{B7-entropy-12-01532}{7}]。


与之前密切相关的工作相比，本文明确研究了 \(\alpha\)、\(\beta\) 和 \(\gamma\) 三个广泛的散度类，讨论了它们的性质，以及它们与 Csiszár–Morimoto 散度、Bregman 散度的联系。本文给出的部分性质，尤其是 \(\gamma\) 散度的一些性质，是新的结果。


利用本文讨论的散度性质，我们开发了一族用于独立成分分析和非负矩阵分解的稳健算法，并在近期的其他工作中给出[\hyperlink{B30-entropy-12-01532}{30}--\hyperlink{B34-entropy-12-01532}{34},\allowbreak \hyperlink{B68-entropy-12-01532}{68}--\hyperlink{B71-entropy-12-01532}{71}]。它们在许多应用中具有潜力，尤其是对统计相关的稀疏源、平滑源或兼具二者的源进行盲分离，以及特征提取和分类问题[\hyperlink{B30-entropy-12-01532}{30}]。


一个重要但仍未解决、本文也未讨论的问题，是如何依据现有数据集的分布，以及噪声或离群值的情况，设定调节参数 \(\alpha,\beta\) 和 \(\gamma\)。近期一些工作研究了这个问题[\hyperlink{B69-entropy-12-01532}{69},\allowbreak \hyperlink{B71-entropy-12-01532}{71},\allowbreak \hyperlink{B84-entropy-12-01532}{84}]。通常，稳健性与统计效率之间存在权衡[\hyperlink{B35-entropy-12-01532}{35},\allowbreak \hyperlink{B37-entropy-12-01532}{37},\allowbreak \hyperlink{B71-entropy-12-01532}{71}]。例如，当调节参数 \(\gamma\) 较大时，\(\gamma\) 散度的稳健性较强，但统计效率可能较低；当该参数取较小的正值时，方法的稳健性较弱，但统计效率较高[\hyperlink{B35-entropy-12-01532}{35}]。统计效率与调节参数之间的关系，已由多位作者讨论，参见 Jones 等[\hyperlink{B85-entropy-12-01532}{85}]、Basu 等[\hyperlink{B67-entropy-12-01532}{67}]以及 Fujisawa 和 Eguchi[\hyperlink{B35-entropy-12-01532}{35}]。

\appendix
\section{\texorpdfstring{从 Tsallis \(\alpha\) 熵导出的散度}{从 Tsallis α 熵导出的散度}}\label{app1-entropy-12-01532}
本附录说明 Tsallis 熵和 Rényi 熵与 \(\alpha\) 散度、\(\beta\) 散度及 \(\gamma\) 散度之间的联系。为简化讨论，先假设 \(\int\overline{p}{(\mathbf{x})}d\mu{(\mathbf{x})} = 1\) 且 \(\int\overline{q}{(\mathbf{x})}d\mu{(\mathbf{x})} = 1\)，这些约束也可以放宽。


利用幂函数 \(h_{\alpha}{(\overline{\mathbf{p}})} = \int{\overline{p}}^{\alpha}d\mu{(\mathbf{x})}\)，Tsallis 熵定义为

\setcounter{equation}{99}
\begin{equation}
\label{FD126-entropy-12-01532}
\fitmath{H_{T}^{(\alpha)}{(\overline{\mathbf{p}})} = \frac{1}{1 - \alpha}\left\{ {h_{\alpha}{(\overline{\mathbf{p}})} - 1} \right\}}
\end{equation}

Rényi 熵则可写为

\setcounter{equation}{100}
\begin{equation}
\label{FD127-entropy-12-01532}
\fitmath{H_{R}^{(\alpha)}{(\overline{\mathbf{p}})} = \frac{1}{1 - \alpha}\log\left\{ {h_{\alpha}{(\overline{\mathbf{p}})}} \right\} = \frac{1}{1 - \alpha}\log\left( {\int{\overline{p}}^{\alpha}{(\mathbf{x})}\, d\mu{(\mathbf{x})}} \right)}
\end{equation}




Tsallis 熵与 \(\alpha\) 散度和 \(\beta\) 散度相对应，而 Rényi 熵与 \(\gamma\) 散度相关。


Rényi 熵和 Tsallis 熵是标准 Boltzmann–Gibbs 熵，即 Shannon 信息熵，的两种不同推广。在 \(\left. \alpha\rightarrow 1 \right.\) 的极限下，有 \(p^{\alpha - 1} = \exp{({(\alpha - 1)}\log p)} \approx 1 + {(\alpha - 1)}\log p\)，二者均退化为 Shannon 熵：

\setcounter{equation}{101}
\begin{equation}
\label{FD128-entropy-12-01532}
\fitmath{H_{S}{(\overline{\mathbf{p}})} = - \int\overline{p}{(\mathbf{x})}\log{(\overline{p}{(\mathbf{x})})}\, d\mu{(\mathbf{x})}}
\end{equation}

Tsallis 熵还可以视为 Shannon 熵的 \(\alpha\) 形变，因为

\setcounter{equation}{102}
\begin{equation}
\label{FD129-entropy-12-01532}
\fitmath{H_{T}^{(\alpha)}{(\overline{\mathbf{p}})} = - \int{\overline{p}}^{\alpha}{(\mathbf{x})}\log_{\,\alpha}\,{(\overline{p}{(\mathbf{x})})}\, d\mu{(\mathbf{x})} = \frac{1}{1 - \alpha}\left( {\int{\overline{p}}^{\,\alpha}{(\mathbf{x})}d\mu{(\mathbf{x})} - 1} \right)}
\end{equation}

其中使用的广义对数函数定义为

\setcounter{equation}{103}
\begin{equation}
\label{FD130-entropy-12-01532}
\fitmath{\log_{\,\alpha}x = \frac{1}{1 - \alpha}\left( {x^{1 - \alpha} - 1} \right)}
\end{equation}




类似地，Tsallis 相对熵或 Tsallis 散度定义为

\setcounter{equation}{104}
\begin{equation}
\label{FD131-entropy-12-01532}
\fitmath{\begin{array}{rcl}
{D_{T}^{(\alpha)}\left( {\overline{\mathbf{p}} \parallel \overline{\mathbf{q}}} \right)} & = & {- \int\overline{p}\log_{\,\alpha}\left( \frac{\overline{q}}{\overline{p}} \right)d\mu{(\mathbf{x})} = \frac{1}{1 - \alpha}\left( {1 - \int{\overline{p}}^{\alpha}{\overline{q}}^{1 - \alpha}d\mu{(\mathbf{x})}} \right)}
\end{array}}
\end{equation}




这是 \(\alpha\) 散度经过重新缩放后的形式：

\setcounter{equation}{105}
\begin{equation}
\label{FD132-entropy-12-01532}
\fitmath{D_{T}^{(\alpha)}{(\overline{\mathbf{p}} \parallel \overline{\mathbf{q}})} = \alpha D_{A}^{(\alpha)}\left( {\overline{\mathbf{p}} \parallel \overline{\mathbf{q}}} \right)}
\end{equation}

\(\alpha\) 散度更为一般，因为它对任意正数组 \(\mathbf{P}\) 和 \(\mathbf{Q}\) 都有定义，而且允许 \(\alpha\) 取任意值，包括奇异值 \(\alpha = 0\) 和 \(\alpha = 1\)。


由于 Tsallis 的 \(\alpha\) 散度是 f 散度，可以将其扩展到相应的一般情形。\hyperref[app2-entropy-12-01532]{附录 B}给出了这种扩展，其中 \(\log_{\,\alpha}{(x)}\) 被另一个幂函数替代。


导出散度的另一种方法是使用凸生成函数。Zhang 对此作了深入研究[\hyperlink{B22-entropy-12-01532}{22}]。


由于

\setcounter{equation}{106}
\begin{equation}
\label{FD133-entropy-12-01532}
\fitmath{\Phi_{\alpha}{(\overline{\mathbf{p}})} = - h_{\alpha}{(\overline{\mathbf{p}})}}
\end{equation}

为凸函数，可以导出与 Tsallis 熵对应的 Bregman 散度：

\setcounter{equation}{107}
\begin{equation}
\label{FD134-entropy-12-01532}
\fitmath{{\widetilde{D}}_{T}^{(\alpha)}{(\overline{\mathbf{p}}\, \parallel \,\overline{\mathbf{q}})} = \int\left\{ {- {\overline{p}}^{\alpha} + {\overline{q}}^{\alpha} + \alpha{\overline{p}}^{\alpha - 1}\left( {\overline{p} - \overline{q}} \right)} \right\} d\mu{(\mathbf{x})}}
\end{equation}




这是在 \(\alpha = \beta\) 条件下，\(\beta\) 散度经过重新缩放和简化后的形式：

\setcounter{equation}{108}
\begin{equation}
\label{FD135-entropy-12-01532}
\fitmath{{\widetilde{D}}_{T}^{(\alpha)}{(\overline{\mathbf{p}}\, \parallel \,\overline{\mathbf{q}})} = \alpha{(1 - \alpha)}D_{B}^{(\beta)}\left( {\overline{\mathbf{p}} \parallel \overline{\mathbf{q}}} \right)}
\end{equation}




\(\beta\) 散度也比 Tsallis 散度更为一般，因为它定义在任意正数组上，并允许 \(\beta\) 取任意值，包括奇异值 \(\beta = 0\) 和 \(\beta = 1\)。要得到完整的 \(\beta\) 散度取值范围，需要扩展 \(h_{\alpha}{(\mathbf{p})}\) 的定义。

\section{\texorpdfstring{熵与散度}{熵与散度}}\label{app2-entropy-12-01532}
熵 \(H(\overline{\mathbf{p}})\) 在均匀分布处取得最大值。因此，给定一个散度函数，可以定义熵 \(H(\overline{\mathbf{p}})\) 为

\setcounter{equation}{109}
\begin{equation}
\label{FD107-entropy-12-01532}
\fitmath{H{(\overline{\mathbf{p}})}{= D(}\overline{\mathbf{p}}{\,||\, 1)} + c}
\end{equation}

其中 c 为某个常数，\(\overline{p}\) 为归一化概率测度。


在许多情形中，\(\Psi{(\overline{\mathbf{p}})}{= - D(}\overline{\mathbf{p}}{\,||\, 1)}\) 是 \(\overline{p}\) 的凸函数。例如，由 KL 散度可得

\setcounter{equation}{110}
\begin{equation}
\label{FD108-entropy-12-01532}
\fitmath{H_{KL}{(\overline{\mathbf{p}})} = - \int\overline{p}{(\mathbf{x})}\log\overline{p}{(\mathbf{x})}d\mu{(\mathbf{x})}}
\end{equation}

忽略一个常数后，这就是 Shannon 熵。


定义概率的幂函数为

\setcounter{equation}{111}
\begin{equation}
\label{FD109-entropy-12-01532}
\fitmath{h_{\alpha}{(\overline{\mathbf{p}})} = \int{\overline{p}}^{\alpha}{(\mathbf{x})}d\mu{(\mathbf{x})}}
\end{equation}

其中 \(0 < \alpha < 1\)。


由 \(\alpha\) 散度 \(D_{A}^{(\alpha)}\) 和 \(\beta\) 散度 \(D_{B}^{(\beta)}\)，得到相应的熵

\setcounter{equation}{112}
\begin{equation}
\label{FD110-entropy-12-01532}
\fitmath{H{(\overline{\mathbf{p}})} = c\, h_{\alpha}{(\overline{\mathbf{p}})} + d}
\end{equation}

其中 c 和 d 为常数，且 \(\alpha = \beta\)。


类似地，由 \(\gamma\) 散度 \(D_{G}^{(\gamma)}\) 得到

\setcounter{equation}{113}
\begin{equation}
\label{FD111-entropy-12-01532}
\fitmath{H{(\overline{\mathbf{p}})} = c\,\log{(h_{\alpha}{(\overline{\mathbf{p}})})} + d}
\end{equation}

其中 \(\alpha = \gamma\)。


不同的散度函数可以生成相似的熵函数族。反过来，从熵函数生成散度也有多种方法。一种典型方法，是利用凸生成函数 \(\Phi{(\overline{p})} = - H{(\overline{p})}\) 导出 Bregman 散度：

\setcounter{equation}{114}
\begin{equation}
\label{FD112-entropy-12-01532}
\fitmath{D_{\Phi}(\overline{\mathbf{p}}{\,||\,}\overline{\mathbf{q}}{) = \int}\left( {\Phi{(\overline{p})} - \Phi{(\overline{q})} - \frac{\delta\Phi(\overline{q})}{\delta\overline{q}}{(\overline{p} - \overline{q})}} \right)\, d\mu{(\mathbf{x})}}
\end{equation}

\section{\texorpdfstring{Tsallis 熵与 Rényi 熵}{Tsallis 熵与 Rényi 熵}}\label{app3-entropy-12-01532}
Rényi 的 \(\alpha\) 熵与 Tsallis 的 \(\alpha\) 熵通过以下单调关系相联系：

\setcounter{equation}{115}
\begin{equation}
\label{FD113-entropy-12-01532}
\fitmath{H_{R}^{(\alpha)}{(\overline{\mathbf{p}})} = \frac{1}{1 - \alpha}\log\left( {1 + {(1 - \alpha)}H_{T}^{(\alpha)}{(\overline{\mathbf{p}})}} \right)}
\end{equation}

或者使用 \(\alpha\) 对数写为

\setcounter{equation}{116}
\begin{equation}
\label{FD114-entropy-12-01532}
\fitmath{H_{T}^{(\alpha)}{(\overline{\mathbf{p}})} = \log_{\,\alpha}e^{\; H_{R}^{(\alpha)}{(\overline{\mathbf{p}})}}}
\end{equation}

\(\Psi_{R}^{(\alpha)} = - H_{R}^{(\alpha)}\) 是 \(\overline{\mathbf{p}}\) 的凸函数，因此可以构造相应的散度。


Bregman 散度可以定义为以下更一般的形式。


若 \(\left. \Phi{(\overline{\mathbf{p}})}:\;\mathbb{R}^{n}\rightarrow\mathbb{R} \right.\) 是严格凸的 \(C^{1}\) 函数，即一阶可微函数，则离散正测度上相应的 Bregman 距离或散度定义为

\setcounter{equation}{117}
\begin{equation}
\label{FD115-entropy-12-01532}
\fitmath{D_{\Phi}(\overline{\mathbf{p}}{\,||\,}\overline{\mathbf{q}}{) = \Phi}{(\overline{\mathbf{p}})} - \Phi{(\overline{\mathbf{q}})} - {(\overline{\mathbf{p}} - \overline{\mathbf{q}})}^{T}\nabla\Phi{(\overline{\mathbf{q}})}}
\end{equation}

其中 \(\nabla\Phi(\overline{\mathbf{q}})\) 表示 \(\Phi\) 在 \(\overline{\mathbf{q}}\) 处的梯度。


由 Rényi 熵导出的离散概率密度上的 Bregman 散度为

\setcounter{equation}{118}
\begin{equation}
\label{FD116-entropy-12-01532}
\fitmath{D_{R}^{(\alpha)}(\overline{\mathbf{p}}{\,||\,}\overline{\mathbf{q}}{) =}\frac{1}{1 - \alpha}\log\frac{h_{\alpha}{(\overline{\mathbf{q}})}}{h_{\alpha}{(\mathbf{p})}} - \frac{\alpha}{1 - \alpha}\frac{1}{h_{\alpha}{(\overline{\mathbf{q}})}}{(1 - \sum\limits_{i}{\overline{p}}_{i}\,{\overline{q}}_{i}^{\alpha - 1})}}
\end{equation}

Tsallis 熵与概率分布的 \(\alpha\) 指数族密切相关。由向量参数 \(\mathbf{\theta}\) 参数化的一族概率分布 \(S = \left\{ {\overline{p}{(\mathbf{x},\mathbf{\theta})}} \right\}\)，若能写为

\setcounter{equation}{119}
\begin{equation}
\label{FD117-entropy-12-01532}
\fitmath{\log_{\,\alpha}\overline{p}{(\mathbf{x},\mathbf{\theta})} = \sum\limits_{i}\theta_{i}x_{i} - \Psi{(\mathbf{\theta})}}
\end{equation}

其中 \(\Psi(\mathbf{\theta})\) 为凸函数，则称为 \(\alpha\) 指数族。这与文献[\hyperlink{B6-entropy-12-01532}{6}]定义的 \(\alpha\) 族相同。对于任意 \(\alpha\)，离散概率分布集合 \(S = \left\{ \mathbf{p} \right\}\) 都构成 \(\alpha\) 指数族。为说明这一点，引入随机变量 \(\mathbf{x} = (x)\)：

\setcounter{equation}{120}
\begin{equation}
\label{FD118-entropy-12-01532}
\fitmath{x_{i} = \delta_{i}}
\end{equation}

当结果为 i 时，它取值为 1，该结果的概率为 \(\overline{p}\)。


于是，概率可以写为

\setcounter{equation}{121}
\begin{equation}
\label{FD119-entropy-12-01532}
\fitmath{\log_{\,\alpha}\overline{p}{(\mathbf{x},\mathbf{\theta})} = \sum\limits_{i}\theta_{i}\delta_{i} - \Psi{(\mathbf{\theta})}}
\end{equation}

其中

\setcounter{equation}{122}
\begin{equation}
\label{FD120-entropy-12-01532}
\fitmath{\theta_{i} = \log_{\,\alpha}{\overline{p}}_{i}}
\end{equation}

由于 \(\sum{\overline{p}}_{i} = 1\) 成立，或者写成

\setcounter{equation}{123}
\begin{equation}
\label{FD121-entropy-12-01532}
\fitmath{\sum\delta_{i} = 1}
\end{equation}

因此，\(\mathbf{\theta}\) 的自由变量为 \(n - 1\)，记为 \(\left( {\theta_{1},\cdots,\theta_{n - 1}} \right)\)。此时，\(\theta_{n}\) 是其余变量的函数，得到

\setcounter{equation}{124}
\begin{equation}
\label{FD122-entropy-12-01532}
\fitmath{\Psi{(\mathbf{\theta})} = - \log_{\,\alpha}\theta_{n}}
\end{equation}




可以由这个凸函数 \(\Psi(\mathbf{\theta})\) 导出散度。不过，使用 \(\Psi(\mathbf{\theta})\) 的对偶函数更直观，该对偶函数为

\setcounter{equation}{125}
\begin{equation}
\label{FD123-entropy-12-01532}
\fitmath{\varphi{(\mathbf{\eta})} = \frac{1}{1 - \alpha}\left( {\frac{1}{h_{\alpha}{(\overline{\mathbf{p}})}} - 1} \right)}
\end{equation}

其中 \(\mathbf{\eta}\) 是 \(\mathbf{\theta}\) 通过 Legendre 变换得到的对偶变量：

\setcounter{equation}{126}
\begin{equation}
\label{FD124-entropy-12-01532}
\fitmath{\mathbf{\eta} = \nabla_{\mathbf{\theta}}\Psi{(\mathbf{\theta})}}
\end{equation}




得到的散度为

\setcounter{equation}{127}
\begin{equation}
\label{FD125-entropy-12-01532}
\fitmath{D_{\exp}^{(\alpha)}{(\overline{\mathbf{p}}\, \parallel \,\overline{\mathbf{q}})} = \frac{1}{h_{\alpha}{(\overline{\mathbf{p}})}}D_{T}^{(\alpha)}{(\overline{\mathbf{p}}\, \parallel \,\overline{\mathbf{q}})}}
\end{equation}

这是 \(\alpha\) 散度的共形变换[\hyperlink{B7-entropy-12-01532}{7}]。共形变换也称共形映射或保角变换，它保持任意两条直线或曲线在交点处的局部夹角。

\begin{thebibliography}{85}
\raggedright
\bibitem{B1-entropy-12-01532}\hypertarget{B1-entropy-12-01532}{} Amari S. Differential-Geometrical Methods in Statistics. Springer Verlag. Berlin, Germany. 1985.
\bibitem{B2-entropy-12-01532}\hypertarget{B2-entropy-12-01532}{} Amari S. Dualistic geometry of the manifold of higher-order neurons. Neural Network. 1991. 4. pp.~443--451. \href{https://doi.org/10.1016/0893-6080(91)90040-C}{doi: 10.1016/0893-6080(91)90040-C}.
\bibitem{B3-entropy-12-01532}\hypertarget{B3-entropy-12-01532}{} Amari S.; Nagaoka H. Methods of Information Geometry. Oxford University Press. New York, NY, USA. 2000.
\bibitem{B4-entropy-12-01532}\hypertarget{B4-entropy-12-01532}{} Amari S. Integration of stochastic models by minimizing alpha-divergence. Neural Comput. 2007. 19. pp.~2780--2796. \href{https://doi.org/10.1162/neco.2007.19.10.2780}{doi: 10.1162/neco.2007.19.10.2780}.
\bibitem{B5-entropy-12-01532}\hypertarget{B5-entropy-12-01532}{} Amari S. Information geometry and its applications: Convex function and dually flat manifold. Emerging Trends in Visual Computing. Nielsen F. Springer. New York, NY, USA. pp.~75--102.
\bibitem{B6-entropy-12-01532}\hypertarget{B6-entropy-12-01532}{} Amari S. Alpha-divergence is unique, belonging to both f-divergence and Bregman divergence classes. IEEE Trans. Inf. Theory. 2009. 55. pp.~4925--4931. \href{https://doi.org/10.1109/TIT.2009.2030485}{doi: 10.1109/TIT.2009.2030485}.
\bibitem{B7-entropy-12-01532}\hypertarget{B7-entropy-12-01532}{} Amari S.; Cichocki A. Information geometry of divergence functions. Bull. Pol. Acad. Sci. 2010. (in print). \href{https://doi.org/10.2478/v10175-010-0019-1}{doi: 10.2478/v10175-010-0019-1}.
\bibitem{B8-entropy-12-01532}\hypertarget{B8-entropy-12-01532}{} Murata N.; Takenouchi T.; Kanamori T.; Eguchi S. Information geometry of U-Boost and Bregman divergence. Neural Comput. 2004. 16. pp.~1437--1481. \href{https://doi.org/10.1162/089976604323057452}{doi: 10.1162/089976604323057452}.
\bibitem{B9-entropy-12-01532}\hypertarget{B9-entropy-12-01532}{} Fujimoto Y.; Murata N. A modified EM Algorithm for mixture models based on Bregman divergence. Ann. Inst. Stat. Math. 2007. 59. pp.~57--75. \href{https://doi.org/10.1007/s10463-006-0097-x}{doi: 10.1007/s10463-006-0097-x}.
\bibitem{B10-entropy-12-01532}\hypertarget{B10-entropy-12-01532}{} Zhu H.; Rohwer R. Bayesian Invariant measurements of generalization. Neural Process. Lett. 1995. 2. pp.~28--31. \href{https://doi.org/10.1007/BF02309013}{doi: 10.1007/BF02309013}.
\bibitem{B11-entropy-12-01532}\hypertarget{B11-entropy-12-01532}{} Zhu H.; Rohwer R. Measurements of generalisation based on information geometry. Mathematics of Neural Networks: Model Algorithms and Applications. Ellacott S.W.; Mason J.C.; Anderson I.J. Kluwer. Norwell, MA, USA. 1997. pp.~394--398.
\bibitem{B12-entropy-12-01532}\hypertarget{B12-entropy-12-01532}{} Nielsen F.; Nock R. Sided and symmetrized Bregman centroids. IEEE Trans. Inf. Theory. 2009. 56. pp.~2882--2903. \href{https://doi.org/10.1109/TIT.2009.2018176}{doi: 10.1109/TIT.2009.2018176}.
\bibitem{B13-entropy-12-01532}\hypertarget{B13-entropy-12-01532}{} Boissonnat J.D.; Nielsen F.; Nock R. Bregman Voronoi diagrams. Discrete and Computational Geometry (Springer). 2010. (in print). \href{https://doi.org/10.1007/s00454-010-9256-1}{doi: 10.1007/s00454-010-9256-1}.
\bibitem{B14-entropy-12-01532}\hypertarget{B14-entropy-12-01532}{} Yamano T. A generalization of the Kullback-Leibler divergence and its properties. J. Math. Phys. 2009. 50. pp.~85--95. \href{https://doi.org/10.1063/1.3116115}{doi: 10.1063/1.3116115}.
\bibitem{B15-entropy-12-01532}\hypertarget{B15-entropy-12-01532}{} Minami M.; Eguchi S. Robust blind source separation by Beta-divergence. Neural Comput. 2002. 14. pp.~1859--1886.
\bibitem{B16-entropy-12-01532}\hypertarget{B16-entropy-12-01532}{} Bregman L. The relaxation method of finding a common point of convex sets and its application to the solution of problems in convex programming. Comp. Math. Phys., USSR. 1967. 7. pp.~200--217. \href{https://doi.org/10.1016/0041-5553(67)90040-7}{doi: 10.1016/0041-5553(67)90040-7}.
\bibitem{B17-entropy-12-01532}\hypertarget{B17-entropy-12-01532}{} Csiszár I. Eine Informations Theoretische Ungleichung und ihre Anwendung auf den Beweis der Ergodizitt von Markoffschen Ketten. Magyar Tud. Akad. Mat. Kutat Int. Kzl. 1963. 8. pp.~85--108.
\bibitem{B18-entropy-12-01532}\hypertarget{B18-entropy-12-01532}{} Csiszár I. Axiomatic characterizations of information measures. Entropy. 2008. 10. pp.~261--273. \href{https://doi.org/10.3390/e10030261}{doi: 10.3390/e10030261}.
\bibitem{B19-entropy-12-01532}\hypertarget{B19-entropy-12-01532}{} Csiszár I. Information measures: A critial survey. Transactions of the 7th Prague Conference. Prague, Czech Republic. 18–23 August 1974. Reidel. Dordrecht, Netherlands. 1977. pp.~83--86.
\bibitem{B20-entropy-12-01532}\hypertarget{B20-entropy-12-01532}{} Ali M.; Silvey S. A general class of coefficients of divergence of one distribution from another. J. Royal Stat. Soc. 1966. Ser B. pp.~131--142.
\bibitem{B21-entropy-12-01532}\hypertarget{B21-entropy-12-01532}{} Hein M.; Bousquet O. Hilbertian metrics and positive definite kernels on probability measures. Proceedings of the Tenth International Workshop on Artificial Intelligence and Statistics. Barbados. 6–8 January 2005. Ghahramani Z.; Cowell R. AISTATS. 2005. 10. pp.~136--143.
\bibitem{B22-entropy-12-01532}\hypertarget{B22-entropy-12-01532}{} Zhang J. Divergence function, duality, and convex analysis. Neural Comput. 2004. 16. pp.~159--195. \href{https://doi.org/10.1162/08997660460734047}{doi: 10.1162/08997660460734047}.
\bibitem{B23-entropy-12-01532}\hypertarget{B23-entropy-12-01532}{} Zhang J. Referential duality and representational duality on statistical manifolds. Proceedings of the Second International Symposium on Information Geometry and its Applications. University of Tokyo, Tokyo, Japan. 12–16 December 2005. University of Tokyo. Tokyo, Japan. 2006. pp.~58--67.
\bibitem{B24-entropy-12-01532}\hypertarget{B24-entropy-12-01532}{} Zhang J. A note on curvature of a-connections of a statistical manifold. Ann. Inst. Stat. Math. 2007. 59. pp.~161--170. \href{https://doi.org/10.1007/s10463-006-0105-1}{doi: 10.1007/s10463-006-0105-1}.
\bibitem{B25-entropy-12-01532}\hypertarget{B25-entropy-12-01532}{} Zhang J.; Matsuzoe H. Dualistic differential geometry associated with a convex function. Springer Series of Advances in Mechanics and Mathematics. 2008. Springer. New York, NY, USA. pp.~58--67.
\bibitem{B26-entropy-12-01532}\hypertarget{B26-entropy-12-01532}{} Lafferty J. Additive models, boosting, and inference for generalized divergences. Proceedings of the Twelfth Annual Conference on Computational Learning Theory. Santa Cruz, CA, USA. 7–9 July 1999. ACM. New York, NY, USA. 1999.
\bibitem{B27-entropy-12-01532}\hypertarget{B27-entropy-12-01532}{} Banerjee A.; Merugu S.; Dhillon I.S.; Ghosh J. Clustering with Bregman divergences. J. Mach. Learn. Res. 2005. 6. pp.~1705--1749.
\bibitem{B28-entropy-12-01532}\hypertarget{B28-entropy-12-01532}{} Villmann T.; Haase S. Divergence based vector quantization using Fréchet derivatives. Neural Comput. 2010. (submitted for publication).
\bibitem{B29-entropy-12-01532}\hypertarget{B29-entropy-12-01532}{} Villmann T.; Haase S.; Schleif F.M.; Hammer B. Divergence based online learning in vector quantization. Proceedings of the International Conference on Artifial Intelligence and Soft Computing (ICAISC2010). LNAI, Zakopane, Poland. 13–17 June 2010.
\bibitem{B30-entropy-12-01532}\hypertarget{B30-entropy-12-01532}{} Cichocki A.; Zdunek R.; Phan A.H.; Amari S. Nonnegative Matrix and Tensor Factorizations. John Wiley \& Sons Ltd. Chichester, UK. 2009.
\bibitem{B31-entropy-12-01532}\hypertarget{B31-entropy-12-01532}{} Cichocki A.; Zdunek R.; Amari S. Csiszár’s divergences for nonnegative matrix factorization: Family of new algorithms. Springer, LNCS-3889. 2006. 3889. pp.~32--39.
\bibitem{B32-entropy-12-01532}\hypertarget{B32-entropy-12-01532}{} Cichocki A.; Amari S.; Zdunek R.; Kompass R.; Hori G.; He Z. Extended SMART algorithms for Nonnegative Matrix Factorization. Springer, LNAI-4029. 2006. 4029. pp.~548--562.
\bibitem{B33-entropy-12-01532}\hypertarget{B33-entropy-12-01532}{} Cichocki A.; Zdunek R.; Choi S.; Plemmons R.; Amari S. Nonnegative tensor factorization using Alpha and Beta divergences. Proceedings of the IEEE International Conference on Acoustics, Speech, and Signal Processing. Tulose, France. May 2007. Volume III. pp.~1393--1396.
\bibitem{B34-entropy-12-01532}\hypertarget{B34-entropy-12-01532}{} Cichocki A.; Zdunek R.; Choi S.; Plemmons R.; Amari S.I. Novel multi-layer nonnegative tensor factorization with sparsity constraints. Springer, LNCS-4432. 2007. 4432. pp.~271--280.
\bibitem{B35-entropy-12-01532}\hypertarget{B35-entropy-12-01532}{} Fujisawa H.; Eguchi S. Robust parameter estimation with a small bias against heavy contamination. J. Multivar. Anal. 2008. 99. pp.~2053--2081. \href{https://doi.org/10.1016/j.jmva.2008.02.004}{doi: 10.1016/j.jmva.2008.02.004}.
\bibitem{B36-entropy-12-01532}\hypertarget{B36-entropy-12-01532}{} Liese F.; Vajda I. Convex Statistical Distances. Teubner-Texte zur Mathematik Teubner Texts in Mathematics. 1987. 95. pp.~1--85.
\bibitem{B37-entropy-12-01532}\hypertarget{B37-entropy-12-01532}{} Eguchi S.; Kato S. Entropy and divergence associated with power function and the statistical application. Entropy. 2010. 12. pp.~262--274. \href{https://doi.org/10.3390/e12020262}{doi: 10.3390/e12020262}.
\bibitem{B38-entropy-12-01532}\hypertarget{B38-entropy-12-01532}{} Taneja I. On generalized entropies with applications. Lectures in Applied Mathematics and Informatics. Ricciardi L. Manchester University Press. Manchester, UK. 1990. pp.~107--169.
\bibitem{B39-entropy-12-01532}\hypertarget{B39-entropy-12-01532}{} Taneja I. New developments in generalized information measures. Advances in Imaging and Electron Physics. Hawkes P. Elsevier. Amsterdam, Netherlands. 1995. Volume 91. pp.~37--135.
\bibitem{B40-entropy-12-01532}\hypertarget{B40-entropy-12-01532}{} Gorban A.N.; Gorban P.A.; Judge G. Entropy: The Markov ordering approach. Entropy. 2010. 12. pp.~1145--1193. \href{https://doi.org/10.3390/e12051145}{doi: 10.3390/e12051145}.
\bibitem{B41-entropy-12-01532}\hypertarget{B41-entropy-12-01532}{} Chernoff H. A measure of asymptotic efficiency for tests of a hypothesis based on a sum of observations. Ann. Math. Statist. 1952. 23. pp.~493--507. \href{https://doi.org/10.1214/aoms/1177729330}{doi: 10.1214/aoms/1177729330}.
\bibitem{B42-entropy-12-01532}\hypertarget{B42-entropy-12-01532}{} Minka T. Divergence measures and message passing. Microsoft Research Technical Report (MSR-TR-2005). 2005.
\bibitem{B43-entropy-12-01532}\hypertarget{B43-entropy-12-01532}{} Taneja I. On measures of information and inaccuarcy. J. Statist. Phys. 1976. 14. pp.~203--270.
\bibitem{B44-entropy-12-01532}\hypertarget{B44-entropy-12-01532}{} Cressie N.; Read T. Goodness-of-Fit Statistics for Discrete Multivariate Data. Springer. New York, NY, USA. 1988.
\bibitem{B45-entropy-12-01532}\hypertarget{B45-entropy-12-01532}{} Cichocki A.; Lee H.; Kim Y.D.; Choi S. Nonnegative matrix factorization with Alpha-divergence. Pattern. Recognit. Lett. 2008. 29. pp.~1433--1440. \href{https://doi.org/10.1016/j.patrec.2008.02.016}{doi: 10.1016/j.patrec.2008.02.016}.
\bibitem{B46-entropy-12-01532}\hypertarget{B46-entropy-12-01532}{} Tsallis C. Possible generalization of Boltzmann-Gibbs statistics. J. Statist. Phys. 1988. 52. pp.~479--487. \href{https://doi.org/10.1007/BF01016429}{doi: 10.1007/BF01016429}.
\bibitem{B47-entropy-12-01532}\hypertarget{B47-entropy-12-01532}{} Havrda J.; Charvát F. Quantification method of classification processes: Concept of structrual a-entropy. Kybernetika. 1967. 3. pp.~30--35.
\bibitem{B48-entropy-12-01532}\hypertarget{B48-entropy-12-01532}{} Cressie N.; Read T. Multinomial goodness-of-fit tests. J. R. Stat. Soc. Ser. B. 1984. 46. pp.~440--464.
\bibitem{B49-entropy-12-01532}\hypertarget{B49-entropy-12-01532}{} Vajda I. Theory of Statistical Inference and Information. Kluwer Academic Press. Amsterdam, Netherland. 1989.
\bibitem{B50-entropy-12-01532}\hypertarget{B50-entropy-12-01532}{} Hellinger E. Neue Begründung der Theorie Quadratischen Formen von unendlichen vielen Veränderlichen. J. Reine Ang. Math. 1909. 136. pp.~210--271.
\bibitem{B51-entropy-12-01532}\hypertarget{B51-entropy-12-01532}{} Morimoto T. Markov processes and the H-theorem. J. Phys. Soc. Jap. 1963. 12. pp.~328--331. \href{https://doi.org/10.1143/JPSJ.18.328}{doi: 10.1143/JPSJ.18.328}.
\bibitem{B52-entropy-12-01532}\hypertarget{B52-entropy-12-01532}{} Österreicher F. Csiszár’s f-divergences-basic properties. Technical report. Research Report Collection. Victoria University. Melbourne, Australia. 2002.
\bibitem{B53-entropy-12-01532}\hypertarget{B53-entropy-12-01532}{} Harremoës P.; Vajda I. Joint range of f-divergences. Accepted for presentation at ISIT 2010. Austin, TX, USA. 13–18 June 2010.
\bibitem{B54-entropy-12-01532}\hypertarget{B54-entropy-12-01532}{} Dragomir S. Inequalities for Csiszár f-Divergence in Information Theory. Victoria University. Melbourne, Australia. 2000. (edited monograph).
\bibitem{B55-entropy-12-01532}\hypertarget{B55-entropy-12-01532}{} Rényi A. On the foundation of information theory. Rev. Inst. Int. Stat. 1965. 33. pp.~1--4. \href{https://doi.org/10.2307/1401301}{doi: 10.2307/1401301}.
\bibitem{B56-entropy-12-01532}\hypertarget{B56-entropy-12-01532}{} Rényi A. On measures of entropy and information. Proceddings of the Fourth Berkeley Symposium on Mathematical Statistics and Probability. Berkeley, CA, USA. 20 June–30 July 1960. University of California Press. Berkeley, CA, USA. Volome 1. pp.~547--561.
\bibitem{B57-entropy-12-01532}\hypertarget{B57-entropy-12-01532}{} Rényi A. Probability Theory. North-Holland. Amsterdam, The Netherlands. 1970.
\bibitem{B58-entropy-12-01532}\hypertarget{B58-entropy-12-01532}{} Harremoës P. Interpretaions of Rényi entropies and divergences. Physica A. 2006. 365. pp.~57--62. \href{https://doi.org/10.1016/j.physa.2006.01.012}{doi: 10.1016/j.physa.2006.01.012}.
\bibitem{B59-entropy-12-01532}\hypertarget{B59-entropy-12-01532}{} Harremoës P. Joint range of Rényi entropies. Kybernetika. 2009. 45. pp.~901--911.
\bibitem{B60-entropy-12-01532}\hypertarget{B60-entropy-12-01532}{} Hero A.; Ma B.; Michel O.; Gorman J. Applications of entropic spanning graphs. IEEE Signal Process. Mag. 2002. 19. pp.~85--95. \href{https://doi.org/10.1109/MSP.2002.1028355}{doi: 10.1109/MSP.2002.1028355}.
\bibitem{B61-entropy-12-01532}\hypertarget{B61-entropy-12-01532}{} Topsoe F. Some inequalities for information divergence and related measuresof discrimination. IEEE Trans. Inf. Theory. 2000. 46. pp.~1602--1609. \href{https://doi.org/10.1109/18.850703}{doi: 10.1109/18.850703}.
\bibitem{B62-entropy-12-01532}\hypertarget{B62-entropy-12-01532}{} Burbea J.; Rao C. Entropy differential metric, distance and divergence measures in probability spaces: A unified approach. J. Multi. Analysis. 1982. 12. pp.~575--596. \href{https://doi.org/10.1016/0047-259X(82)90065-3}{doi: 10.1016/0047-259X(82)90065-3}.
\bibitem{B63-entropy-12-01532}\hypertarget{B63-entropy-12-01532}{} Burbea J.; Rao C. On the convexity of some divergence measures based on entropy functions. IEEE Trans. Inf. Theory. 1982. IT-28. pp.~489--495. \href{https://doi.org/10.1109/TIT.1982.1056497}{doi: 10.1109/TIT.1982.1056497}.
\bibitem{B64-entropy-12-01532}\hypertarget{B64-entropy-12-01532}{} Sibson R. Information radius. Probability Theory and Related Fields. 1969. 14. pp.~149--160. \href{https://doi.org/10.1007/BF00537520}{doi: 10.1007/BF00537520}.
\bibitem{B65-entropy-12-01532}\hypertarget{B65-entropy-12-01532}{} Jeffreys H. An invariant form for the prior probability in estimation problems. Proc. Roy. Soc. Lon., Ser. A. 1946. 186. pp.~453--461. \href{https://doi.org/10.1098/rspa.1946.0056}{doi: 10.1098/rspa.1946.0056}.
\bibitem{B66-entropy-12-01532}\hypertarget{B66-entropy-12-01532}{} Kullback S.; Leibler R. On information and sufficiency. Ann. Math. Statist. 1951. 22. pp.~79--86. \href{https://doi.org/10.1214/aoms/1177729694}{doi: 10.1214/aoms/1177729694}.
\bibitem{B67-entropy-12-01532}\hypertarget{B67-entropy-12-01532}{} Basu A.; Harris I.R.; Hjort N.; Jones M. Robust and efficient estimation by minimising a density power divergence. Biometrika. 1998. 85. pp.~549--559. \href{https://doi.org/10.1093/biomet/85.3.549}{doi: 10.1093/biomet/85.3.549}.
\bibitem{B68-entropy-12-01532}\hypertarget{B68-entropy-12-01532}{} Mollah M.; Minami M.; Eguchi S. Exploring latent structure of mixture ICA models by the minimum Beta-divergence method. Neural Comput. 2006. 16. pp.~166--190. \href{https://doi.org/10.1162/089976606774841549}{doi: 10.1162/089976606774841549}.
\bibitem{B69-entropy-12-01532}\hypertarget{B69-entropy-12-01532}{} Mollah M.; Eguchi S.; Minami M. Robust prewhitening for ICA by minimizing Beta-divergence and its application to FastICA. Neural Process. Lett. 2007. 25. pp.~91--110. \href{https://doi.org/10.1007/s11063-006-9023-8}{doi: 10.1007/s11063-006-9023-8}.
\bibitem{B70-entropy-12-01532}\hypertarget{B70-entropy-12-01532}{} Kompass R. A Generalized divergence measure for Nonnegative Matrix Factorization. Neural Comput. 2006. 19. pp.~780--791. \href{https://doi.org/10.1162/neco.2007.19.3.780}{doi: 10.1162/neco.2007.19.3.780}.
\bibitem{B71-entropy-12-01532}\hypertarget{B71-entropy-12-01532}{} Mollah M.; Sultana N.; Minami M.; Eguchi S. Robust extraction of local structures by the minimum of Beta-divergence method. Neural Netw. 2010. 23. pp.~226--238. \href{https://doi.org/10.1016/j.neunet.2009.11.011}{doi: 10.1016/j.neunet.2009.11.011}.
\bibitem{B72-entropy-12-01532}\hypertarget{B72-entropy-12-01532}{} Nielsen F.; Nock R. The dual Voronoi diagrams with respect to representational Bregman divergences. Proceedings of the International Symposium on Voronoi Diagrams (ISVD). Copenhagen, Denmark. 23–26 June 2009.
\bibitem{B73-entropy-12-01532}\hypertarget{B73-entropy-12-01532}{} Cichocki A.; Phan A. Fast local algorithms for large scale nonnegative matrix and tensor factorizations. IEICE (invited paper). 2009. E92-A (3). pp.~708--721. \href{https://doi.org/10.1587/transfun.E92.A.708}{doi: 10.1587/transfun.E92.A.708}.
\bibitem{B74-entropy-12-01532}\hypertarget{B74-entropy-12-01532}{} Cichocki A.; Phan A.; Caiafa C. Flexible HALS algorithms for sparse non-negative matrix/tensor factorization. Proceedings of the 18th IEEE workshops on Machine Learning for Signal Processing. Cancun, Mexico. 16–19 October 2008.
\bibitem{B75-entropy-12-01532}\hypertarget{B75-entropy-12-01532}{} Dhillon I.; Sra S. Generalized Nonnegative Matrix Approximations with Bregman Divergences. Neural Information Processing Systems. MIT Press. Vancouver, Canada. 2005. pp.~283--290.
\bibitem{B76-entropy-12-01532}\hypertarget{B76-entropy-12-01532}{} Févotte C.; Bertin N.; Durrieu J.L. Nonnegative matrix factorization with the Itakura-Saito divergence with application to music analysis. Neural Comput. 2009. 21. pp.~793--830. \href{https://doi.org/10.1162/neco.2008.04-08-771}{doi: 10.1162/neco.2008.04-08-771}.
\bibitem{B77-entropy-12-01532}\hypertarget{B77-entropy-12-01532}{} Itakura F.; Saito F. Analysis synthesis telephony based on the maximum likelihood method. Proceedings of the of the 6th International Congress on Acoustics. Tokyo, Japan. 1968. pp.~17--20.
\bibitem{B78-entropy-12-01532}\hypertarget{B78-entropy-12-01532}{} Eggermont P.; LaRiccia V. On EM-like algorithms for minimum distance estimation. Technical report. Mathematical Sciences. University of Delaware. Newark, DE, USA. 1998.
\bibitem{B79-entropy-12-01532}\hypertarget{B79-entropy-12-01532}{} Févotte C.; Cemgil A.T. Nonnegative matrix factorizations as probabilistic inference in composite models. Proceedings of the 17th European Signal Processing Conference (EUSIPCO-09). Glasgow, Scotland, UK. 24–28 August 2009.
\bibitem{B80-entropy-12-01532}\hypertarget{B80-entropy-12-01532}{} Banerjee A.; Dhillon I.; Ghosh J.; Merugu S.; Modha D. A generalized maximum entropy approach to Bregman co-clustering and matrix approximation. Proceedings of the tenth ACM SIGKDD International Conference on Knowledge Discovery and Data Mining. Seattle, WA, USA. 22–25 August 2004. ACM Press. New York, NY, USA. 2004. pp.~509--514.
\bibitem{B81-entropy-12-01532}\hypertarget{B81-entropy-12-01532}{} Lafferty J. Additive models, boosting, and inference for generalized divergences. Proceedings of the 12th Annual Conference on Computational Learning Theory. Santa Cruz, CA, USA. 7–9 July 1999. ACM Press. New York, USA. 1999. pp.~125--133.
\bibitem{B82-entropy-12-01532}\hypertarget{B82-entropy-12-01532}{} Frigyik. B.A.; Srivastan S.; Gupta M. Functional Bregman divergence and Bayesion estimation of distributions. IEEE Trans. Inf. Theory. 2008. 54. pp.~5130--5139. \href{https://doi.org/10.1109/TIT.2008.929943}{doi: 10.1109/TIT.2008.929943}.
\bibitem{B83-entropy-12-01532}\hypertarget{B83-entropy-12-01532}{} Principe J. Information Theoretic Learning: Renyi’s Entropy and Kernel Perspectives. Springer. Berlin, Germany. 2010.
\bibitem{B84-entropy-12-01532}\hypertarget{B84-entropy-12-01532}{} Choi H.; Choi S.; Katake A.; Choe Y. Learning alpha-integration with partially-labeled data. Proceedings of the IEEE International Conference on Acoustics, Speech, and Signal Processing (ICASSP-2010). Dallas, TX, USA. 14–19 March 2010.
\bibitem{B85-entropy-12-01532}\hypertarget{B85-entropy-12-01532}{} Jones M.; Hjort N.; Harris I.R.; Basu A. A comparison of related density-based minimum divergence estimators. Biometrika. 1998. 85. pp.~865--873. \href{https://doi.org/10.1093/biomet/88.3.865}{doi: 10.1093/biomet/88.3.865}.
\end{thebibliography}

\bigskip\noindent\small © 2010 作者；被许可方：MDPI，瑞士巴塞尔。原文为开放获取文章，按知识共享署名许可（CC BY 3.0）发布。
\end{document}
